% Traduction : la phrase impérative avec ou sans pronoms compléments ; emisión de salud / podcast ecológico — Espagnol Tle A — Cours signé OnBuch+
% Programme officiel MINESEC. Compiler avec : tectonic E08-imperativo-pronombres-podcast.tex
\newcommand{\DOCMATIERE}{Espagnol}
\newcommand{\DOCNIVEAU}{Tle A}
\newcommand{\DOCMODULE}{Módulo 2 : Environnement, santé et bien-être}
\newcommand{\DOCLECON}{E08}
\newcommand{\DOCTITRE}{Traduction : la phrase impérative avec ou sans pronoms compléments ; emisión de salud / podcast ecológico}
% OnBuch+ — préambule « cours signés » ENRICHI (compilé avec tectonic / XeLaTeX)
% Système visuel « Pop » d'OnBuch+ : fond crème (jamais blanc), bordures encre
% épaisses, ombres « dures » décalées, titres Archivo Black en capitales.
% Version MATHÉMATIQUES : graphes (pgfplots/tikz), boîtes pédagogiques
% (définition, propriété, méthode, attention, le savais-tu, exercice résolu,
% exercice, TP, bilan), page de garde et signature OnBuch+ ancrée en bas.
%
% Macros attendues dans main.tex AVANT \input{preamble} :
%   \DOCMATIERE  \DOCNIVEAU  \DOCMODULE  \DOCLECON (numéro)  \DOCTITRE
\documentclass[11pt]{article}

\usepackage{fontspec}
\usepackage[spanish,french]{babel} % français = langue principale ; l'espagnol via \esp{..} / espagnol
\usepackage[a4paper,margin=2cm,top=3.4cm,bottom=2.8cm,headheight=1.4cm,footskip=1.3cm]{geometry}
\usepackage{amsmath,amssymb,amsfonts}
\usepackage{mathtools} % \xmapsto, \coloneqq... souvent utilisés par les modèles
\usepackage{cancel} % simplifications barrées (bilans de forces)
% \abs{z} (module, valeur absolue) et \norm{u} (norme), utilisés par les modèles
\providecommand{\abs}{}\let\abs\relax\DeclarePairedDelimiter{\abs}{\lvert}{\rvert}
\providecommand{\norm}{}\let\norm\relax\DeclarePairedDelimiter{\norm}{\lVert}{\rVert}
\usepackage[table]{xcolor} % [table] : fournit aussi \rowcolors (alternance de lignes)
\usepackage{pagecolor}
\usepackage{tikz}
\usetikzlibrary{arrows.meta,positioning,calc,shapes.geometric,shapes.misc,shapes.arrows,decorations.pathmorphing,decorations.pathreplacing,decorations.markings,patterns,shadows,mindmap,trees,fit,backgrounds,matrix,intersections,angles,quotes}
\usepackage{pgfplots}
\usepackage[version=4]{mhchem} % équations nucléaires \ce{^{235}_{92}U}
\usepackage[european,siunitx]{circuitikz} % schémas électriques
\pgfplotsset{compat=1.18}
\usepgfplotslibrary{fillbetween,groupplots} % aires entre courbes (calcul intégral)
\usepackage{enumitem}
\usepackage{fancyhdr}
% [most] charge toutes les bibliothèques tcolorbox (listings, minted, raster,
% documentation...) et gonfle inutilement la mémoire TeX sur un document long
% (~300 boîtes) : on ne charge que ce qui est réellement utilisé.
\usepackage[skins,breakable]{tcolorbox}
\usepackage{graphicx}
\usepackage{colortbl}
\usepackage{booktabs}
\usepackage{tabularx}
\usepackage{diagbox} % case d'en-tête diagonale des tableaux à double entrée
% Colonnes extensibles centrée / gauche / droite (C, L, R), souvent utilisées par les modèles
\newcolumntype{C}{>{\centering\arraybackslash}X}
\newcolumntype{L}{>{\raggedright\arraybackslash}X}
\newcolumntype{R}{>{\raggedleft\arraybackslash}X}
\usepackage{array}
\usepackage{multicol}
\usepackage{siunitx}
% parse-numbers=false : accepte aussi \SI{2,5 \times 10^{-3}}{...}
\sisetup{locale=FR,output-decimal-marker={,},group-minimum-digits=4,parse-numbers=false}
\DeclareSIUnit\ohm{\ensuremath{\symup{\Omega}}} % U+2126 absent de la police
\DeclareSIUnit\division{div} % réglages d'oscilloscope (ms/div, V/div)
\DeclareSIUnit\tr{tr} % tours (vitesse de rotation, ex. \SI{1500}{\tr\per\minute})
\DeclareSIUnit\molar{M} % concentration molaire (ex. \SI{3}{\molar}, \SI{1}{\micro\molar})
\DeclareSIUnit\an{an} % année (le modèle confond parfois avec \year, un compteur TeX) — ex. \SI{5}{\kilo\meter\per\an}
\DeclareSIUnit\FCFA{FCFA} % franc CFA (ex. \SI{500}{\FCFA\per\kilo\gram})
\DeclareSIUnit\degreeBrix{\textdegree Brix} % teneur en sucre d'un sirop/jus (réfractométrie)
\DeclareSIUnit\month{mois} % le modèle utilise parfois \month, un compteur TeX, comme unité de durée
\usepackage{qrcode}
% \degree : commande fréquemment utilisée par les modèles pour un angle ou une
% température en dehors de \SI{}{} (non fournie par les paquets chargés).
\providecommand{\degree}{^{\circ}}
% \qed (fin de démonstration), utilisé par les modèles sans amsthm
\providecommand{\qed}{\hfill$\square$}
% \barre : double barre des valeurs interdites dans un tableau de variations (tkz-tab)
\providecommand{\barre}{\ensuremath{\|}}
% environnement proof (amsthm non chargé) utilisé par les modèles
\ifdefined\proof\else\newenvironment{proof}[1][Démonstration]{\par\noindent\textit{#1.}\ }{\qed\par}\fi
\usepackage{lastpage}
\usepackage{hyperref}
\hypersetup{hidelinks}

% ============================================================
% Palette « Pop » OnBuch+ — mêmes hex que la classe P (plus_ui.dart)
% ============================================================
\definecolor{poporange}{HTML}{F59321}
\definecolor{popdark}{HTML}{E07A0C}
\definecolor{popcream}{HTML}{FFF6E8}
\definecolor{popink}{HTML}{1C1714}
\definecolor{poppurple}{HTML}{7A5AE0}
\definecolor{popgreen}{HTML}{2E9E6A}
\definecolor{poppink}{HTML}{E0457B}
\definecolor{popblue}{HTML}{2F7DE1}
\definecolor{popgold}{HTML}{C98A12}
\definecolor{popmuted}{HTML}{8A7F76}
\definecolor{popline}{HTML}{EADFD2}
\definecolor{popgray}{HTML}{8A7F76}
\colorlet{popgrey}{popgray}
% Alias tolérants (le modèle nomme parfois les couleurs différemment)
\colorlet{popwhite}{white}
\colorlet{popblack}{popink}
\colorlet{popred}{poppink}
\colorlet{popyellow}{popgold}
% Teintes claires (fonds de tableaux/encadrés) — le modèle nomme parfois
% les couleurs avec un suffixe L (ex. popblueL) sans les avoir définies.
\colorlet{poporangeL}{poporange!20}
\colorlet{popdarkL}{popdark!20}
\colorlet{poppurpleL}{poppurple!20}
\colorlet{popgreenL}{popgreen!20}
\colorlet{poppinkL}{poppink!20}
\colorlet{popblueL}{popblue!20}
\colorlet{popgoldL}{popgold!20}
\colorlet{popredL}{popred!20}
\colorlet{popgrayL}{popgray!20}
% Teintes claires pour fonds de tableaux / zones
\colorlet{popblueL}{popblue!10!white}
\colorlet{poporangeL}{poporange!14!white}
\colorlet{popgreenL}{popgreen!12!white}
\colorlet{poppurpleL}{poppurple!12!white}
\colorlet{poppinkL}{poppink!12!white}
\colorlet{popgoldL}{popgold!14!white}

\pagecolor{popcream}

% ============================================================
% Typographie
% ============================================================
\defaultfontfeatures{Ligatures=TeX}
\setmainfont{PlusJakartaSans-Medium.ttf}[Path=../../../fonts/,BoldFont=PlusJakartaSans-Bold.ttf]
% Maths en sans-serif (Fira Math) avec chiffres et lettres droites de Plus
% Jakarta, pour que les formules se fondent dans le texte.
\usepackage[math-style=ISO,bold-style=ISO]{unicode-math}
\setmathfont{FiraMath-Regular.otf}[Path=../../../fonts/]
\setmathfont{PlusJakartaSans-Medium.ttf}[Path=../../../fonts/,range={up/num,up/{Latin,latin},\mathup}]
\usepackage{icomma} % virgule décimale sans espace en mode math
% Caractères mathématiques Unicode absents des polices de texte (Plus Jakarta,
% Archivo Black) : rendus par leur équivalent mathématique, en texte comme en
% formule — sinon ils s'affichent en carré vide (ex. « Arithmétique dans ℤ »).
\usepackage{newunicodechar}
\newunicodechar{ℝ}{\ensuremath{\mathbb{R}}}
\newunicodechar{ℤ}{\ensuremath{\mathbb{Z}}}
\newunicodechar{ℂ}{\ensuremath{\mathbb{C}}}
\newunicodechar{ℕ}{\ensuremath{\mathbb{N}}}
\newunicodechar{ℚ}{\ensuremath{\mathbb{Q}}}
\newunicodechar{𝔻}{\ensuremath{\mathbb{D}}}
\newunicodechar{α}{\ensuremath{\alpha}}
\newunicodechar{β}{\ensuremath{\beta}}
\newunicodechar{γ}{\ensuremath{\gamma}}
\newunicodechar{δ}{\ensuremath{\delta}}
\newunicodechar{ε}{\ensuremath{\varepsilon}}
\newunicodechar{θ}{\ensuremath{\theta}}
\newunicodechar{λ}{\ensuremath{\lambda}}
\newunicodechar{μ}{\ensuremath{\mu}}
\newunicodechar{σ}{\ensuremath{\sigma}}
\newunicodechar{τ}{\ensuremath{\tau}}
\newunicodechar{φ}{\ensuremath{\varphi}}
\newunicodechar{ψ}{\ensuremath{\psi}}
\newunicodechar{ω}{\ensuremath{\omega}}
\newunicodechar{Σ}{\ensuremath{\Sigma}}
% majuscules produites par \MakeUppercase dans les titres (α-aminés -> Α)
\newunicodechar{Α}{\ensuremath{\alpha}}
\newunicodechar{Β}{\ensuremath{\beta}}
\newunicodechar{Γ}{\ensuremath{\Gamma}}
\newunicodechar{Δ}{\ensuremath{\Delta}}
\newunicodechar{Ε}{\ensuremath{\varepsilon}}
\newunicodechar{Θ}{\ensuremath{\Theta}}
\newunicodechar{Λ}{\ensuremath{\Lambda}}
\newunicodechar{Μ}{\ensuremath{\mu}}
\newunicodechar{Τ}{\ensuremath{\tau}}
\newunicodechar{Φ}{\ensuremath{\Phi}}
\newunicodechar{Ψ}{\ensuremath{\Psi}}
\newunicodechar{Ω}{\ensuremath{\Omega}}
\newunicodechar{↦}{\ensuremath{\mapsto}}
\newunicodechar{∈}{\ensuremath{\in}}
\newunicodechar{∉}{\ensuremath{\notin}}
\newunicodechar{∧}{\ensuremath{\wedge}}
\newunicodechar{⇔}{\ensuremath{\Leftrightarrow}}
\newunicodechar{⟺}{\ensuremath{\Longleftrightarrow}}
\newunicodechar{⇒}{\ensuremath{\Rightarrow}}
\newunicodechar{⟹}{\ensuremath{\Longrightarrow}}
\newunicodechar{∘}{\ensuremath{\circ}}
\newunicodechar{∪}{\ensuremath{\cup}}
\newunicodechar{∩}{\ensuremath{\cap}}
\newunicodechar{≡}{\ensuremath{\equiv}}
\newunicodechar{⊂}{\ensuremath{\subset}}
\newunicodechar{⊥}{\ensuremath{\perp}}
\newunicodechar{∃}{\ensuremath{\exists}}
\newunicodechar{∀}{\ensuremath{\forall}}
\newunicodechar{∛}{\ensuremath{\sqrt[3]{\,}}}
\newunicodechar{∜}{\ensuremath{\sqrt[4]{\,}}}
\newunicodechar{⃗}{\ensuremath{{}^{\to}}}
\newunicodechar{ⁿ}{\ifmmode^{n}\else\textsuperscript{\ensuremath{n}}\fi}
\newunicodechar{ˣ}{\ifmmode^{x}\else\textsuperscript{\ensuremath{x}}\fi}
\newunicodechar{ᵇ}{\ifmmode^{b}\else\textsuperscript{\ensuremath{b}}\fi}
\newunicodechar{ʳ}{\ifmmode^{r}\else\textsuperscript{\ensuremath{r}}\fi}
\newunicodechar{ᵘ}{\ifmmode^{u}\else\textsuperscript{\ensuremath{u}}\fi}
\newunicodechar{ʸ}{\ifmmode^{y}\else\textsuperscript{\ensuremath{y}}\fi}
\newunicodechar{ᵃ}{\ifmmode^{a}\else\textsuperscript{\ensuremath{a}}\fi}
\newunicodechar{ᵉ}{\ifmmode^{e}\else\textsuperscript{\ensuremath{e}}\fi}
\newunicodechar{ᵏ}{\ifmmode^{k}\else\textsuperscript{\ensuremath{k}}\fi}
\newunicodechar{ᵖ}{\ifmmode^{p}\else\textsuperscript{\ensuremath{p}}\fi}
\newunicodechar{ᴺ}{\ifmmode^{N}\else\textsuperscript{\ensuremath{N}}\fi}
\newunicodechar{ᶜ}{\ifmmode^{c}\else\textsuperscript{\ensuremath{c}}\fi}
\newunicodechar{ʰ}{\ifmmode^{h}\else\textsuperscript{\ensuremath{h}}\fi}
\newunicodechar{ᵠ}{\ifmmode^{q}\else\textsuperscript{\ensuremath{q}}\fi}
\newunicodechar{ₙ}{\ifmmode_{n}\else\textsubscript{\ensuremath{n}}\fi}
\newunicodechar{ₐ}{\ifmmode_{a}\else\textsubscript{\ensuremath{a}}\fi}
\newunicodechar{ᵢ}{\ifmmode_{i}\else\textsubscript{\ensuremath{i}}\fi}
\newunicodechar{ᵣ}{\ifmmode_{r}\else\textsubscript{\ensuremath{r}}\fi}
\newunicodechar{ₖ}{\ifmmode_{k}\else\textsubscript{\ensuremath{k}}\fi}
\newunicodechar{ᵦ}{\ifmmode_{\beta}\else\textsubscript{\ensuremath{\beta}}\fi}
\newunicodechar{ₓ}{\ifmmode_{x}\else\textsubscript{\ensuremath{x}}\fi}
\newfontfamily\popdisplay{ArchivoBlack-Regular.ttf}[Path=../../../fonts/]
\newfontfamily\poplogo{SpaceGrotesk-Bold.ttf}[Path=../../../fonts/]
\newfontfamily\popsemi{PlusJakartaSans-SemiBold.ttf}[Path=../../../fonts/]
\newfontfamily\popxbold{PlusJakartaSans-ExtraBold.ttf}[Path=../../../fonts/]
\newfontfamily\popxboldls{PlusJakartaSans-ExtraBold.ttf}[Path=../../../fonts/,LetterSpace=3.0]

\setlist[enumerate]{leftmargin=1.7em,itemsep=5pt,topsep=4pt}
\setlist[itemize]{leftmargin=1.7em,itemsep=4pt,topsep=4pt,label={\color{poporange}\textbullet}}
\setlength{\parindent}{0pt}
\setlength{\parskip}{5pt}
\renewcommand{\arraystretch}{1.25}

% pgfplots : style Pop par défaut
\pgfplotsset{
  every axis/.append style={axis line style={popink,line width=1pt},tick style={popink},
    grid style={popline,line width=0.5pt},label style={font=\small},tick label style={font=\footnotesize}},
  popcurve/.style={poporange,line width=1.8pt,no markers,smooth},
}
% Même style disponible dans un \draw TikZ simple (hors pgfplots)
\tikzset{popcurve/.style={poporange,line width=1.8pt}}
\tikzset{popgrid/.style={popline,very thin}}
\colorlet{popcurve}{poporange} % le nom du style est parfois utilisé comme couleur

% ============================================================
% Pill / squiggle
% ============================================================
\newcommand{\poppill}[3]{%
  \tikz[baseline=-0.55ex]\node[draw=popink,line width=1.2pt,rounded corners=2.6pt,%
    fill=#2,inner xsep=6.5pt,inner ysep=3pt]{%
    \popxboldls\fontsize{7}{7}\selectfont\color{#3}\MakeUppercase{#1}};%
}
\newcommand{\popsquiggle}[1]{%
  \tikz[baseline=-0.4ex]{
    \draw[#1,line width=2.6pt,line cap=round]
      plot [smooth] coordinates {(0,0.09) (0.34,0.01) (0.68,0.21) (1.02,0.01) (1.36,0.10)};
  }%
}
% Mot-clé surligné (marqueur orange sous le texte)
\newcommand{\cle}[1]{{\popxbold\color{popdark}#1}}

% ============================================================
% En-tête / pied
% ============================================================
\pagestyle{fancy}
\fancyhf{}
\renewcommand{\headrulewidth}{0pt}
\renewcommand{\footrulewidth}{0pt}
\lhead{\raisebox{-0.15ex}{\poplogo\fontsize{13}{13}\selectfont\color{popink}OnBuch\color{poporange}+}}
\rhead{\poppill{\DOCMATIERE\ \textbullet\ \DOCNIVEAU\ \textbullet\ Leçon \DOCLECON}{white}{popink}}
\lfoot{\popxbold\fontsize{7.6}{7.6}\selectfont\color{popmuted}\MakeUppercase{Cours signé OnBuch+}\ \textbullet\ {\popsemi\DOCTITRE}}
\rfoot{\tikz[baseline=-0.6ex]\node[rounded corners=8.5pt,fill=popink,minimum height=17pt,minimum width=17pt,inner xsep=5pt,inner sep=0]{\popxbold\fontsize{8}{8}\selectfont\color{popcream}\thepage\,/\,\pageref*{LastPage}};}

% ============================================================
% Titres
% ============================================================
\newcommand{\chaptitre}[2]{%
  \begin{tcolorbox}[enhanced,colback=poporange,colframe=popink,boxrule=2.6pt,arc=11pt,
      drop shadow=popink,left=15pt,right=15pt,top=12pt,bottom=11pt,before skip=3pt,after skip=18pt]
    \poppill{#2}{popink}{popcream}\par\vspace{9pt}
    {\raggedright\hyphenpenalty=10000\exhyphenpenalty=10000\popdisplay\fontsize{22}{24}\selectfont\color{popcream}\MakeUppercase{#1}\par}
  \end{tcolorbox}
}
% Section numérotée : pastille orange avec le numéro + titre Archivo
\newcounter{popsec}
\newcommand{\coursec}[1]{%
  \stepcounter{popsec}\setcounter{popsub}{0}%
  \par\vspace{16pt}\noindent
  \begin{minipage}{\linewidth}
  \tikz[baseline=-0.7ex]\node[draw=popink,line width=1.6pt,fill=poporange,rounded corners=4pt,minimum size=22pt,inner sep=0]{\popdisplay\fontsize{12}{12}\selectfont\color{popcream}\thepopsec};%
  \hspace{8pt}{\popdisplay\fontsize{15}{17}\selectfont\color{popink}\MakeUppercase{#1}}\par
  \vspace{4pt}\noindent{\color{poporange}\rule{36pt}{3.6pt}}
  \end{minipage}\par\nopagebreak\vspace{8pt}
}
\newcounter{popsub}
\newcommand{\courssub}[1]{%
  \stepcounter{popsub}%
  \par\vspace{9pt}\noindent{\popxbold\fontsize{11.5}{13}\selectfont\color{popink}{\color{poporange}\thepopsec.\thepopsub}\hspace{6pt}#1}\par\nopagebreak\vspace{4pt}
}

% ============================================================
% Boîtes pédagogiques — même recette : fond blanc, bordure épaisse colorée,
% ombre dure encre, badge plein. Titre optionnel [#1].
% ============================================================
\tcbset{popbase/.style={breakable,enhanced,colback=white,boxrule=2.3pt,arc=9pt,
  shadow={3pt}{-3pt}{0pt}{popink},left=14pt,right=14pt,top=11pt,bottom=11pt,before skip=11pt,after skip=15pt,
  fontupper=\popsemi\fontsize{10.6}{14.5}\selectfont\color{popink},
  fontlower=\popsemi\fontsize{10.6}{14.5}\selectfont\color{popink}}}
% Coche dessinée (le glyphe ✓ n'existe pas dans les polices du document)
\newcommand{\popcheck}[1]{\tikz[baseline=-0.5ex]\draw[#1,line width=1.6pt,line cap=round,line join=round](-0.13,0.0)--(-0.04,-0.09)--(0.13,0.1);}
\providecommand{\coche}{\popcheck{popgreen}}
\providecommand{\resultat}[1]{\par\smallskip\noindent\fcolorbox{popgreen}{popgreenL}{\parbox{\dimexpr\linewidth-2\fboxsep-2\fboxrule\relax}{\textbf{Résultat.} #1}}\par\smallskip}
\newcommand{\popboxhead}[3]{\poppill{#1}{#2}{white}\ifx\relax#3\relax\else\hspace{6pt}{\popxbold\fontsize{10.6}{12}\selectfont\color{popink}#3}\fi\par\vspace{7pt}}

\newtcolorbox{aretenir}[1][]{popbase,colframe=popgreen,before upper={\popboxhead{À retenir}{popgreen}{#1}}}
% Texte en espagnol (document de lecture, dialogue, poème, extrait) : cadre
% d'étude. Un titre optionnel (source, auteur) s'affiche dans l'en-tête.
\newtcolorbox{textebox}[1][]{popbase,colframe=popblue,before upper={\popboxhead{Texto}{popblue}{#1}\begin{otherlanguage}{spanish}},after upper={\end{otherlanguage}}}
% Marques ✓ / ✗ (forme correcte / fautive) souvent utilisées par les modèles
\providecommand{\cross}{\textcolor{poppink}{\ensuremath{\times}}}
\providecommand{\xmark}{\cross}\providecommand{\crossmark}{\cross}\providecommand{\cmark}{\ensuremath{\checkmark}}
% Ligne de réponse / justification à compléter (inventée par les modèles)
\providecommand{\justification}[1]{\par\noindent\textit{Justification~:} \dotfill\par}
% \textipa (package tipa non chargé) : la police ne couvre pas l'API -> simple passage
\providecommand{\textipa}[1]{#1}
% Soulignement (ulem non chargé) : \uline, \ul
\providecommand{\uline}[1]{\underline{#1}}\providecommand{\ul}[1]{\underline{#1}}
% Mot ou phrase en espagnol à l'intérieur d'un paragraphe français
\newcommand{\esp}[1]{\ifmmode\text{\textit{#1}}\else\begingroup\selectlanguage{spanish}\itshape #1\endgroup\fi}
\newtcolorbox{exemplebox}[1][]{popbase,colframe=poporange,before upper={\popboxhead{Exemple}{poporange}{#1}}}
\newtcolorbox{definition}[1][]{popbase,colframe=popblue,before upper={\popboxhead{Définition}{popblue}{#1}}}
\newtcolorbox{propriete}[1][]{popbase,colframe=poppurple,before upper={\popboxhead{Propriété}{poppurple}{#1}}}
\newtcolorbox{methode}[1][]{popbase,colframe=popgold,before upper={\popboxhead{Méthode}{popgold}{#1}}}
\newtcolorbox{attention}[1][]{popbase,colframe=poppink,before upper={\popboxhead{Attention}{poppink}{#1}}}
\newtcolorbox{experience}[1][]{popbase,colframe=popblue,colback=popblueL,before upper={\popboxhead{Expérience}{popblue}{#1}}}
% « Le savais-tu ? » : carte violette pleine (contexte, culture, Cameroun)
\newtcolorbox{savaistu}[1][]{popbase,colback=poppurple,colframe=popink,
  fontupper=\popsemi\fontsize{10.4}{14.2}\selectfont\color{white},
  before upper={\poppill{Le savais-tu\,?}{white}{poppurple}\ifx\relax#1\relax\else\hspace{6pt}{\popxbold\color{white}#1}\fi\par\vspace{7pt}}}
% Exercice résolu : énoncé en haut, \tcblower puis la solution sur fond crème
\newcounter{popexo}\newcounter{popexr}
\newtcolorbox{exoresolu}[1][]{popbase,colframe=poporange,colbacklower=popcream,
  segmentation style={popink,line width=1.2pt,dashed},
  before upper={\stepcounter{popexr}\popboxhead{Exercice résolu \thepopexr}{poporange}{#1}},
  before lower={\poppill{Solution}{popink}{popcream}\par\vspace{6pt}}}
% Exercice (non résolu en place) — la correction est dans la partie Corrigés
\newtcolorbox{exercice}[1][]{popbase,colframe=popink,
  before upper={\stepcounter{popexo}\popboxhead{Exercice \thepopexo}{popink}{#1}}}
% Corrigé
\newtcolorbox{corrige}[1][]{popbase,colframe=popgreen,colback=popgreenL,
  before upper={\popboxhead{Corrigé}{popgreen}{#1}}}
% Objectifs / prérequis / bilan
\newtcolorbox{objectifs}{popbase,colframe=popink,colback=poporangeL,
  before upper={\popboxhead{Objectifs de la leçon}{popink}{}}}
\newtcolorbox{prerequis}{popbase,colframe=popink,colback=popblueL,
  before upper={\popboxhead{Prérequis}{popblue}{}}}
\newtcolorbox{bilan}{popbase,colframe=popink,colback=white,boxrule=2.8pt,
  before upper={\popboxhead{Fiche bilan}{poporange}{L'essentiel à connaître pour l'examen}}}
% Figure encadrée avec légende
\newtcolorbox{popfigure}[1][]{enhanced,breakable=false,colback=white,colframe=popline,boxrule=1.4pt,arc=8pt,
  left=10pt,right=10pt,top=10pt,bottom=8pt,before skip=10pt,after skip=12pt,halign=center,
  lower separated=false,halign lower=center,
  fontlower=\popsemi\fontsize{9}{11.5}\selectfont\color{popmuted},
  #1}
\newcommand{\legende}[1]{\par\vspace{6pt}{\popsemi\fontsize{9}{11.5}\selectfont\color{popmuted}#1}}

% ============================================================
% Signature OnBuch+
% ============================================================
\newcommand{\poplogoinline}[1]{{\poplogo\fontsize{#1}{#1}\selectfont\color{popink}OnBuch\color{poporange}+}}
\newcommand{\signatureonbuch}{%
  \begin{tcolorbox}[enhanced,colback=popink,colframe=popink,boxrule=2.2pt,arc=10pt,
      drop shadow=poporange,left=14pt,right=14pt,top=10pt,bottom=10pt,before skip=0pt,after skip=0pt]
    \tikz[baseline=-0.7ex]\node[circle,fill=popgreen,draw=popcream,line width=1.3pt,minimum size=22pt,inner sep=0]{\popcheck{white}};%
    \hspace{9pt}\begin{minipage}[c]{0.62\linewidth}
      {\popxbold\fontsize{12}{13}\selectfont\color{popcream}Signé\ \textbullet\ {\poplogo OnBuch\color{poporange}+}}\par\vspace{2pt}
      {\raggedright\popsemi\fontsize{8}{10.5}\selectfont\color{popline}\MakeUppercase{Équipe pédagogique OnBuch+}\\\MakeUppercase{Conforme au programme officiel MINESEC}\par}
    \end{minipage}\hfill
    {\poplogo\fontsize{22}{22}\selectfont\color{popcream}OnBuch\color{poporange}+}
  \end{tcolorbox}%
}

% ============================================================
% Page de garde
% #1 titre · #2 matière · #3 niveau · #4 module · #5 n° leçon · #6 durée ·
% #7 description / objectifs (texte ou itemize)
% ============================================================
\newcommand{\pagegarde}[7]{%
  \thispagestyle{empty}
  \newgeometry{a4paper,margin=1.7cm,top=1.6cm,bottom=1.4cm}%
  \begin{tikzpicture}[remember picture,overlay]
    \fill[poporange!18] ([xshift=-3.2cm,yshift=-3.6cm]current page.north east) circle (4.6cm);
    \fill[poppurple!14] ([xshift=2.4cm,yshift=7.6cm]current page.south west) circle (3.8cm);
  \end{tikzpicture}%
  {\poplogo\fontsize{30}{30}\selectfont\color{popink}OnBuch\color{poporange}+}\hfill
  \raisebox{0.6ex}{\poppill{Cours signé \textbullet\ Premium}{popink}{popcream}}\par
  \vspace{1pt}\popsquiggle{poppurple}\par
  \vspace{8pt}
  \begin{tcolorbox}[enhanced,colback=poporange,colframe=popink,boxrule=2.8pt,arc=13pt,
      drop shadow=popink,left=18pt,right=18pt,top=12pt,bottom=13pt,after skip=10pt]
    \poppill{#2 \textbullet\ #3}{popink}{popcream}\hspace{5pt}\poppill{#4}{white}{popink}\par\vspace{8pt}
    {\popdisplay\fontsize{14}{14}\selectfont\color{popink}LEÇON #5}\par\vspace{4pt}
    {\raggedright\hyphenpenalty=10000\exhyphenpenalty=10000\popdisplay\fontsize{25}{27}\selectfont\color{popcream}\MakeUppercase{#1}\par}
    \vspace{8pt}
    {\popxbold\fontsize{9.5}{9.5}\selectfont\color{popink}\MakeUppercase{Durée conseillée : #6}}
  \end{tcolorbox}
  \begin{tcolorbox}[enhanced,colback=white,colframe=popink,boxrule=2.2pt,arc=10pt,
      drop shadow=popink,left=16pt,right=16pt,top=10pt,bottom=10pt,after skip=8pt]
    \poppill{Dans cette leçon}{poppurple}{white}\par\vspace{5pt}
    \popsemi\fontsize{9.3}{12.6}\selectfont\color{popink}#7
  \end{tcolorbox}
  \noindent\begin{minipage}[t]{0.487\linewidth}
    \begin{tcolorbox}[enhanced,colback=popgreen,colframe=popink,boxrule=2.2pt,arc=10pt,
        drop shadow=popink,left=11pt,right=11pt,top=10pt,bottom=10pt,height=3.1cm,valign=center]
      \tcbox[colback=white,colframe=popink,boxrule=1.3pt,arc=5pt,boxsep=0pt,nobeforeafter,
        left=3pt,right=3pt,top=3pt,bottom=3pt,valign=center]{\qrcode[height=1.9cm]{https://play.google.com/store/apps/details?id=com.luvvix.onbuch&hl=fr}}%
      \hspace{8pt}\begin{minipage}[c]{0.5\linewidth}
        {\popdisplay\fontsize{11}{12.5}\selectfont\color{white}\MakeUppercase{Télécharge OnBuch}}\par\vspace{4pt}
        {\raggedright\popsemi\fontsize{8.4}{11}\selectfont\color{white}Gratuit \textbullet\ Google Play\\Tuteur IA, annales, concours, résultats\par}
      \end{minipage}
    \end{tcolorbox}
  \end{minipage}\hfill
  \begin{minipage}[t]{0.487\linewidth}
    \begin{tcolorbox}[enhanced,colback=poppurple,colframe=popink,boxrule=2.2pt,arc=10pt,
        drop shadow=popink,left=11pt,right=11pt,top=10pt,bottom=10pt,height=3.1cm,valign=center]
      \tikz[baseline=-0.7ex]\node[circle,fill=white,draw=popink,line width=1.3pt,minimum size=1.6cm,inner sep=0]{\popdisplay\fontsize{24}{24}\selectfont\color{poppurple}+};%
      \hspace{8pt}\begin{minipage}[c]{0.6\linewidth}
        {\popdisplay\fontsize{11}{12.5}\selectfont\color{white}\MakeUppercase{Passe à OnBuch+}}\par\vspace{4pt}
        {\raggedright\popsemi\fontsize{8.4}{11}\selectfont\color{white}Tous les cours signés, Tuteur IA illimité, fiches PDF — onglet OnBuch+ de l'app\par}
      \end{minipage}
    \end{tcolorbox}
  \end{minipage}
  \vspace{8pt}
  \popcontenu
  \vfill
  \signatureonbuch
  \restoregeometry
  \newpage
}

% Rangée « Ce cours contient » (page de garde)
\newcommand{\poptile}[3]{% #1 couleur #2 gros texte #3 légende
  \begin{tcolorbox}[enhanced,colback=#1,colframe=popink,boxrule=2pt,arc=9pt,shadow={2.5pt}{-2.5pt}{0pt}{popink},
      left=6pt,right=6pt,top=9pt,bottom=9pt,halign=center,valign=center,height=1.75cm,width=\linewidth,nobeforeafter]
    {\popdisplay\fontsize{12.5}{13}\selectfont\color{white}\MakeUppercase{#2}}\par\vspace{3pt}
    {\centering\popsemi\fontsize{7.8}{9.5}\selectfont\color{white}#3\par}
  \end{tcolorbox}}
\newcommand{\popcontenu}{%
  \par\noindent{\popxboldls\fontsize{7.5}{7.5}\selectfont\color{popmuted}CE COURS SIGNÉ CONTIENT}\par\vspace{7pt}
  \noindent\begin{minipage}[t]{0.235\linewidth}\poptile{popblue}{Cours}{complet, illustré, pas à pas}\end{minipage}\hfill
  \begin{minipage}[t]{0.235\linewidth}\poptile{poporange}{Méthodes}{et exercices résolus}\end{minipage}\hfill
  \begin{minipage}[t]{0.235\linewidth}\poptile{poppink}{Exercices}{type examen + corrigés}\end{minipage}\hfill
  \begin{minipage}[t]{0.235\linewidth}\poptile{popgreen}{Fiche bilan}{l'essentiel à retenir}\end{minipage}\par}

% Sommaire
\newtcolorbox{sommaire}{enhanced,colback=white,colframe=popink,boxrule=2.2pt,arc=10pt,
  drop shadow=popink,left=18pt,right=18pt,top=13pt,bottom=13pt,before skip=6pt,after skip=20pt,
  before upper={\poppill{Sommaire}{poppurple}{white}\par\vspace{9pt}\popsemi\fontsize{10.6}{14.5}\selectfont\color{popink}}}

% Signature de fin de cours — poussée tout en bas de la dernière page
\newcommand{\findecours}{%
  \par\vfill\nopagebreak
  \begin{center}\popsquiggle{poporange}\end{center}\vspace{4pt}
  \signatureonbuch
}
\AtBeginDocument{\def\llbracket{\mathopen{[\mkern-3.5mu[}}\def\rrbracket{\mathclose{]\mkern-3.5mu]}}} % intervalles d'entiers (glyphe ⟦ absent de la police)
% Glyphes absents de Fira Math : redéfinis à partir de glyphes disponibles
\AtBeginDocument{%
  \def\setminus{\mathbin{\backslash}}\let\smallsetminus\setminus
  \def\vdots{\mathord{\vbox{\baselineskip3.2pt\lineskiplimit0pt\kern4pt\hbox{.}\hbox{.}\hbox{.}}}}%
  \def\ddots{\mathinner{\mkern1mu\raise6.5pt\hbox{.}\mkern2mu\raise3.6pt\hbox{.}\mkern2mu\raise0.7pt\hbox{.}\mkern1mu}}%
  \def\triangle{\mathord{\scalebox{0.9}{$\symup{\Delta}$}}}%
  \def\nabla{\mathord{\raisebox{\depth}{\scalebox{1}[-1]{$\symup{\Delta}$}}}}%
  \def\checkmark{\popcheck{popdark}}%
  \def\star{\mathbin{\ast}}\def\bigstar{\mathord{\ast}}%
  \def\diamond{\mathbin{\scalebox{0.6}{$\Diamond$}}}%
  \def\bigcup{\mathop{\mathchoice{\raisebox{-0.3em}{\scalebox{1.7}{$\cup$}}}{\scalebox{1.25}{$\cup$}}{\cup}{\cup}}\displaylimits}%
  \def\bigcap{\mathop{\mathchoice{\raisebox{-0.3em}{\scalebox{1.7}{$\cap$}}}{\scalebox{1.25}{$\cap$}}{\cap}{\cap}}\displaylimits}%
  \def\lesseqqgtr{\mathrel{\lessgtr}}\def\bigoplus{\mathop{\oplus}\displaylimits}%
}
\providecommand{\ssi}{si et seulement si} % abréviation fréquente
\AtBeginDocument{\providecommand{\wideparen}{\overparen}} % arc de cercle

%% FIGFIX-BEGIN v3 — correctifs globaux des figures (docs/onbuch-generation/tools/figfix)
\usepackage{adjustbox}\usepackage{etoolbox}
% toute figure TikZ/pgfplots plus large que la ligne est réduite à la largeur disponible
\BeforeBeginEnvironment{tikzpicture}{\begin{adjustbox}{max width=\linewidth}}
\AfterEndEnvironment{tikzpicture}{\end{adjustbox}}
% légendes pgfplots lisibles par-dessus les courbes
\pgfplotsset{every axis/.append style={legend style={fill=white,fill opacity=0.92,text opacity=1,draw=black!15,inner sep=3pt}}}
% étiquettes d'arêtes (syntaxe edge["..."]) avec fond blanc
\tikzset{every edge quotes/.append style={fill=white,inner sep=1.2pt}}
% bulles de cartes mentales : texte à la ligne dans la bulle, bulle un peu plus grande
\tikzset{every concept/.append style={text width=2.0cm,align=center,minimum size=2.3cm,inner sep=2pt}}
%% FIGFIX-END
\begin{document}

\pagegarde{Traduction : la phrase impérative avec ou sans pronoms compléments ; emisión de salud / podcast ecológico}{Espagnol}{Tle A}{Módulo 2 : Environnement, santé et bien-être}{E08}{2 h}{Leçon mixte du Módulo 2 (Environnement, santé et bien-être) consacrée à la traduction de la phrase impérative espagnole, avec ou sans pronoms compléments (enclise, proclise, accentuation), et à la production d'une émission orale de conseils de santé et d'écologie. Elle combine grammaire contrastive français/espagnol, entraînement au thème et à la version, et tâche finale de communication authentique : un podcast écologique.
\vspace{4pt}
\begin{multicols}{2}\small
\begin{itemize}
\item À la fin de cette leçon, je sais conjuguer l'impératif affirmatif et négatif aux quatre personnes (tú, vosotros, usted, ustedes)
\item je sais former l'impératif négatif à partir du subjonctif présent
\item je sais placer les pronoms compléments à l'impératif affirmatif (enclise : dámelo) et négatif (proclise : no me lo des)
\item je sais appliquer les règles d'accentuation écrite liées à l'enclise (dámelo, dígaselo, levántate)
\item je sais traduire correctement une phrase impérative du français vers l'espagnol et inversement, avec ou sans pronoms
\item je sais repérer et éviter les pièges fréquents des francophones (place du pronom, ordre des pronoms, accent oublié)
\end{itemize}\end{multicols}}

\begin{sommaire}
\begin{multicols}{2}
\begin{enumerate}
\item Mise en route : écouter des consignes (texte support)
\item Former l'impératif affirmatif et négatif
\item Les pronoms compléments avec l'impératif : enclise et proclise
\item Lexique thématique : santé et environnement
\item Entraînement à la traduction : thème et version
\item Production : émission de santé / podcast écologique
\item Activité d'intégration
\item Méthodes et exercices résolus
\item Exercices
\item Corrigés des exercices
\item Fiche bilan
\end{enumerate}
\end{multicols}
\end{sommaire}

\begin{objectifs}
\begin{itemize}
\item À la fin de cette leçon, je sais conjuguer l'impératif affirmatif et négatif aux quatre personnes (tú, vosotros, usted, ustedes)
\item je sais former l'impératif négatif à partir du subjonctif présent
\item je sais placer les pronoms compléments à l'impératif affirmatif (enclise : dámelo) et négatif (proclise : no me lo des)
\item je sais appliquer les règles d'accentuation écrite liées à l'enclise (dámelo, dígaselo, levántate)
\item je sais traduire correctement une phrase impérative du français vers l'espagnol et inversement, avec ou sans pronoms
\item je sais repérer et éviter les pièges fréquents des francophones (place du pronom, ordre des pronoms, accent oublié)
\item je sais utiliser le lexique de la santé et de l'environnement dans des consignes
\item je sais rédiger et présenter oralement un court podcast donnant des conseils de santé et d'écologie
\end{itemize}
\end{objectifs}

\begin{prerequis}
\begin{itemize}
\item Le présent de l'indicatif des verbes réguliers et irréguliers fréquents (classe de Première)
\item Les pronoms compléments d'objet direct (me, te, lo, la, nos, os, los, las) et indirects (me, te, le→se, nos, os, les→se)
\item L'ordre des pronoms combinés : le pronom indirect précède le pronom direct (me lo, te la, se lo)
\item Le subjonctif présent des verbes réguliers (base de l'impératif négatif et de usted/ustedes)
\item Le lexique de base du corps, de la santé et de la nature (Première, Módulo 1 et 2)
\end{itemize}
\end{prerequis}

\newpage

\coursec{Mise en route : écouter des consignes (texte support)}

\courssub{Situation d'accroche et problématique}

Imaginez la scène : il est sept heures du matin à Yaoundé, quartier Mvog-Ada. Sur les ondes d'une radio communautaire, une voix dynamique annonce un grand concours : \esp{¡Cuida tu barrio, cuida tu salud!} — « Prends soin de ton quartier, prends soin de ta santé ! ». Le défi : enregistrer le meilleur podcast écologique et sanitaire en espagnol. Les élèves de Terminale A du lycée voisin veulent participer. Ils ont déjà de bonnes idées : \esp{Recicla las botellas} (« Recycle les bouteilles »), \esp{No tires basura en la calle} (« Ne jette pas d'ordures dans la rue »), \esp{Bebe agua limpia} (« Bois de l'eau propre »).

Mais très vite, un problème surgit. Comment dire « donne-le-moi » ? Faut-il dire \esp{me lo da}, \esp{dámelo} ou \esp{da me lo} ? Et pour « ne me le donne pas », où placer les petits mots \esp{me} et \esp{lo} ? En français, nous disons « donne-le-moi » avec des traits d'union après le verbe, mais « ne me le donne pas » avec les pronoms avant. L'espagnol a-t-il la même logique ?

\textbf{Problématique de la leçon :} comment former et traduire correctement la phrase impérative en espagnol, avec ou sans pronoms compléments, afin de produire une émission orale claire et convaincante sur la santé et l'environnement ?

Pour répondre à cette question, nous allons d'abord \emph{écouter} un extrait d'émission de radio, observer les phrases qui donnent des ordres ou des conseils, puis dégager nous-mêmes la règle avant de la vérifier.

\courssub{Écoute et lecture du texte support}

Fermez les yeux : vous êtes devant votre poste de radio. L'animateur et l'animatrice de l'émission \emph{Salud y medio ambiente en tu barrio} prennent le micro. Lisez le texte à voix haute, puis écoutez la lecture du professeur.

\begin{textebox}[Salud y medio ambiente en tu barrio — émission de radio communautaire]
Buenos días, queridos oyentes. Bienvenidos a «Salud y medio ambiente en tu barrio», su programa de los martes. Hoy les damos consejos para su barrio, para su familia y para su cuerpo.

Primero, la salud. Beban agua limpia todos los días; hiervanla si es necesario, porque el agua sucia trae enfermedades. Lávense las manos antes de comer: es un gesto pequeño, pero salva vidas. Protéjanse de los mosquitos, porque el mosquito transmite el dengue y el paludismo; usen mosquiteras por la noche y no dejen agua estancada cerca de la casa. Si un niño tiene fiebre, llévenlo rápidamente al centro de salud: no esperen, no lo dejen para mañana.

Segundo, el medio ambiente. Reciclen las botellas de plástico; no las tiren al río ni a la calle. Recojan la basura delante de su puerta y deposítenla en el contenedor. Planten un árbol, cuídenlo, riéguenlo cada semana: un árbol da sombra, aire puro y frutas. No quemen la basura: el humo enferma a los niños y a los ancianos.

Y recuerden, amigos: su barrio es su casa. Cuídenlo. Su cuerpo es su tesoro. Cuídense. Escríbannos, llámenos, mándennos sus ideas para el próximo programa. ¡Hasta la semana que viene, y no lo olviden: cuida tu barrio, cuida tu salud!
\end{textebox}

\textbf{Traduction du texte.} « Bonjour, chers auditeurs. Bienvenue à "Santé et environnement dans ton quartier", votre émission du mardi. Aujourd'hui nous vous donnons des conseils pour votre quartier, pour votre famille et pour votre corps.

D'abord, la santé. Buvez de l'eau propre tous les jours ; faites-la bouillir si nécessaire, car l'eau sale apporte des maladies. Lavez-vous les mains avant de manger : c'est un petit geste, mais il sauve des vies. Protégez-vous des moustiques, car le moustique transmet la dengue et le paludisme ; utilisez des moustiquaires la nuit et ne laissez pas d'eau stagnante près de la maison. Si un enfant a de la fièvre, emmenez-le rapidement au centre de santé : n'attendez pas, ne le remettez pas à demain.

Ensuite, l'environnement. Recyclez les bouteilles en plastique ; ne les jetez ni dans la rivière ni dans la rue. Ramassez les ordures devant votre porte et déposez-les dans la poubelle. Plantez un arbre, prenez-en soin, arrosez-le chaque semaine : un arbre donne de l'ombre, de l'air pur et des fruits. Ne brûlez pas les ordures : la fumée rend malades les enfants et les personnes âgées.

Et rappelez-vous, amis : votre quartier est votre maison. Prenez-en soin. Votre corps est votre trésor. Prenez soin de vous. Écrivez-nous, appelez-nous, envoyez-nous vos idées pour la prochaine émission. À la semaine prochaine, et ne l'oubliez pas : prends soin de ton quartier, prends soin de ta santé ! »

\courssub{Glossaire du texte}

\begin{definition}[Vocabulaire de la santé et de l'environnement]
Voici les mots essentiels du texte, à apprendre par cœur :
\begin{itemize}
\item \esp{el cuerpo} : le corps
\item \esp{la basura} : les ordures, les déchets
\item \esp{reciclar} : recycler
\item \esp{el grifo} : le robinet
\item \esp{el dengue} : la dengue (maladie transmise par le moustique)
\item \esp{el mosquito} : le moustique
\item \esp{la botella} : la bouteille
\item \esp{el consejo} : le conseil
\item \esp{hervir (ie)} : faire bouillir ; ici \esp{hiervanla} = « faites-la bouillir »
\item \esp{la mosquitera} : la moustiquaire
\item \esp{el agua estancada} : l'eau stagnante
\item \esp{el contenedor} : le conteneur, la poubelle publique
\end{itemize}
\end{definition}

\begin{attention}[Faux amis et pièges lexicaux]
\begin{itemize}
\item \esp{el consejo} signifie « le conseil » (avis), mais aussi « le conseil » (assemblée, comme \esp{el consejo municipal}). Ne pas confondre avec \esp{el concierto} (le concert).
\item \esp{el grifo} est « le robinet » en Espagne ; dans plusieurs pays d'Amérique latine on dit \esp{la llave} ou \esp{el chorro}. Un francophone ne doit jamais traduire \esp{grifo} par « griffe » !
\item \esp{la basura} : ordures. Le verbe « jeter » se dit \esp{tirar} ou \esp{echar}, jamais \esp{jetar} (ce verbe n'existe pas).
\item Attention à l'accent : \esp{hiervanla} porte un accent écrit à cause du pronom ajouté au verbe. Nous y reviendrons.
\end{itemize}
\end{attention}

\courssub{Questions de compréhension}

\textbf{Compréhension globale.}
\begin{enumerate}
\item Qui parle dans ce texte ? À qui s'adresse-t-on ?
\item Quel est le type de document (lettre, article de presse, émission de radio, poème) ? Justifiez avec deux indices du texte.
\item Quels sont les deux grands thèmes de l'émission ?
\item Quel est le but des animateurs : informer, convaincre, donner des ordres, raconter ?
\end{enumerate}

\textbf{Compréhension détaillée.}
\begin{enumerate}
\item Citez trois conseils concernant la santé.
\item Citez trois conseils concernant l'environnement.
\item Que faut-il faire si un enfant a de la fièvre ?
\item Pourquoi ne faut-il pas brûler les ordures ?
\item Comment les auditeurs peuvent-ils participer à la prochaine émission ?
\end{enumerate}

\begin{corrige}[Compréhension]
\textbf{Globale.} 1. Deux animateurs de radio s'adressent aux auditeurs (\esp{queridos oyentes}, \esp{les damos}, \esp{su barrio}). 2. C'est une émission de radio : indices \esp{bienvenidos a... su programa}, \esp{escríbannos, llámenos}, \esp{hasta la semana que viene}. 3. La santé (\esp{primero, la salud}) et l'environnement (\esp{segundo, el medio ambiente}). 4. Donner des conseils et des consignes : le texte est presque entièrement à l'impératif.

\textbf{Détaillée.} 1. Boire de l'eau propre, se laver les mains, se protéger des moustiques avec des moustiquaires. 2. Recycler les bouteilles, ramasser les ordures, planter et arroser un arbre. 3. L'emmener rapidement au centre de santé, sans attendre. 4. Parce que la fumée rend malades les enfants et les personnes âgées. 5. En écrivant, en téléphonant et en envoyant leurs idées.
\end{corrige}

\courssub{Repérage : à la recherche des impératifs}

Avant d'énoncer la règle, observons. Reprenez le texte avec deux couleurs : \textbf{soulignez} tous les verbes qui donnent un ordre, un conseil ou une interdiction ; \textbf{entourez} les petits mots comme \esp{me}, \esp{te}, \esp{lo}, \esp{la}, \esp{se}, \esp{nos} collés au verbe ou placés devant.

\begin{methode}[Comment repérer un impératif dans un texte]
\begin{enumerate}
\item Cherchez un verbe \textbf{sans sujet exprimé} : pas de \esp{yo}, \esp{tú}, \esp{él} devant.
\item Demandez-vous : ce verbe donne-t-il un \textbf{ordre}, un \textbf{conseil} ou une \textbf{interdiction} ? Si oui, c'est un impératif.
\item Vérifiez la forme : à l'affirmatif, le verbe est souvent « nu » (\esp{bebe}, \esp{recicla}) ou porte des pronoms collés (\esp{dámelo}, \esp{cuídate}) ; au négatif, il est précédé de \esp{no} et ressemble au subjonctif (\esp{no tires}, \esp{no lo dejen}).
\item Attention aux infinitifs d'affichage (\esp{no fumar}) et aux présents de vérité générale : ce ne sont pas des impératifs.
\end{enumerate}
\end{methode}

\begin{exemplebox}[Deux exemples traduits]
\esp{Bebe agua limpia.} = « Bois de l'eau propre. » (impératif affirmatif, sans pronom)

\esp{No la tires.} = « Ne la jette pas. » (impératif négatif : le pronom \esp{la} est placé \emph{avant} le verbe, entre \esp{no} et le verbe)

\esp{Cuídenlo.} = « Prenez-en soin. » (impératif affirmatif : le pronom \esp{lo} est \emph{collé après} le verbe)
\end{exemplebox}

\courssub{Classement des exemples relevés}

Voici les impératifs du texte, classés en deux colonnes. Observez bien la place des pronoms dans chaque colonne.

\begin{center}
\begin{tabularx}{\linewidth}{|X|X|}
\hline
\rowcolor{popblueL} \textbf{Impératif affirmatif} & \textbf{Impératif négatif} \\
\hline
\esp{Beban agua limpia.} & \esp{No las tiren al río.} \\
\hline
\esp{Hiervanla si es necesario.} & \esp{No dejen agua estancada.} \\
\hline
\esp{Lávense las manos.} & \esp{No lo dejen para mañana.} \\
\hline
\esp{Protéjanse de los mosquitos.} & \esp{No quemen la basura.} \\
\hline
\esp{Reciclen las botellas.} & \esp{No lo olviden.} \\
\hline
\esp{Cuídenlo. / Cuídense.} & \esp{No esperen.} \\
\hline
\esp{Escríbannos, llámenos, mándennos sus ideas.} & \\
\hline
\end{tabularx}
\end{center}

\courssub{Hypothèse des élèves : démarche inductive}

Regardez attentivement le tableau. Où sont les pronoms (\esp{la}, \esp{se}, \esp{lo}, \esp{nos}) dans la colonne de gauche ? Et dans la colonne de droite ?

\begin{experience}[Formuler l'hypothèse en équipes]
Par groupes de quatre, observez le tableau pendant trois minutes et complétez oralement les phrases suivantes :
\begin{enumerate}
\item À l'impératif \textbf{affirmatif}, les pronoms compléments se placent \ldots{} le verbe et s'écrivent \ldots{} (collés / séparés).
\item À l'impératif \textbf{négatif}, les pronoms compléments se placent \ldots{} le verbe, après le mot \ldots{}
\item Quand on colle un pronom après le verbe, on ajoute souvent un \ldots{} sur la voyelle accentuée (\esp{hiervan} $\rightarrow$ \esp{hiervanla}).
\end{enumerate}
Chaque groupe présente son hypothèse à la classe. Le professeur la note au tableau : elle sera vérifiée et transformée en règle officielle dans la section suivante.
\end{experience}

\begin{aretenir}[Ce que l'observation nous a déjà appris]
\begin{itemize}
\item L'impératif sert à donner un \cle{ordre}, un \cle{conseil} ou une \cle{interdiction} ; le sujet n'est pas exprimé.
\item À l'\cle{affirmatif}, le pronom est \cle{enclitique} : collé après le verbe (\esp{dámelo}, \esp{cuídate}), souvent avec un accent écrit ajouté.
\item Au \cle{négatif}, le pronom est \cle{proclitique} : placé avant le verbe, entre \esp{no} et le verbe (\esp{no me lo des}).
\item C'est presque la même logique qu'en français : « donne-le-moi » / « ne me le donne pas ». La grande différence : en espagnol, on écrit le pronom \emph{collé} au verbe affirmatif, sans trait d'union.
\end{itemize}
\end{aretenir}

\begin{exoresolu}[Repérage et traduction]
Énoncé. Dans la phrase \esp{Recojan la basura y deposítenla en el contenedor}, relevez l'impératif avec pronom, indiquez la place du pronom, puis traduisez la phrase en français.

\tcblower

Solution détaillée.
\begin{enumerate}
\item Les deux verbes \esp{recojan} et \esp{deposítenla} sont sans sujet exprimé et donnent des consignes : ce sont des impératifs affirmatifs (personne \esp{ustedes}, « vous »).
\item L'impératif avec pronom est \esp{deposítenla} = \esp{depositen} + \esp{la}. Le pronom \esp{la} (qui reprend \esp{la basura}) est \textbf{enclitique} : collé après le verbe.
\item Accentuation : \esp{depositen} s'accentue naturellement sur \emph{po} ; en ajoutant \esp{la}, la syllabe accentuée devient la troisième en partant de la fin, donc on écrit un accent : \esp{deposítenla}.
\item Traduction : « Ramassez les ordures et déposez-les dans la poubelle. » Remarquez que le français fait exactement pareil : pronom après le verbe affirmatif, avec trait d'union (« déposez-les »).
\end{enumerate}
\end{exoresolu}

\begin{exercice}[À vous de jouer]
Relevez dans le texte trois autres impératifs avec pronom, classez-les (affirmatif / négatif) et traduisez-les en français. Puis formulez avec vos mots l'hypothèse sur la place des pronoms.
\end{exercice}

\begin{savaistu}[L'impératif des campagnes de santé en Guinée équatoriale]
En Guinée équatoriale, pays voisin du Cameroun et seul pays hispanophone d'Afrique subsaharienne, l'espagnol est langue officielle. Les campagnes de santé publique y utilisent massivement l'impératif, sur les affiches et à la radio : \esp{Lávate las manos} (« Lave-toi les mains »), \esp{Vacúnate} (« Fais-toi vacciner »), \esp{Bebe agua hervida} (« Bois de l'eau bouillie »), \esp{No tires basura en la calle} (« Ne jette pas d'ordures dans la rue »). Remarquez la forme \esp{vacúnate} : l'accent écrit apparaît dès qu'on colle le pronom réfléchi \esp{te} à l'impératif \esp{vacuna}. La grammaire de cette leçon est donc un outil de communication réel, utilisé chaque jour à Malabo et à Bata — et utile à tout Camerounais qui voyage, étudie ou travaille avec le monde hispanophone.
\end{savaistu}

\begin{popfigure}
\begin{tikzpicture}[mindmap, grow cyclic, every node/.style={concept, text=white, align=center, font=\small}, concept color=popblue, level 1/.style={sibling angle=90, level distance=3.2cm}]
\node {La consigna}
  child[concept color=popgreen] { node[align=center]{Salud \\ \esp{Bebe agua} \\ \esp{Lávate las manos}} }
  child[concept color=popgold] { node[align=center]{Medio ambiente \\ \esp{Recicla las botellas} \\ \esp{Planta un árbol}} }
  child[concept color=poporange] { node[align=center]{Imperativo afirmativo \\ pronom collé : \\ \esp{dámelo}} }
  child[concept color=poppink] { node[align=center]{Imperativo negativo \\ pronom avant : \\ \esp{no me lo des}} };
\end{tikzpicture}
\legende{Carte mentale : la consigne, entre santé, environnement et grammaire de l'impératif.}
\end{popfigure}

\textbf{Transition.} Nous avons entendu les consignes, compris le message, repéré les impératifs et formulé notre hypothèse sur la place des pronoms. Il est temps maintenant de vérifier cette hypothèse : comment se forment exactement l'impératif affirmatif et l'impératif négatif à toutes les personnes (\esp{tú}, \esp{vosotros}, \esp{usted}, \esp{ustedes}) ? C'est l'objet de la section suivante.

\coursec{Former l'impératif affirmatif et négatif}

Dans la section précédente, nous avons écouté des consignes de santé et d'écologie : \esp{Bebe agua. No tires basura al suelo. Haga deporte todos los días.} Vous avez remarqué que toutes ces phrases donnent un ordre ou un conseil, et qu'aucune ne contient de sujet exprimé. Elles sont construites avec le \cle{mode impératif}. Il est temps maintenant d'apprendre à former ce mode de manière rigoureuse, à l'affirmatif comme au négatif, pour les quatre personnes qui existent en espagnol : \esp{tú}, \esp{vosotros}, \esp{usted} et \esp{ustedes}.

\courssub{Qu'est-ce que l'impératif ?}

\begin{definition}[L'impératif]
L'\cle{impératif} (\esp{el imperativo}) est le mode verbal de l'\textbf{ordre}, du \textbf{conseil}, de la \textbf{consigne}, de la \textbf{prière} ou de l'\textbf{interdiction}. Sa caractéristique essentielle : le sujet n'est \textbf{jamais exprimé}. En espagnol, il possède quatre personnes :
\begin{itemize}
\item \esp{tú} : on s'adresse à une personne qu'on tutoie ;
\item \esp{vosotros/vosotras} : on s'adresse à plusieurs personnes qu'on tutoie (usage surtout espagnol) ;
\item \esp{usted} : on s'adresse à une personne qu'on vouvoie ;
\item \esp{ustedes} : on s'adresse à plusieurs personnes qu'on vouvoie (en Amérique latine, \esp{ustedes} remplace aussi \esp{vosotros}).
\end{itemize}
\end{definition}

\begin{attention}[Un point capital]
En espagnol, l'impératif \textbf{négatif} n'est pas la forme affirmative précédée de \esp{no}. C'est une forme totalement différente, empruntée au \textbf{subjonctif présent}. C'est la grande différence avec le français, où « parle » et « ne parle pas » utilisent la même forme verbale. Retenez dès maintenant : \esp{¡Habla!} mais \esp{¡No hables!}
\end{attention}

\courssub{Observation : découvrir les formes dans un texte}

\begin{textebox}[Consejos del doctor Mbarga]
Buenos días, jóvenes de Yaoundé. Escuchad con atención. \textbf{Bebed} mucha agua durante el día, sobre todo cuando hace calor. \textbf{Comed} frutas y verduras frescas del mercado. \textbf{No comáis} demasiados dulces. \textbf{Haced} deporte tres veces por semana : corred, jugad al fútbol, caminad. \textbf{No paséis} toda la tarde delante del teléfono. Y tú, joven, \textbf{duerme} ocho horas cada noche, \textbf{no veas} la televisión hasta muy tarde. Señora, \textbf{traiga} a sus hijos al centro de salud para la vacunación y \textbf{no espere} a que estén enfermos. Señores, \textbf{protejan} el medio ambiente : \textbf{planten} árboles, \textbf{no quemen} la basura, \textbf{reciclen} el plástico. ¡La salud y la naturaleza son un tesoro!
\end{textebox}

\textbf{Glossaire :} \esp{escuchad} = écoutez ; \esp{los dulces} = les sucreries ; \esp{delante de} = devant ; \esp{la vacunación} = la vaccination ; \esp{quemar} = brûler ; \esp{reciclar} = recycler ; \esp{un tesoro} = un trésor.

Observons les verbes en gras. À l'\textbf{affirmatif}, on trouve \esp{bebed}, \esp{comed}, \esp{haced}, \esp{duerme}, \esp{traiga}, \esp{planten}. Au \textbf{négatif}, on trouve \esp{no comáis}, \textbf{no paséis}, \esp{no veas}, \esp{no espere}, \esp{no quemen}, \esp{no reciclen}. Deux familles de formes se dessinent clairement. Dégagéns la règle.

\courssub{L'impératif affirmatif de \esp{tú}}

\begin{propriete}[Formation de l'impératif affirmatif de tú]
L'impératif affirmatif de \esp{tú} est identique à la \textbf{3\up{e} personne du singulier du présent de l'indicatif} :
\begin{itemize}
\item \esp{hablar} $\rightarrow$ \esp{habla} (il parle $\rightarrow$ parle !)
\item \esp{comer} $\rightarrow$ \esp{come} (il mange $\rightarrow$ mange !)
\item \esp{escribir} $\rightarrow$ \esp{escribe} (il écrit $\rightarrow$ écris !)
\end{itemize}
Les verbes à diphtongaison suivent la même logique : \esp{duerme} (dormir), \esp{cierra} (fermer), \esp{piensa} (penser), \esp{vuelve} (revenir).
\end{propriete}

\begin{exemplebox}[Exemples traduits]
\esp{Habla más despacio, por favor.} « Parle plus lentement, s'il te plaît. »\\
\esp{Come verduras todos los días.} « Mange des légumes tous les jours. »\\
\esp{Escribe la carta ahora.} « Écris la lettre maintenant. »\\
\esp{Duerme ocho horas.} « Dors huit heures. »\\
\esp{Cierra la puerta.} « Ferme la porte. »
\end{exemplebox}

\textbf{Les huit impératifs irréguliers de \esp{tú}.} Huit verbes très fréquents possèdent une forme courte irrégulière qu'il faut apprendre par cœur :

\begin{center}
\begin{tabularx}{\linewidth}{|X|X|X|}
\hline
\rowcolor{popblueL} Infinitif & Impératif \esp{tú} & Exemple traduit \\
\hline
\esp{decir} & \esp{di} & \esp{Di la verdad.} « Dis la vérité. » \\
\hline
\esp{hacer} & \esp{haz} & \esp{Haz los deberes.} « Fais tes devoirs. » \\
\hline
\esp{ir} & \esp{ve} & \esp{Ve al mercado.} « Va au marché. » \\
\hline
\esp{poner} & \esp{pon} & \esp{Pon la mesa.} « Mets la table. » \\
\hline
\esp{salir} & \esp{sal} & \esp{Sal de aquí.} « Sors d'ici. » \\
\hline
\esp{ser} & \esp{sé} & \esp{Sé bueno.} « Sois bon. » \\
\hline
\esp{tener} & \esp{ten} & \esp{Ten paciencia.} « Aie patience. » \\
\hline
\esp{venir} & \esp{ven} & \esp{Ven aquí.} « Viens ici. » \\
\hline
\end{tabularx}
\end{center}

\begin{savaistu}[Une phrase mnémotechnique]
Les huit impératifs irréguliers de \esp{tú} se mémorisent facilement avec la phrase : \textbf{« Di haz ve, pon sal, sé ten ven »}. Répétez-la à voix haute trois fois : elle contient, dans l'ordre, \esp{decir}, \esp{hacer}, \esp{ir}, \esp{poner}, \esp{salir}, \esp{ser}, \esp{tener}, \esp{venir}. Beaucoup d'élèves espagnols l'apprennent ainsi !
\end{savaistu}

\courssub{L'impératif affirmatif de \esp{vosotros}}

\begin{propriete}[Formation de l'impératif de vosotros]
L'impératif affirmatif de \esp{vosotros} se forme en remplaçant le \textbf{-r final de l'infinitif par -d} :
\begin{itemize}
\item \esp{hablar} $\rightarrow$ \esp{hablad}
\item \esp{comer} $\rightarrow$ \esp{comed}
\item \esp{escribir} $\rightarrow$ \esp{escribid}
\end{itemize}
Cette formation est \textbf{absolument régulière pour tous les verbes}, même les irréguliers : \esp{ir} $\rightarrow$ \esp{id}, \esp{tener} $\rightarrow$ \esp{tened}, \esp{hacer} $\rightarrow$ \esp{haced}, \esp{venir} $\rightarrow$ \esp{venid}.
\end{propriete}

\begin{exemplebox}[Exemples traduits]
\esp{Hablad más alto, chicos.} « Parlez plus fort, les garçons. »\\
\esp{Comed despacio.} « Mangez lentement. »\\
\esp{Id a la farmacia.} « Allez à la pharmacie. »\\
\esp{Tened cuidado con el fuego.} « Faites attention au feu. »
\end{exemplebox}

\courssub{L'impératif de \esp{usted} et \esp{ustedes} : le subjonctif}

\begin{propriete}[Usted et ustedes : formes du subjonctif présent]
Pour s'adresser à une personne qu'on vouvoie (\esp{usted}) ou à plusieurs (\esp{ustedes}), l'impératif affirmatif emprunte ses formes au \textbf{subjonctif présent} : on échange les voyelles thématiques.
\begin{itemize}
\item Verbes en \textbf{-ar} : voyelle \textbf{e}. \esp{hablar} $\rightarrow$ \esp{hable} (usted), \esp{hablen} (ustedes).
\item Verbes en \textbf{-er} et \textbf{-ir} : voyelle \textbf{a}. \esp{comer} $\rightarrow$ \esp{coma}, \esp{coman} ; \esp{escribir} $\rightarrow$ \esp{escriba}, \esp{escriban}.
\end{itemize}
Verbes irréguliers au subjonctif : \esp{ir} $\rightarrow$ \esp{vaya}, \esp{vayan} ; \esp{ser} $\rightarrow$ \esp{sea}, \esp{sean} ; \esp{tener} $\rightarrow$ \esp{tenga}, \esp{tengan} ; \esp{hacer} $\rightarrow$ \esp{haga}, \esp{hagan} ; \esp{estar} $\rightarrow$ \esp{esté}, \esp{estén}.
\end{propriete}

\begin{exemplebox}[Exemples traduits (vouvoiement)]
\esp{Tenga cuidado, señora.} « Faites attention, madame. »\\
\esp{Pase, por favor.} « Entrez, je vous en prie. »\\
\esp{No fume aquí.} « Ne fumez pas ici. »\\
\esp{Escuchen las instrucciones.} « Écoutez les instructions. »\\
\esp{Vayan al centro de salud.} « Allez au centre de santé. »
\end{exemplebox}

\courssub{L'impératif négatif : \esp{no} + subjonctif présent}

\begin{propriete}[La règle d'or de l'impératif négatif]
L'impératif négatif emprunte \textbf{TOUTES} ses formes au \textbf{subjonctif présent}, précédées de \esp{no}, à \textbf{toutes les personnes} :
\begin{itemize}
\item \esp{tú} : \esp{no hables}, \esp{no comas}, \esp{no escribas}, \esp{no vayas}
\item \esp{vosotros} : \esp{no habléis}, \esp{no comáis}, \esp{no escribáis}, \esp{no vayáis}
\item \esp{usted} : \esp{no hable}, \esp{no coma}, \esp{no escriba}, \esp{no vaya}
\item \esp{ustedes} : \esp{no hablen}, \esp{no coman}, \esp{no escriban}, \esp{no vayan}
\end{itemize}
Même les verbes dont l'impératif affirmatif de \esp{tú} est irrégulier redeviennent « normaux » au négatif : \esp{¡Di la verdad!} mais \esp{¡No digas mentiras!} ; \esp{¡Haz los deberes!} mais \esp{¡No hagas ruido!}
\end{propriete}

\begin{attention}[L'erreur numéro un des francophones]
En français, « parle » et « ne parle pas » ont la même forme. Les élèves francophones transposent ce réflexe et disent \esp{no habla}, \esp{no come}, \esp{no ve}. C'est \textbf{faux} ! Le négatif n'est \textbf{jamais} la forme de l'indicatif. Dites : \esp{no hables} (ne parle pas), \esp{no comas} (ne mange pas), \esp{no vayas} (ne va pas). Astuce : à l'affirmatif, \esp{tú} prend la voyelle de l'indicatif ; au négatif, il \textbf{change de voyelle}.
\end{attention}

\begin{exemplebox}[Affirmatif contre négatif : exemples traduits]
\esp{Haz los deberes.} « Fais tes devoirs. » $\rightarrow$ \esp{No hagas los deberes ahora.} « Ne fais pas tes devoirs maintenant. »\\
\esp{Ve al río.} « Va à la rivière. » $\rightarrow$ \esp{No vayas al río solo.} « Ne va pas seul à la rivière. »\\
\esp{Pon la mesa.} « Mets la table. » $\rightarrow$ \esp{No pongas los pies en la mesa.} « Ne mets pas les pieds sur la table. »\\
\esp{Sal ahora.} « Sors maintenant. » $\rightarrow$ \esp{No salgas de noche.} « Ne sors pas la nuit. »
\end{exemplebox}

\courssub{Les changements orthographiques}

\begin{propriete}[Préserver le son de l'infinitif]
Au subjonctif (donc à l'impératif négatif et aux formes de \esp{usted}), certains verbes modifient leur orthographe pour conserver le son de l'infinitif :
\begin{itemize}
\item \textbf{c $\rightarrow$ z} devant a : \esp{empezar} $\rightarrow$ \esp{no empieces}, \esp{no empiece}. (Le son [s] de \esp{empezar} s'écrit z devant a.)
\item \textbf{g $\rightarrow$ gu} devant e : \esp{llegar} $\rightarrow$ \esp{no llegues}, \esp{no llegue}. (Le gu garde le son [g] dur.)
\item \textbf{c $\rightarrow$ qu} devant e : \esp{practicar} $\rightarrow$ \esp{practica} (affirmatif) mais \esp{no practiques}, \esp{no practique}. (Le qu garde le son [k].)
\item \textbf{z $\rightarrow$ c} devant e : \esp{cruzar} $\rightarrow$ \esp{no cruces}, \esp{no cruce}.
\item \textbf{gu $\rightarrow$ gü} : \esp{averiguar} $\rightarrow$ \esp{no averigües}.
\end{itemize}
\end{propriete}

\begin{exemplebox}[Exemples traduits]
\esp{Empieza el trabajo.} « Commence le travail. » $\rightarrow$ \esp{No empieces tarde.} « Ne commence pas tard. »\\
\esp{Llega temprano.} « Arrive tôt. » $\rightarrow$ \esp{No llegues tarde a clase.} « N'arrive pas en retard en classe. »\\
\esp{Practica deporte.} « Fais du sport. » $\rightarrow$ \esp{No practiques bajo la lluvia.} « Ne fais pas de sport sous la pluie. »\\
\esp{Cruza con cuidado.} « Traverse prudemment. » $\rightarrow$ \esp{No cruces la calle sin mirar.} « Ne traverse pas la rue sans regarder. »
\end{exemplebox}

\courssub{Tableaux récapitulatifs}

\begin{center}
\begin{tabularx}{\linewidth}{|l|X|X|X|X|X|}
\hline
\rowcolor{popblueL} Verbe & \esp{tú} (+) & \esp{tú} (--) & \esp{vosotros} (+) & \esp{usted} (+) & \esp{ustedes} (+) \\
\hline
\esp{hablar} & \esp{habla} & \esp{no hables} & \esp{hablad} & \esp{hable} & \esp{hablen} \\
\hline
\esp{comer} & \esp{come} & \esp{no comas} & \esp{comed} & \esp{coma} & \esp{coman} \\
\hline
\esp{escribir} & \esp{escribe} & \esp{no escribas} & \esp{escribid} & \esp{escriba} & \esp{escriban} \\
\hline
\esp{ir} & \esp{ve} & \esp{no vayas} & \esp{id} & \esp{vaya} & \esp{vayan} \\
\hline
\esp{tener} & \esp{ten} & \esp{no tengas} & \esp{tened} & \esp{tenga} & \esp{tengan} \\
\hline
\end{tabularx}
\end{center}

Pour le négatif des trois dernières personnes, il suffit d'ajouter \esp{no} : \esp{no habléis}, \esp{no hable}, \esp{no hablen}, etc.

\begin{center}
\begin{tabularx}{\linewidth}{|X|X|X|}
\hline
\rowcolor{popgreenL} Français & Español & Remarque \\
\hline
Parle ! / Ne parle pas ! & \esp{¡Habla! / ¡No hables!} & changement de voyelle au négatif \\
\hline
Mangez ! (tutoiement pluriel) & \esp{¡Comed!} & infinitif --r + d \\
\hline
Faites attention ! (vouvoiement) & \esp{¡Tenga cuidado!} & forme de subjonctif \\
\hline
N'allez pas là-bas ! (vouvoiement pluriel) & \esp{¡No vayan allí!} & \esp{ir} : subjonctif irrégulier \\
\hline
Fais tes devoirs ! & \esp{¡Haz los deberes!} & impératif irrégulier de \esp{hacer} \\
\hline
\end{tabularx}
\end{center}

\begin{popfigure}
\begin{center}
\begin{tikzpicture}[
  pers/.style={rectangle, rounded corners, draw=popblue, fill=popblueL, font=\bfseries\small, minimum width=1.7cm, minimum height=0.8cm},
  aff/.style={rectangle, rounded corners, draw=popgreen, fill=popgreenL, font=\small, align=center, minimum width=3.2cm, minimum height=1.5cm},
  neg/.style={rectangle, rounded corners, draw=poporange, fill=poporangeL, font=\small, align=center, minimum width=3.2cm, minimum height=1.5cm},
  head/.style={font=\bfseries\small, text=popdark}
]
\node[head] at (3.2,3.6) {Affirmatif};
\node[head] at (7.2,3.6) {Négatif};
\node[pers] at (0,2.4) {tú};
\node[pers] at (0,0.6) {vosotros};
\node[pers] at (0,-1.2) {usted};
\node[pers] at (0,-3.0) {ustedes};
\node[aff, align=center] at (3.2,2.4){habla\\ come\\ escribe};
\node[neg, align=center] at (7.2,2.4){no hables\\ no comas\\ no escribas};
\node[aff, align=center] at (3.2,0.6){hablad\\ comed\\ escribid};
\node[neg, align=center] at (7.2,0.6){no habléis\\ no comáis\\ no escribáis};
\node[aff, align=center] at (3.2,-1.2){hable\\ coma\\ escriba};
\node[neg, align=center] at (7.2,-1.2){no hable\\ no coma\\ no escriba};
\node[aff, align=center] at (3.2,-3.0){hablen\\ coman\\ escriban};
\node[neg, align=center] at (7.2,-3.0){no hablen\\ no coman\\ no escriban};
\draw[popline, thick] (1.6,3.1) -- (8.8,3.1);
\end{tikzpicture}
\end{center}
\legende{Tableau à double entrée : les quatre personnes de l'impératif à l'affirmatif (vert) et au négatif (orange) pour \esp{hablar}, \esp{comer}, \esp{escribir}. Au négatif, toutes les formes viennent du subjonctif présent.}
\end{popfigure}

\begin{aretenir}[L'essentiel en quatre phrases]
\begin{enumerate}
\item \esp{tú} affirmatif = 3\up{e} personne du singulier du présent (\esp{habla}, \esp{come}, \esp{escribe}) + 8 irréguliers : \esp{di, haz, ve, pon, sal, sé, ten, ven}.
\item \esp{vosotros} affirmatif = infinitif --r + d (\esp{hablad}, \esp{comed}, \esp{escribid}), toujours régulier.
\item \esp{usted} et \esp{ustedes} = formes du subjonctif présent (\esp{hable}, \esp{hablen}).
\item Tout le négatif = \esp{no} + subjonctif présent (\esp{no hables}, \esp{no comáis}, \esp{no vaya}).
\end{enumerate}
\end{aretenir}

\begin{exoresolu}[Traduire en espagnol]
Énoncé. Traduisez en espagnol : 1. « Bois de l'eau, mon enfant. » 2. « Ne mangez pas trop de sucreries, les enfants. » 3. « Faites du sport, madame. » 4. « Ne va pas seul à la rivière. » 5. « N'arrivez pas en retard, messieurs. »
\tcblower
Solution détaillée.\\
1. On tutoie une personne $\rightarrow$ impératif affirmatif de \esp{tú} : 3\up{e} personne du singulier de \esp{beber}. \esp{Bebe agua, hijo mío.}\\
2. On s'adresse à plusieurs personnes qu'on tutoie, au négatif $\rightarrow$ \esp{no} + subjonctif de \esp{vosotros} : \esp{No comáis demasiados dulces, niños.}\\
3. On vouvoie une personne, à l'affirmatif $\rightarrow$ subjonctif de \esp{hacer} : \esp{Haga deporte, señora.}\\
4. Tutoiement au négatif $\rightarrow$ \esp{no} + subjonctif de \esp{tú} ; \esp{ir} est irrégulier : \esp{No vayas al río solo.}\\
5. Vouvoiement pluriel au négatif ; \esp{llegar} a un changement orthographique g $\rightarrow$ gu : \esp{No lleguen tarde, señores.}
\end{exoresolu}

\begin{exercice}[À vous de conjuguer]
Mettez les verbes entre parenthèses à l'impératif demandé, puis traduisez la phrase en français :
\begin{enumerate}
\item \esp{(cerrar, tú, +) \ldots\ la ventana, hace frío.}
\item \esp{(fumar, usted, --) No \ldots\ en el hospital.}
\item \esp{(reciclar, vosotros, +) \ldots\ el plástico, chicos.}
\item \esp{(venir, tú, +) \ldots\ conmigo al mercado.}
\item \esp{(empezar, ustedes, --) No \ldots\ sin mí.}
\end{enumerate}
\end{exercice}

\begin{corrige}[Exercice : À vous de conjuguer]
1. \esp{Cierra la ventana, hace frío.} « Ferme la fenêtre, il fait froid. » (3\up{e} pers. du singulier, diphtongaison e $\rightarrow$ ie.)\\
2. \esp{No fume en el hospital.} « Ne fumez pas à l'hôpital. » (subjonctif de \esp{usted}.)\\
3. \esp{Reciclad el plástico, chicos.} « Recyclez le plastique, les enfants. » (infinitif --r + d.)\\
4. \esp{Ven conmigo al mercado.} « Viens avec moi au marché. » (impératif irrégulier de \esp{venir}.)\\
5. \esp{No empiecen sin mí.} « Ne commencez pas sans moi. » (subjonctif de \esp{ustedes}, changement c $\rightarrow$ z.)
\end{corrige}

Maintenant que les quatre personnes de l'impératif n'ont plus de secret pour vous, une question se pose : où placer les pronoms quand on dit « donne-le-moi » ou « ne me le donne pas » ? C'est l'objet de la section suivante, consacrée à l'enclise et à la proclise.

\coursec{Les pronoms compléments avec l'impératif : enclise et proclise}

Dans la section précédente, nous avons appris à \cle{former} l'impératif affirmatif (\esp{compra}, \esp{haced}, \esp{escuche}) et négatif (\esp{no compres}, \esp{no hagáis}, \esp{no escuchen}). Mais dans la vraie vie, on donne rarement des ordres « à vide » : on dit « donne-\textbf{le}-moi », « lève-\textbf{toi} », « ne \textbf{les} jette pas ». Autrement dit, l'impératif adore les \cle{pronoms compléments}. Or leur place en espagnol obéit à une mécanique très précise, différente du français, et c'est l'un des points les plus rentables — et les plus piégeux — de l'épreuve de traduction au Bac. Observons d'abord.

\courssub{Observation : un texte pour voir la règle en action}

\begin{textebox}[Consejos de la doctora Belén (consulta de salud)]
— Buenos días, señor Etoundi. Veo que tiene usted el colesterol alto. Voy a darle unos consejos muy sencillos. Primero: \textbf{haga} deporte cada día, pero \textbf{hágalo} por la mañana, no por la noche. Coma más verduras: \textbf{cómprelas} en el mercado de Mvog-Mbi, son frescas y baratas. ¿El aceite de palma? \textbf{Redúzcalo}, no \textbf{lo} elimine del todo. ¿Y la sal? \textbf{Quítela} poco a poco, pero \textbf{no se la} quite a los niños de golpe. ¿Tiene usted la receta? \textbf{Démela}, por favor, quiero revisarla. \textbf{No me la} devuelva firmada sin leerla. Ah, y una cosa más: \textbf{acuéstese} temprano y \textbf{no se acueste} nunca después de medianoche. ¿Entendido? \textbf{Repítamelo} todo, quiero estar segura.
\end{textebox}

\textbf{Glossaire :} \esp{el colesterol} : le cholestérol ; \esp{las verduras} : les légumes ; \esp{el aceite de palma} : l'huile de palme ; \esp{la sal} : le sel ; \esp{quitar} : enlever, retirer ; \esp{de golpe} : d'un coup ; \esp{la receta} : l'ordonnance ; \esp{devolver} : rendre ; \esp{acostarse} : se coucher ; \esp{después de medianoche} : après minuit.

\textbf{Questions d'observation :}
\begin{enumerate}
\item Relevez tous les impératifs affirmatifs du texte. Où se trouvent les pronoms : avant ou après le verbe ? Sont-ils soudés ou séparés ?
\item Relevez les impératifs négatifs. Où sont les pronoms cette fois ?
\item Observez \esp{hágalo}, \esp{cómprelas}, \esp{démela}, \esp{repítamelo} : que remarquez-vous sur les accents ?
\end{enumerate}

\textbf{Constat :} à l'affirmatif, le pronom est \emph{après} le verbe et \emph{collé} à lui (\esp{hágalo}) ; au négatif, il est \emph{avant} le verbe et \emph{séparé} (\esp{no lo elimine}). C'est toute la leçon.

\courssub{La règle d'or : enclise à l'affirmatif, proclise au négatif}

\begin{propriete}[Enclise obligatoire à l'impératif affirmatif]
À l'impératif \cle{affirmatif}, le ou les pronoms compléments se placent \textbf{après} le verbe et s'écrivent \textbf{soudés} à lui, en un seul mot : c'est l'\cle{enclise}.

\esp{Compra el pan.} → \esp{Cómpralo.} « Achète-le. »

\esp{Dame el libro.} → \esp{Dámelo.} « Donne-le-moi. »

On n'écrit jamais \esp{dame lo} en deux mots.
\end{propriete}

\begin{propriete}[Proclise obligatoire à l'impératif négatif]
À l'impératif \cle{négatif}, le ou les pronoms compléments se placent \textbf{avant} le verbe, \textbf{après} la négation \esp{no}, et restent des mots \textbf{séparés} : c'est la \cle{proclise}.

\esp{No compres el pan.} → \esp{No lo compres.} « Ne l'achète pas. »

\esp{No me des el libro.} → \esp{No me lo des.} « Ne me le donne pas. »
\end{propriete}

\begin{exemplebox}[Les deux mécanismes face à face]
\begin{itemize}
\item \esp{Dámelo.} « Donne-le-moi. » / \esp{No me lo des.} « Ne me le donne pas. »
\item \esp{Díselo a María.} « Dis-le à Marie. » / \esp{No se lo digas a María.} « Ne le dis pas à Marie. »
\item \esp{Levántate.} « Lève-toi. » / \esp{No te levantes.} « Ne te lève pas. »
\item \esp{Acuéstense.} « Couchez-vous. » / \esp{No se acuesten.} « Ne vous couchez pas. »
\item \esp{Apaga la luz. → Apágala.} « Éteins-la. » / \esp{No la apagues.} « Ne l'éteins pas. »
\end{itemize}
\end{exemplebox}

\begin{popfigure}
\begin{tikzpicture}[font=\small, >=Stealth, node distance=0.6cm]
\node[draw=popblue, fill=popblueL, rounded corners, align=center, minimum width=3.2cm] (imp) {\textbf{Impératif}\\ \textbf{+ pronom}};
\node[draw=popgreen, fill=popgreenL, rounded corners, align=center, right=2.2cm of imp, yshift=1.1cm, minimum width=4.6cm] (aff) {\textbf{AFFIRMATIF : enclise}\\ verbe + pronom soudé\\ \esp{dámelo}, \esp{levántate}};
\node[draw=poppink, fill=poppinkL, rounded corners, align=center, right=2.2cm of imp, yshift=-1.1cm, minimum width=4.6cm] (neg) {\textbf{NÉGATIF : proclise}\\ no + pronom + verbe\\ \esp{no me lo des}, \esp{no te levantes}};
\draw[->, very thick, popgreen] (imp) -- (aff) node[midway, above, sloped, popdark]{affirmatif ?};
\draw[->, very thick, poppink] (imp) -- (neg) node[midway, below, sloped, popdark]{négatif ?};
\end{tikzpicture}
\legende{Les deux destins du pronom avec l'impératif : soudé après (affirmatif), séparé avant (négatif).}
\end{popfigure}

\begin{popfigure}
\begin{tikzpicture}[font=\small, >=Stealth, node distance=0.55cm and 1.4cm]
\node[draw=popblue, fill=popblueL, rounded corners, align=center] (start) {Phrase à traduire\\ avec impératif + pronom};
\node[draw=popgold, fill=popgoldL, diamond, aspect=2.2, align=center, below=of start] (q1) {Affirmatif ?};
\node[draw=popgreen, fill=popgreenL, rounded corners, align=center, right=2.6cm of q1] (oui) {\textbf{ENCLISE}\\ verbe+pronom en un mot\\ + accent écrit si besoin\\ \esp{cómpramelo}};
\node[draw=poppink, fill=poppinkL, rounded corners, align=center, below=of q1] (non) {\textbf{PROCLISE}\\ no + pronom(s) + verbe\\ mots séparés\\ \esp{no me lo compres}};
\node[draw=popmuted, fill=popcream, rounded corners, align=center, below=of non] (fin) {Vérifier : ordre COI puis COD\\ (\esp{me lo}, \esp{se la})};
\draw[->, thick] (start) -- (q1);
\draw[->, thick, popgreen] (q1) -- node[above]{oui} (oui);
\draw[->, thick, poppink] (q1) -- node[right]{non} (non);
\draw[->, thick, popmuted] (non) -- (fin);
\draw[->, thick, popmuted] (oui.south) |- (fin.east);
\end{tikzpicture}
\legende{Organigramme de décision : où placer le pronom avec l'impératif ?}
\end{popfigure}

\courssub{Combiner deux pronoms : l'ordre indirect avant direct}

Quand l'impératif porte deux pronoms (un objet direct et un objet indirect), l'espagnol impose un ordre fixe : \cle{COI avant COD}.

\begin{propriete}[Ordre des pronoms combinés]
Le pronom d'objet \textbf{indirect} (\esp{me, te, se, nos, os, se}) précède toujours le pronom d'objet \textbf{direct} (\esp{lo, la, los, las}) : \esp{me lo}, \esp{te la}, \esp{se lo}, \esp{nos las}. À l'affirmatif, le bloc entier se soude au verbe ; au négatif, le bloc entier précède le verbe.
\end{propriete}

\begin{propriete}[Le cas spécial : le/les deviennent se]
Devant \esp{lo, la, los, las}, les pronoms indirects \esp{le} et \esp{les} se transforment en \esp{se} (pour éviter la cacophonie \esp{le lo}). On écrit donc \esp{díselo} « dis-le-lui » et jamais \esp{dile lo} ; \esp{no se lo digas} « ne le lui dis pas ».
\end{propriete}

\begin{center}
\begin{tabularx}{\linewidth}{|X|X|X|}
\hline
\rowcolor{popblueL} \textbf{Français} & \textbf{Español} & \textbf{Remarque} \\
\hline
Donne-le-moi. & \esp{Dámelo.} & COI (\esp{me}) + COD (\esp{lo}) : ordre inverse du français affirmatif \\
\hline
Ne me le donne pas. & \esp{No me lo des.} & même ordre qu'en français au négatif \\
\hline
Dis-le-lui. & \esp{Díselo.} & \esp{le} → \esp{se} devant \esp{lo} \\
\hline
Ne les leur envoyez pas. & \esp{No se los envíen.} & \esp{les} → \esp{se} devant \esp{los} \\
\hline
Apporte-nous-les. & \esp{Tráenoslas.} & \esp{nos} + \esp{las}, soudés, accent sur \esp{tra-} \\
\hline
\end{tabularx}
\end{center}

\begin{attention}[Le piège majeur du Bac : l'ordre inversé à l'affirmatif]
En français, à l'impératif affirmatif, on dit « donne-\textbf{le}-moi » : le COD (\emph{le}) vient \emph{avant} le COI (\emph{moi}). En espagnol, c'est l'inverse : \esp{dá-\textbf{me}-\textbf{lo}}. Traduire « donne-le-moi » par \esp{dalo me} ou par \esp{damelo} sans accent est une faute classique. Retenez : en espagnol, c'est \textbf{toujours} la personne d'abord, la chose ensuite — \esp{me lo} = « moi, ça ».
\end{attention}

\courssub{L'accentuation écrite liée à l'enclise}

\begin{propriete}[L'enclise déplace l'accent tonique]
Quand on soude un ou deux pronoms au verbe, le mot s'allonge mais l'accent tonique \textbf{reste sur la même voyelle du verbe}. Pour le conserver, il faut souvent ajouter un \cle{accent écrit} :

\esp{DA} → \esp{DAme} → \esp{DÁmelo}

\esp{COMpra} → \esp{COMprala} → \esp{CÓmpramela}

\esp{exPLIca} → \esp{exPLIcame} → \esp{exPLÍcamelo}
\end{propriete}

\textbf{Comment savoir s'il faut un accent ?} Appliquez les règles générales d'accentuation au mot complet obtenu. En pratique :
\begin{itemize}
\item Un seul pronom ajouté : l'accent écrit n'est nécessaire que si le verbe est long (plus d'une syllabe avant la terminaison) : \esp{dame} (pas d'accent) mais \esp{levántate}, \esp{cómprala}, \esp{escúchame}.
\item Deux pronoms ajoutés : l'accent écrit est presque toujours obligatoire : \esp{dámelo}, \esp{díselo}, \esp{tráenoslas}, \esp{explícamelo}.
\item Formes d'\esp{usted/ustedes} : l'accent du verbe (forme de subjonctif) est déjà marqué, il ne bouge pas : \esp{dígamelo}, \esp{háganlo}, \esp{acuéstense}.
\end{itemize}

\begin{exemplebox}[La chaîne de transformation]
\begin{itemize}
\item \esp{da} → \esp{dame} « donne-moi » → \esp{dámelo} « donne-le-moi »
\item \esp{compra} → \esp{cómprala} « achète-la » → \esp{cómpramela} « achète-la-moi »
\item \esp{diga} → \esp{dígale} « dites-lui » → \esp{dígaselo} « dites-le-lui »
\item \esp{explica} → \esp{explícame} « explique-moi » → \esp{explícamelo} « explique-le-moi »
\end{itemize}
\end{exemplebox}

\courssub{Cas particuliers : vosotros et verbes pronominaux}

\begin{propriete}[Vosotros : la chute du -d devant os]
À la 2\textsuperscript{e} personne du pluriel (\esp{vosotros}), l'impératif affirmatif se termine en \textbf{-d} (\esp{levantad}, \esp{dad}). Devant le pronom réfléchi \esp{os}, cette \textbf{-d tombe} : \esp{levantaos} « levez-vous », \esp{callaos} « taisez-vous », \esp{daoslo} « donnez-le-vous ».

\textbf{Exception :} le verbe \esp{ir} garde son -d : \esp{idos} « allez-vous-en » (et non \esp{íos}).
\end{propriete}

\begin{propriete}[Verbes pronominaux à l'impératif]
Avec les verbes pronominaux, le pronom réfléchi suit exactement la même règle : soudé après à l'affirmatif, séparé avant au négatif.

\esp{Levántate.} « Lève-toi. » / \esp{No te levantes.} « Ne te lève pas. »

\esp{Acuéstense.} « Couchez-vous. » / \esp{No se acuesten.} « Ne vous couchez pas. »

\esp{Lavaos las manos.} « Lavez-vous les mains. » / \esp{No os lavéis con agua fría.} « Ne vous lavez pas à l'eau froide. »
\end{propriete}

\begin{center}
\begin{tabularx}{\linewidth}{|X|X|X|}
\hline
\rowcolor{popblueL} \textbf{Personne} & \textbf{Affirmatif (enclise)} & \textbf{Négatif (proclise)} \\
\hline
tú & \esp{levántate} & \esp{no te levantes} \\
\hline
usted & \esp{levántese} & \esp{no se levante} \\
\hline
vosotros & \esp{levantaos} & \esp{no os levantéis} \\
\hline
ustedes & \esp{levántense} & \esp{no se levanten} \\
\hline
\end{tabularx}
\end{center}

\courssub{Méthode et entraînement à la traduction}

\begin{methode}[Traduire un impératif avec pronoms en cinq pas]
\begin{enumerate}
\item \textbf{Affirmatif ou négatif ?} Cherchez le « ne... pas » français.
\item \textbf{Formez l'impératif} à la bonne personne (\esp{tú}, \esp{usted}, \esp{vosotros}, \esp{ustedes}) : présent d'indicatif sans -s pour l'affirmatif de \esp{tú} ; subjonctif pour le négatif et pour \esp{usted/ustedes}.
\item \textbf{Identifiez les pronoms} : qui reçoit ? (COI : \esp{me, te, le}→\esp{se}, \esp{nos}) — quoi ? (COD : \esp{lo, la, los, las}).
\item \textbf{Placez selon la règle} : affirmatif → soudés après, COI + COD ; négatif → séparés avant, après \esp{no}.
\item \textbf{Vérifiez l'accent écrit} : l'accent tonique du verbe doit être conservé.
\end{enumerate}
\end{methode}

\begin{exoresolu}[Thème : traduisez en espagnol]
Énoncé. Traduisez : 1) « Donne-le-moi. » 2) « Ne me le donne pas. » 3) « Lavez-vous les mains, les enfants ! » 4) « Ne le lui dites pas. » 5) « Explique-nous-la, s'il te plaît. » (la = la règle, \esp{la regla})

\tcblower

Solution détaillée.

1) « Donne-le-moi. » Affirmatif → enclise. Impératif de \esp{dar} à \esp{tú} : \esp{da}. COI \esp{me} + COD \esp{lo} soudés : \esp{damelo}. L'accent tonique reste sur \esp{da-} : il faut l'accent écrit. \textbf{\esp{Dámelo.}}

2) « Ne me le donne pas. » Négatif → proclise + subjonctif : \esp{no des}. Pronoms avant le verbe, après \esp{no} : pas d'accent nouveau. \textbf{\esp{No me lo des.}}

3) « Lavez-vous les mains, les enfants ! » \esp{vosotros} affirmatif : \esp{lavad} ; chute du -d devant \esp{os} : \esp{lavaos}. Accent écrit pour conserver la tonique sur \esp{la-} : \textbf{\esp{¡Lavaos las manos, niños!}}

4) « Ne le lui dites pas. » Négatif, \esp{usted} : \esp{no diga}. \esp{le} → \esp{se} devant \esp{lo} : \textbf{\esp{No se lo diga.}}

5) « Explique-nous-la. » Affirmatif : \esp{explica} + \esp{nos} + \esp{la} (\esp{la regla}, féminin) : \esp{explicanosla} → accent écrit obligatoire : \textbf{\esp{Explícanosla, por favor.}}
\end{exoresolu}

\begin{attention}[Erreurs fréquentes des francophones]
\begin{itemize}
\item Écrire \esp{me lo des} sans \esp{no} : au négatif, la particule \esp{no} est \textbf{obligatoire} et le pronom se place entre \esp{no} et le verbe : \esp{no me lo des}.
\item Écrire \esp{dame lo} en deux mots : à l'affirmatif, la soudure est \textbf{obligatoire} : \esp{dámelo}.
\item Oublier l'accent écrit : \esp{damelo} au lieu de \esp{dámelo} change la prononciation (on dirait \esp{daMElo}) : faute d'orthographe sanctionnée.
\item Garder l'ordre français « le-moi » : \esp{dalo me} n'existe pas ; toujours COI puis COD.
\item Oublier la transformation \esp{le} → \esp{se} : \esp{dile lo} est incorrect, on dit \esp{díselo}.
\end{itemize}
\end{attention}

\begin{exercice}[À votre tour : thème et version]
\textbf{A. Thème.} Traduisez en espagnol : 1) « Jette-le à la poubelle ! » (\esp{tirar}) 2) « Ne les jetez pas par terre. » (\esp{ustedes}) 3) « Lève-toi et fais du sport. » 4) « Ne te couche pas tard. » 5) « Apporte-la-nous demain. » (la = l'ordonnance, \esp{la receta})

\textbf{B. Version.} Traduisez en français : 1) \esp{Cómpraselo a tu hermana.} 2) \esp{No se lo compres.} 3) \esp{Recicladlo todo.} 4) \esp{No os acostéis sin lavaros los dientes.}
\end{exercice}

\begin{corrige}[Exercice : thème et version]
\textbf{A.} 1) \esp{¡Tíralo a la papelera!} 2) \esp{No los tiren al suelo.} 3) \esp{Levántate y haz deporte.} 4) \esp{No te acuestes tarde.} 5) \esp{Tráenosla mañana.}

\textbf{B.} 1) « Achète-le à ta sœur. » (\esp{compra} + \esp{se} + \esp{lo} ; \esp{se} = « à elle », précisé par \esp{a tu hermana}) 2) « Ne le lui achète pas. » 3) « Recyclez-le tout entier. » 4) « Ne vous couchez pas sans vous laver les dents. »
\end{corrige}

\begin{aretenir}[L'essentiel de la section]
\begin{itemize}
\item Impératif \textbf{affirmatif} : pronom(s) \textbf{après}, soudés → \esp{dámelo}, \esp{levántate} (enclise).
\item Impératif \textbf{négatif} : \esp{no} + pronom(s) + verbe → \esp{no me lo des} (proclise).
\item Ordre fixe : COI avant COD ; \esp{le/les} → \esp{se} devant \esp{lo/la/los/las}.
\item L'enclise exige souvent un accent écrit : \esp{da} → \esp{dámelo}.
\item \esp{vosotros} : chute du -d devant \esp{os} (\esp{levantaos}), sauf \esp{idos}.
\end{itemize}
\end{aretenir}

\coursec{Lexique thématique : santé et environnement}

Dans la section précédente, nous avons appris à placer les pronoms compléments avec l'impératif : \esp{Hierve el agua} devient \esp{Hiérvela}, et \esp{no tires la basura} devient \esp{no la tires}. Mais pour donner de vrais conseils de santé et d'écologie dans un podcast, il faut d'abord posséder les mots. C'est l'objet de cette section : construire, champ par champ, un vocabulaire riche, précis et correctement orthographié, que vous réutiliserez dans la traduction (section 5) et dans votre émission orale (section 6).

\courssub{Observation : un extrait d'émission de radio}

Lisons d'abord ce court texte, début d'un podcast écologique et sanitaire enregistré par des élèves d'un lycée de Yaoundé. Soulignez mentalement tous les mots liés à la santé et à l'environnement.

\begin{textebox}[« Salud y planeta », podcast del colegio]
Bienvenidos a nuestro programa « Salud y planeta ». Hoy hablamos de un tema que nos toca a todos: cómo cuidar nuestro cuerpo y nuestro medio ambiente al mismo tiempo.

Primero, la salud. En nuestra ciudad, el paludismo sigue siendo una enfermedad peligrosa. Por eso, les damos tres consejos: duerman bajo una mosquitera, hiervan el agua antes de beberla y acudan al médico en cuanto tengan fiebre. No olviden que la vacuna protege contra muchas enfermedades: vacunen a sus hijos en el hospital del barrio.

Segundo, el medio ambiente. Nuestro barrio está lleno de basura: botellas de plástico, bolsas, latas. Recuerden: no tiren la basura al río, porque contamina el agua y mata a los peces. Reciclen las botellas, reutilicen las bolsas y reduzcan el plástico. Es importante plantar árboles: los árboles limpian el aire y protegen la tierra.

Cuidar la naturaleza es cuidar nuestra propia salud. Un río limpio, un bosque verde y un aire puro son la mejor medicina. Hasta la próxima emisión, y recuerden: ¡la salud empieza con la higiene!
\end{textebox}

\textbf{Traduction :} « Bienvenue dans notre émission "Santé et planète". Aujourd'hui nous parlons d'un sujet qui nous concerne tous : comment prendre soin de notre corps et de notre environnement en même temps. D'abord, la santé. Dans notre ville, le paludisme reste une maladie dangereuse. C'est pourquoi nous vous donnons trois conseils : dormez sous une moustiquaire, faites bouillir l'eau avant de la boire et allez chez le médecin dès que vous avez de la fièvre. N'oubliez pas que le vaccin protège contre beaucoup de maladies : vaccinez vos enfants à l'hôpital du quartier. Ensuite, l'environnement. Notre quartier est plein d'ordures : bouteilles en plastique, sacs, canettes. Rappelez-vous : ne jetez pas les ordures dans la rivière, car elles polluent l'eau et tuent les poissons. Recyclez les bouteilles, réutilisez les sacs et réduisez le plastique. Il est important de planter des arbres : les arbres purifient l'air et protègent la terre. Prendre soin de la nature, c'est prendre soin de notre propre santé. Une rivière propre, une forêt verte et un air pur sont le meilleur des médicaments. À la prochaine émission, et rappelez-vous : la santé commence par l'hygiène ! »

\textbf{Glossaire :} \esp{la mosquitera} = la moustiquaire ; \esp{hervir} = faire bouillir ; \esp{en cuanto} = dès que ; \esp{la fiebre} = la fièvre ; \esp{la lata} = la canette ; \esp{los peces} = les poissons ; \esp{puro} = pur ; \esp{propio} = propre (à soi).

\textbf{Questions de compréhension :}
\begin{enumerate}
\item ¿Cuáles son los tres consejos contra el paludismo?
\item ¿Por qué no debemos tirar la basura al río?
\item Según el texto, ¿qué relación hay entre la naturaleza y la salud?
\end{enumerate}

\courssub{Champ lexical de la santé}

Le champ lexical de la \cle{santé} s'organise autour de quatre pôles : le corps, la maladie, les soins et les professionnels. Observez les exemples du texte : \esp{cuidar nuestro cuerpo}, \esp{una enfermedad peligrosa}, \esp{acudan al médico}, \esp{en el hospital}.

\begin{center}
\begin{tabularx}{\linewidth}{|X|X|X|X|}
\hline
\rowcolor{popblueL} \textbf{Español} & \textbf{Français} & \textbf{Español} & \textbf{Français} \\
\hline
la salud & la santé & el cuerpo & le corps \\
\hline
la enfermedad & la maladie & el médico / la médica & le médecin \\
\hline
el hospital & l'hôpital & la medicina & le médicament ; la médecine \\
\hline
cuidarse & se soigner, prendre soin de soi & lavarse & se laver \\
\hline
vacunarse & se faire vacciner & descansar & se reposer \\
\hline
la fiebre & la fièvre & estar enfermo & être malade \\
\hline
curar & guérir & la farmacia & la pharmacie \\
\hline
recetar & prescrire (une ordonnance) & la gripe & la grippe \\
\hline
\end{tabularx}
\end{center}

\begin{exemplebox}[Exemples traduits]
\esp{Es importante descansar ocho horas cada noche.} « Il est important de se reposer huit heures chaque nuit. »

\esp{Lávate las manos antes de comer.} « Lave-toi les mains avant de manger. » (Notez l'enclise du pronom réfléchi : \esp{lávate}, avec accent écrit.)

\esp{El médico me recetó una medicina contra la fiebre.} « Le médecin m'a prescrit un médicament contre la fièvre. »
\end{exemplebox}

\begin{attention}[Faux amis : medicina et constipado]
\esp{La medicina} signifie « le médicament » (ce qu'on avale) mais aussi « la médecine » (la science). Le contexte décide : \esp{Estudio medicina} = « J'étudie la médecine » ; \esp{Toma la medicina} = « Prends ton médicament ».

Piège classique : \esp{estar constipado} signifie « être enrhumé » et non « être constipé » ! La constipation se dit \esp{el estreñimiento}. De même, \esp{una intoxicación} désigne une intoxication alimentaire, pas l'ivresse.
\end{attention}

\courssub{Champ lexical de l'hygiène et de la prévention}

Ce champ est essentiel dans le contexte camerounais et africain : la lutte contre le paludisme et la dengue, l'accès à l'eau potable, la vaccination.

\begin{center}
\begin{tabularx}{\linewidth}{|X|X|X|X|}
\hline
\rowcolor{popgreenL} \textbf{Español} & \textbf{Français} & \textbf{Español} & \textbf{Français} \\
\hline
el agua potable & l'eau potable & hervir (ie) & faire bouillir \\
\hline
la mosquitera & la moustiquaire & el mosquito & le moustique \\
\hline
el dengue & la dengue & el paludismo & le paludisme \\
\hline
la vacuna & le vaccin & la higiene & l'hygiène \\
\hline
el jabón & le savon & prevenir (ie) & prévenir \\
\hline
la picadura & la piqûre & limpio / sucio & propre / sale \\
\hline
\end{tabularx}
\end{center}

\begin{propriete}[Particularités à retenir]
\begin{itemize}
\item \esp{El agua} est féminin mais prend l'article masculin \esp{el} devant \emph{a} accentuée : \esp{el agua potable}, mais \esp{las aguas}. Ne dites jamais \esp{*la agua}. L'adjectif reste au féminin : \esp{el agua está fría}.
\item \esp{Hervir} est un verbe à changement radical (e $\to$ ie) : \esp{hierve el agua} (impératif \esp{tú}), \esp{hiervan el agua} (impératif \esp{ustedes}). Même changement pour \esp{prevenir} : \esp{previene}, \esp{prevengan}.
\item \esp{El paludismo} est le mot courant en Afrique hispanophone (Guinée équatoriale) ; \esp{la malaria} existe aussi.
\item \esp{El dengue} est masculin : \esp{el dengue}, comme \esp{el paludismo}.
\end{itemize}
\end{propriete}

\begin{exemplebox}[Exemples traduits]
\esp{Hierve el agua antes de beberla.} « Fais bouillir l'eau avant de la boire. » (Notez l'enclise : \esp{beberla}.)

\esp{No olvides dormir bajo la mosquitera.} « N'oublie pas de dormir sous la moustiquaire. »

\esp{La higiene es la primera medicina.} « L'hygiène est la première des médecines. »
\end{exemplebox}

\courssub{Champ lexical de l'environnement et de la nature}

Ces deux champs se complètent : l'\cle{environnement} (les problèmes et les solutions) et la \cle{nature} (les éléments à protéger).

\begin{center}
\begin{tabularx}{\linewidth}{|X|X|X|X|}
\hline
\rowcolor{poporangeL} \textbf{Español} & \textbf{Français} & \textbf{Español} & \textbf{Français} \\
\hline
el medio ambiente & l'environnement & la basura & les ordures \\
\hline
la contaminación & la pollution & reciclar & recycler \\
\hline
reutilizar & réutiliser & reducir & réduire \\
\hline
el vertedero & la décharge & el plástico & le plastique \\
\hline
la botella & la bouteille & la bolsa & le sac \\
\hline
tirar & jeter & contaminar & polluer \\
\hline
\end{tabularx}
\end{center}

\begin{center}
\begin{tabularx}{\linewidth}{|X|X|X|X|}
\hline
\rowcolor{popgreenL} \textbf{Español} & \textbf{Français} & \textbf{Español} & \textbf{Français} \\
\hline
el río & la rivière, le fleuve & el bosque & la forêt \\
\hline
la tierra & la terre & el aire & l'air \\
\hline
los árboles & les arbres & plantar & planter \\
\hline
proteger & protéger & la naturaleza & la nature \\
\hline
el mar & la mer & el pez / los peces & le poisson / les poissons \\
\hline
\end{tabularx}
\end{center}

\begin{exemplebox}[Exemples traduits]
\esp{No tires la basura al río.} « Ne jette pas les ordures dans la rivière. »

\esp{Recicla las botellas de plástico y reutiliza las bolsas.} « Recycle les bouteilles en plastique et réutilise les sacs. »

\esp{Planten árboles para proteger la tierra y limpiar el aire.} « Plantez des arbres pour protéger la terre et purifier l'air. » (Notez \esp{para} + infinitif pour exprimer le but.)
\end{exemplebox}

\begin{savaistu}[« Medio ambiente » en deux mots]
En espagnol, « environnement » se dit \esp{el medio ambiente}, en \textbf{deux mots}. Sous l'influence de l'anglais \emph{environment} et de certaines graphies journalistiques, on voit parfois le calque fautif \esp{*medioambiente} en un seul mot : c'est une erreur à éviter absolument au Bac. Retenez aussi que \esp{la contaminación} est le mot courant pour « pollution » ; \esp{la polución} existe mais est moins fréquent. En Guinée équatoriale, seul pays hispanophone d'Afrique subsaharienne, la protection des forêts tropicales est un sujet d'actualité majeur : un excellent thème pour votre podcast !
\end{savaistu}

\courssub{Carte mentale : organiser le lexique}

Pour mémoriser durablement, rien ne vaut une carte mentale. Voici la nôtre, à recopier et à enrichir dans votre cahier.

\begin{popfigure}
\begin{center}
\begin{tikzpicture}[mindmap, grow cyclic, every node/.style={concept, align=center},
root concept/.append style={concept color=popblue, font=\bfseries},
level 1/.append style={level distance=3.4cm, sibling angle=90},
level 2/.append style={level distance=2.2cm, sibling angle=45}]
\node[root concept]{Salud y medio ambiente}
  child[concept color=popgreen]{ node {salud}
      child { node[concept color=popgreenL] {cuerpo} }
      child { node[concept color=popgreenL] {médico} }
      child { node[concept color=popgreenL] {descansar} } }
  child[concept color=poporange]{ node {higiene}
      child { node[concept color=poporangeL] {agua potable} }
      child { node[concept color=poporangeL] {mosquitera} }
      child { node[concept color=poporangeL] {vacuna} } }
  child[concept color=poppurple]{ node {medio ambiente}
      child { node[concept color=poppurpleL] {basura} }
      child { node[concept color=poppurpleL] {reciclar} }
      child { node[concept color=poppurpleL] {plástico} } }
  child[concept color=popgold]{ node {naturaleza}
      child { node[concept color=popgoldL] {río} }
      child { node[concept color=popgoldL] {bosque} }
      child { node[concept color=popgoldL] {plantar} } };
\end{tikzpicture}
\end{center}
\legende{Carte mentale du lexique : quatre branches, trois mots-clés chacune. Complétez-la avec les tableaux bilingues.}
\end{popfigure}

\courssub{Les expressions de consigne pour une émission}

Un podcast de conseils repose sur des formules fixes qui encadrent l'impératif et le subjonctif. Mémorisez-les : elles donneront à votre émission un ton professionnel.

\begin{center}
\begin{tabularx}{\linewidth}{|X|X|X|}
\hline
\rowcolor{popblueL} \textbf{Español} & \textbf{Français} & \textbf{Remarque} \\
\hline
Les aconsejamos que + subj. & Nous vous conseillons de + inf. & \esp{Les aconsejamos que duerman más.} \\
\hline
Recuerden + inf. / que... & Rappelez-vous de... & \esp{Recuerden lavarse las manos.} \\
\hline
No olviden + inf. & N'oubliez pas de... & \esp{No olviden la mosquitera.} \\
\hline
Es importante + inf. & Il est important de... & \esp{Es importante beber agua potable.} \\
\hline
Es necesario + inf. & Il est nécessaire de... & \esp{Es necesario reciclar.} \\
\hline
Hay que + inf. & Il faut... & \esp{Hay que proteger el bosque.} \\
\hline
Les damos tres consejos & Nous vous donnons trois conseils & formule d'annonce \\
\hline
\end{tabularx}
\end{center}

\begin{attention}[Subjonctif ou infinitif ?]
Après \esp{es importante}, \esp{es necesario}, \esp{hay que} (tournures impersonnelles sans sujet précis), on utilise l'\textbf{infinitif} : \esp{Es importante descansar}. Mais après \esp{les aconsejamos que}, il y a un sujet précis (vous, les auditeurs) : l'espagnol exige alors le \textbf{subjonctif} : \esp{Les aconsejamos que se vacunen} (« Nous vous conseillons de vous faire vacciner »). Ne dites jamais \esp{*les aconsejamos vacunarse}. Rappel : l'impératif négatif et l'impératif de politesse (\esp{usted}, \esp{ustedes}) empruntent justement leurs formes au subjonctif présent : \esp{no tires}, \esp{reciclen}, \esp{duerman}.
\end{attention}

\begin{textebox}[Mini-lexique de survie du podcasteur]
Bienvenidos a nuestro programa... = Bienvenue dans notre émission...

Hoy hablamos de... = Aujourd'hui nous parlons de...

Les damos tres consejos. = Nous vous donnons trois conseils.

Primero... Segundo... Finalmente... = D'abord... Ensuite... Enfin...

Recuerden: la salud es lo primero. = Rappelez-vous : la santé d'abord.

Gracias por escucharnos. = Merci de nous avoir écoutés.

Hasta la próxima emisión. = À la prochaine émission.
\end{textebox}

\courssub{Entraînement guidé}

\begin{exoresolu}[Traduire en espagnol]
Énoncé : Traduisez en espagnol : « Nous vous conseillons de faire bouillir l'eau avant de la boire. N'oubliez pas de recycler les bouteilles en plastique et ne jetez pas les ordures dans la rivière. »
\tcblower
Solution détaillée :
\begin{itemize}
\item « Nous vous conseillons de... » + sujet « vous » : \esp{Les aconsejamos que} + subjonctif : \esp{Les aconsejamos que hiervan el agua} (subjonctif de \esp{hervir}, changement radical e $\to$ ie).
\item « avant de la boire » : infinitif + enclise : \esp{antes de beberla}.
\item « N'oubliez pas de recycler » : \esp{No olviden reciclar las botellas de plástico}.
\item « ne jetez pas » : impératif négatif, 2\textsuperscript{e} pers. du pluriel : \esp{no tiren la basura al río}.
\end{itemize}
Traduction complète : \esp{Les aconsejamos que hiervan el agua antes de beberla. No olviden reciclar las botellas de plástico y no tiren la basura al río.}
\end{exoresolu}

\begin{exercice}[À vous de jouer]
\begin{enumerate}
\item Traduisez : « Il est important de dormir sous la moustiquaire et de se laver les mains avec du savon. »
\item Traduisez : « Recyclez les bouteilles, réutilisez les sacs et réduisez le plastique. »
\item Complétez avec le mot du champ lexical qui convient : \esp{El \dots\ es una enfermedad que transmite el mosquito.} / \esp{Hay que hervir el \dots\ antes de beberla.} / \esp{Los \dots\ limpian el aire y protegen la tierra.}
\end{enumerate}
\end{exercice}

\begin{corrige}[Exercice « À vous de jouer »]
1. \esp{Es importante dormir bajo la mosquitera y lavarse las manos con jabón.} (Tournure impersonnelle : infinitif.)

2. \esp{Reciclen las botellas, reutilicen las bolsas y reduzcan el plástico.} (Impératif \esp{ustedes} ; attention à \esp{reduzcan}, changement c $\to$ zc devant a, comme au subjonctif présent.)

3. \esp{El paludismo es una enfermedad que transmite el mosquito.} / \esp{Hay que hervir el agua antes de beberla.} / \esp{Los árboles limpian el aire y protegen la tierra.}
\end{corrige}

\begin{aretenir}[L'essentiel de la section]
\begin{itemize}
\item Quatre champs lexicaux : \esp{salud}, \esp{higiene}, \esp{medio ambiente}, \esp{naturaleza} — à mémoriser avec la carte mentale.
\item Faux amis : \esp{constipado} = enrhumé ; \esp{medicina} = médicament ou médecine selon le contexte.
\item Orthographe : \esp{medio ambiente} en deux mots ; \esp{el agua} avec article masculin mais adjectif féminin.
\item Formules d'émission : \esp{les aconsejamos que} + subjonctif ; \esp{recuerden}, \esp{no olviden}, \esp{es importante} + infinitif.
\end{itemize}
\end{aretenir}

Vous disposez maintenant de tout le vocabulaire nécessaire. Dans la section suivante, nous le mettrons à l'épreuve dans des exercices de thème et de version, avant de passer au micro pour votre propre podcast écologique.

\coursec{Entraînement à la traduction : thème et version}

Dans la section précédente, nous avons constitué notre boîte à outils lexicale : les mots de la santé (\esp{lavarse las manos, cuidarse, hacer ejercicio}) et ceux de l'environnement (\esp{tirar la basura, reciclar, proteger el planeta}). Il est temps maintenant de mettre ce vocabulaire en mouvement dans l'exercice roi du baccalauréat : la \cle{traduction}. Traduire une phrase impérative avec pronoms compléments exige une méthode rigoureuse : une erreur d'accent, de place ou d'ordre des pronoms suffit à fausser toute la phrase. Cette section vous donne une méthode en cinq étapes, puis un entraînement guidé en version (espagnol $\rightarrow$ français) et en thème (français $\rightarrow$ espagnol), avec correction collective justifiée.

\courssub{La méthode en cinq étapes}

\begin{methode}[Traduire une phrase impérative avec pronoms]
\begin{enumerate}
\item \textbf{Repérer le mode et la polarité} : la phrase est-elle à l'impératif affirmatif ou négatif ? C'est ce qui décide de la \emph{place} des pronoms.
\item \textbf{Identifier la personne} : \esp{tú}, \esp{vosotros}, \esp{usted}, \esp{ustedes} ? La forme verbale change selon le destinataire de la consigne.
\item \textbf{Identifier les pronoms} : y a-t-il un ou deux pronoms compléments ? Le verbe est-il pronominal (\esp{lavarse}, \esp{sentarse}) ?
\item \textbf{Placer et ordonner les pronoms} : à l'affirmatif, ils se soudent \emph{après} le verbe (enclise) dans l'ordre COI + COD (\esp{dámelo} : \esp{me} = à moi, \esp{lo} = le) ; au négatif, ils se placent \emph{avant} le verbe, séparés (proclise) : \esp{no me lo des}.
\item \textbf{Vérifier l'accent écrit} : l'enclise déplace l'accent tonique ; il faut souvent ajouter un accent graphique (\esp{diga} $\rightarrow$ \esp{díganselo}). Enfin, contrôler le sens global de la phrase traduite.
\end{enumerate}
\end{methode}

\begin{popfigure}
\begin{tikzpicture}[
  etape/.style={rectangle, rounded corners=3pt, draw=popblue, fill=popblueL, text width=3.1cm, align=center, font=\small, minimum height=1.3cm},
  fleche/.style={-{Stealth[length=2.5mm]}, thick, poporange}
]
\node[etape] (e1) at (0,0) {\textbf{1. Mode} affirmatif ou négatif ?};
\node[etape, right=0.5cm of e1] (e2) {\textbf{2. Personne} tú / vosotros / usted / ustedes};
\node[etape, right=0.5cm of e2] (e3) {\textbf{3. Pronoms} un, deux, verbe pronominal ?};
\node[etape, below=0.6cm of e3] (e4) {\textbf{4. Place} enclise (aff.) / proclise (nég.)};
\node[etape, left=0.5cm of e4] (e5) {\textbf{5. Accent} et sens global};
\draw[fleche] (e1) -- (e2);
\draw[fleche] (e2) -- (e3);
\draw[fleche] (e3) -- (e4);
\draw[fleche] (e4) -- (e5);
\end{tikzpicture}
\legende{Les cinq étapes de la méthode de traduction de l'impératif avec pronoms.}
\end{popfigure}

\courssub{Version guidée : de l'espagnol vers le français}

\begin{methode}[Réussir la version]
Avant de traduire mot à mot, posez-vous deux questions : \textbf{(1)} l'impératif est-il affirmatif ou négatif ? \textbf{(2)} quels pronoms sont soudés au verbe ou placés devant lui ? Détachez mentalement chaque pronom, identifiez-le (personne, COD ou COI), puis seulement traduisez. En français, l'ordre des pronoms à l'impératif affirmatif est différent de l'espagnol : méfiez-vous !
\end{methode}

Voici cinq phrases de difficulté croissante, tirées du contexte d'une émission de santé et d'écologie.

\begin{exemplebox}[Cinq phrases à traduire (version)]
\begin{enumerate}
\item \esp{Come más fruta y bebe agua.} (sans pronom)
\item \esp{Lávate las manos antes de comer.} (verbe pronominal)
\item \esp{No la tires en la calle.} (un pronom, négatif — \esp{la} = la basura)
\item \esp{Dámela, por favor.} (deux pronoms, affirmatif — \esp{la} = la botella)
\item \esp{No se lo digan a los niños.} (deux pronoms, négatif, \esp{ustedes})
\end{enumerate}
\end{exemplebox}

\begin{exoresolu}[Correction guidée de la version]
Traduire : \esp{No se lo digan a los niños.}
\tcblower
\textbf{Étape 1} : la phrase est négative (\esp{no}) $\rightarrow$ les pronoms sont \emph{avant} le verbe.\par
\textbf{Étape 2} : \esp{digan} = subjonctif présent de \esp{decir}, 3\textsuperscript{e} pers. plur. $\rightarrow$ impératif \esp{ustedes} = « vous ».\par
\textbf{Étape 3} : deux pronoms : \esp{se} (= \esp{les}, COI, ici « aux enfants ») et \esp{lo} (COD, « le » = la nouvelle, l'information).\par
\textbf{Étape 4} : en français, à l'impératif négatif, l'ordre est : ne + le/la/les + lui/leur + verbe.\par
\textbf{Traduction} : « Ne le leur dites pas aux enfants. » (ou plus naturel : « Ne le dites pas aux enfants. »)\par
\textbf{Vérification} : aucun accent à contrôler ici (proclise), le sens global est une consigne d'interdiction.
\end{exoresolu}

Voici maintenant la correction complète des cinq phrases, avec la justification de chaque choix.

\begin{center}
\begin{tabularx}{\linewidth}{|X|X|X|}
\hline
\rowcolor{popblueL} \textbf{Espagnol} & \textbf{Français} & \textbf{Justification} \\
\hline
\esp{Come más fruta y bebe agua.} & Mange plus de fruits et bois de l'eau. & Impératif affirmatif de \esp{comer} et \esp{beber} (\esp{tú}), sans pronom : traduction directe. \\
\hline
\esp{Lávate las manos antes de comer.} & Lave-toi les mains avant de manger. & Verbe pronominal \esp{lavarse} ; à l'affirmatif, \esp{te} se soude au verbe ; accent ajouté : \esp{lava} $\rightarrow$ \esp{lávate}. \\
\hline
\esp{No la tires en la calle.} & Ne la jette pas dans la rue. & Négatif $\rightarrow$ proclise : \esp{la} (COD, la basura) avant \esp{tires} (subjonctif). \\
\hline
\esp{Dámela, por favor.} & Donne-la-moi, s'il te plaît. & Affirmatif $\rightarrow$ enclise, ordre COI + COD : \esp{me} (à moi) + \esp{la} (la botella) ; accent sur \esp{dá-}. \\
\hline
\esp{No se lo digan a los niños.} & Ne le leur dites pas aux enfants. & Négatif, \esp{ustedes} ; \esp{se} = COI (« leur ») devant \esp{lo} = COD (« le »). \\
\hline
\end{tabularx}
\end{center}

\begin{attention}[L'ordre des pronoms : le piège majeur de la version]
En espagnol, l'ordre est toujours \textbf{COI puis COD} : \esp{dámelo} = \esp{me} (à moi) + \esp{lo} (le). En français, à l'impératif affirmatif, c'est l'inverse : le COD passe d'abord, puis le COI : « donne-\textbf{le}-\textbf{moi} ». Donc \esp{dámelo} se traduit « donne-le-moi » et jamais « donne-moi-le ». De même, \esp{díganselo} = « dites-le-leur » (\esp{se} = leur, \esp{lo} = le). En thème, appliquez le chemin inverse : « donne-le-moi » = \esp{dámelo}.
\end{attention}

\courssub{Thème guidé : du français vers l'espagnol}

Appliquons maintenant la méthode en cinq étapes à la traduction de phrases françaises. Les quatre phrases imposées de notre séquence santé/environnement serviront de fil conducteur, complétées par une cinquième phrase à deux pronoms.

\begin{exemplebox}[Cinq phrases à traduire (thème)]
\begin{enumerate}
\item Lave-toi les mains. (impératif affirmatif, pronominal)
\item Ne jette pas les ordures par terre. (impératif négatif, \esp{tú})
\item Donne-la-moi. (la bouteille = \esp{la botella})
\item Ne la leur donnez pas. (la bouteille, aux enfants)
\item Asseyez-vous, s'il vous plaît. (formule de politesse, \esp{usted})
\end{enumerate}
\end{exemplebox}

\textbf{Correction collective détaillée.}

\textbf{Phrase 1 — « Lave-toi les mains. »}\par
Étape 1 : affirmatif $\rightarrow$ enclise. Étape 2 : « toi » $\rightarrow$ \esp{tú}. Étape 3 : verbe pronominal \esp{lavarse}, pronom \esp{te}. Étape 4 : \esp{lava} + \esp{te} soudés. Étape 5 : l'accent tonique tombe sur la première syllabe, il faut un accent écrit : \esp{\textbf{Lávate} las manos.} Sans accent, \esp{lavate} serait fautif.

\textbf{Phrase 2 — « Ne jette pas les ordures par terre. »}\par
Étape 1 : négatif $\rightarrow$ proclise et subjonctif. Étape 2 : \esp{tú}. Étape 3 : pas de pronom complément ici (« les ordures » est un groupe nominal : \esp{la basura} ou \esp{los residuos}). Étape 4 : \esp{no} + subjonctif présent de \esp{tirar} : \esp{tires}. Traduction : \esp{No tires la basura al suelo.} Remarquez : le français « ne... pas » se rend par le seul \esp{no}.

\textbf{Phrase 3 — « Donne-la-moi. » (la bouteille)}\par
Étape 1 : affirmatif $\rightarrow$ enclise. Étape 2 : \esp{tú} $\rightarrow$ \esp{da} (impératif irrégulier de \esp{dar}). Étape 3 : deux pronoms : « la » = COD (\esp{la}, car \esp{la botella} est féminin) et « moi » = COI (\esp{me}). Étape 4 : en espagnol, le COI passe d'abord : \esp{da} + \esp{me} + \esp{la}. Étape 5 : accent obligatoire pour conserver l'accent tonique : \esp{\textbf{Dámela}.} L'ordre français « la-moi » devient \esp{me-la} : c'est l'inversion signalée dans l'encadré d'attention.

\textbf{Phrase 4 — « Ne la leur donnez pas. »}\par
Étape 1 : négatif $\rightarrow$ proclise. Étape 2 : « vous » $\rightarrow$ \esp{ustedes} (registre soigné) ou \esp{vosotros} ; choisissons \esp{ustedes} : subjonctif \esp{den}. Étape 3 : « la » = COD (\esp{la}) ; « leur » = COI, qui devient \esp{se} devant \esp{la} (on ne dit jamais \esp{les la}). Étape 4 : \esp{no} + \esp{se} + \esp{la} + verbe. Traduction : \esp{\textbf{No se la den}.} (ou \esp{No se la deis.} avec \esp{vosotros}). Étape 5 : pas d'accent, car les pronoms sont détachés.

\textbf{Phrase 5 — « Asseyez-vous, s'il vous plaît. »}\par
Étape 1 : affirmatif. Étape 2 : « vous » de politesse $\rightarrow$ \esp{usted}. Étape 3 : verbe pronominal \esp{sentarse} ; avec \esp{usted}, le pronom réflexif est \esp{se}. Étape 4 : \esp{siente} (subjonctif présent de \esp{sentarse}, 3\textsuperscript{e} pers. sing.) + \esp{se}. Étape 5 : accent : \esp{\textbf{Siéntese}, por favor.}

\begin{propriete}[Rappel : le pronom réflexif avec usted / ustedes]
Avec \esp{usted} et \esp{ustedes}, le pronom réflexif est toujours \esp{se}, jamais \esp{te} ni \esp{os} :
\begin{itemize}
\item \esp{Siéntese, por favor.} « Asseyez-vous, s'il vous plaît. »
\item \esp{No se sienten aquí.} « Ne vous asseyez pas ici. » (attention : avec \esp{ustedes}, le verbe est au pluriel : \esp{no se sienten})
\item \esp{Cúidense mucho.} « Prenez bien soin de vous. »
\end{itemize}
L'impératif de \esp{usted/ustedes} emprunte ses formes au subjonctif présent : \esp{sentarse} $\rightarrow$ \esp{siéntese / siéntense} ; \esp{cuidarse} $\rightarrow$ \esp{cuídese / cuídense}.
\end{propriete}

\begin{exoresolu}[Thème complet : phrase à deux pronoms]
Traduire en espagnol : « Ne la leur donnez pas » (la bouteille d'eau, aux enfants — registre \esp{ustedes}).
\tcblower
\textbf{Étape 1} : phrase négative $\rightarrow$ proclise : les pronoms se placent avant le verbe, détachés.\par
\textbf{Étape 2} : « vous » $\rightarrow$ \esp{ustedes} ; l'impératif négatif = \esp{no} + subjonctif présent de \esp{dar} : \esp{den}.\par
\textbf{Étape 3} : deux pronoms : COD « la » = \esp{la} (la botella) ; COI « leur » = théoriquement \esp{les}, mais devant un COD \esp{lo/la/los/las}, \esp{le/les} devient \esp{se}.\par
\textbf{Étape 4} : ordre espagnol : COI puis COD $\rightarrow$ \esp{no se la den}.\par
\textbf{Étape 5} : pas d'accent écrit (proclise) ; sens global : interdiction adressée à plusieurs personnes.\par
\textbf{Traduction} : \esp{No se la den (a los niños).}
\end{exoresolu}

\begin{center}
\begin{tabularx}{\linewidth}{|X|X|X|}
\hline
\rowcolor{popblueL} \textbf{Français} & \textbf{Espagnol} & \textbf{Point de méthode} \\
\hline
Lave-toi les mains. & \esp{Lávate las manos.} & Enclise + accent (\esp{lava} $\rightarrow$ \esp{lávate}). \\
\hline
Ne jette pas les ordures par terre. & \esp{No tires la basura al suelo.} & Négatif = \esp{no} + subjonctif. \\
\hline
Donne-la-moi. & \esp{Dámela.} & Inversion de l'ordre : COI + COD ; accent. \\
\hline
Ne la leur donnez pas. & \esp{No se la den.} & \esp{les} $\rightarrow$ \esp{se} devant \esp{la} ; proclise. \\
\hline
Asseyez-vous, s'il vous plaît. & \esp{Siéntese, por favor.} & \esp{usted} + réflexif \esp{se} + accent. \\
\hline
\end{tabularx}
\end{center}

\courssub{Grille d'auto-évaluation}

Après chaque traduction, interrogez votre production à l'aide de cette grille. Un « oui » à chaque ligne garantit une phrase juste.

\begin{center}
\begin{tabularx}{\linewidth}{|X|X|}
\hline
\rowcolor{popblueL} \textbf{Critère} & \textbf{Question à se poser} \\
\hline
Forme de l'impératif & Ai-je utilisé la bonne personne (\esp{tú}, \esp{vosotros}, \esp{usted}, \esp{ustedes}) et, au négatif, le subjonctif ? \\
\hline
Place du pronom & Affirmatif $\rightarrow$ pronom soudé après le verbe ? Négatif $\rightarrow$ pronom détaché avant le verbe ? \\
\hline
Ordre des pronoms & En espagnol : COI avant COD (\esp{me lo}, \esp{se la}) ? Ai-je transformé \esp{le/les} en \esp{se} ? \\
\hline
Accent écrit & L'enclise a-t-elle exigé un accent (\esp{dámelo}, \esp{siéntese}, \esp{lávate}) ? \\
\hline
Sens global & Ma phrase traduite signifie-t-elle exactement la même chose que l'originale ? \\
\hline
\end{tabularx}
\end{center}

\courssub{Complément utile Bac : l'impératif de courtoisie avec usted}

\begin{aretenir}[Les consignes officielles : « Tenga en cuenta que... »]
Dans les documents officiels, les notices de santé publique et les consignes d'examen, l'espagnol utilise l'impératif de courtoisie avec \esp{usted}, formé sur le subjonctif présent. Expressions fréquentes au Bac :
\begin{itemize}
\item \esp{Tenga en cuenta que...} « Veuillez noter que... / Tenez compte du fait que... »
\item \esp{No dude en consultarnos.} « N'hésitez pas à nous consulter. »
\item \esp{Consulte a su médico.} « Consultez votre médecin. »
\item \esp{Mantenga este producto fuera del alcance de los niños.} « Maintenez ce produit hors de portée des enfants. »
\item \esp{Evite tirar residuos al suelo.} « Évitez de jeter des déchets par terre. »
\end{itemize}
Ces formules sont idéales pour conclure votre podcast écologique ou votre émission de santé (section suivante) : elles donnent à votre production orale un ton professionnel et respectueux.
\end{aretenir}

\begin{exercice}[À vous de jouer]
Traduisez en espagnol, en appliquant les cinq étapes et en vous auto-évaluant avec la grille :
\begin{enumerate}
\item Ne te lève pas tout de suite. (conseil du médecin)
\item Recyclez les bouteilles en plastique. (\esp{ustedes})
\item Apporte-moi-la. (la trousse de secours = \esp{el botiquín} — attention au genre !)
\item Ne nous le répétez pas.
\item Prenez soin de vous et protégez la planète. (\esp{ustedes})
\end{enumerate}
\end{exercice}

\begin{corrige}[Exercice « À vous de jouer »]
\begin{enumerate}
\item \esp{No te levantes enseguida.} — négatif, \esp{tú}, pronominal : proclise, subjonctif \esp{levantes}.
\item \esp{Reciclen las botellas de plástico.} — affirmatif \esp{ustedes}, sans pronom : forme en \esp{-en}.
\item \esp{Tráemelo.} — \esp{el botiquín} est masculin $\rightarrow$ \esp{lo} ; ordre COI + COD ; accent sur \esp{trá-} (impératif irrégulier de \esp{traer}).
\item \esp{No nos lo repitan.} — négatif, \esp{ustedes} : proclise, \esp{nos} (COI) avant \esp{lo} (COD), subjonctif \esp{repitan}.
\item \esp{Cúidense y protejan el planeta.} — pronominal avec \esp{ustedes} : \esp{se} soudé, accent sur \esp{cuí-} ; coordination de deux impératifs.
\end{enumerate}
\end{corrige}

Vous disposez désormais de tous les outils grammaticaux et méthodologiques pour traduire sans faute les consignes impératives. La prochaine étape consiste à exploiter cette compétence dans une production orale authentique : l'enregistrement d'une émission de santé ou d'un podcast écologique, où chaque conseil que vous donnerez à vos auditeurs mobilisera naturellement l'impératif et ses pronoms.

\coursec{Production : émission de santé / podcast écologique}

Après avoir consolidé la traduction de l'impératif et le placement des pronoms compléments dans la section précédente, nous passons maintenant à la mise en œuvre communicative. L'objectif est de produire un message oral authentique, structuré et grammaticalement correct, conforme aux attendus du Bac (épreuve d'expression orale) et aux formats médiatiques actuels. Le podcast écologique ou sanitaire est un genre court, dynamique, qui mobilise l'impératif pour donner des consignes, des conseils, des ordres ou des interdictions.

\courssub{Analyse d'un modèle de script : structure et contraintes linguistiques}

Avant de rédiger, observons un modèle complet respectant toutes les contraintes grammaticales imposées : au moins 4 impératifs affirmatifs, 2 impératifs négatifs, 3 pronoms enclitiques, 2 pronoms proclitiques, avec cohérence de la personne (ici \cle{ustedes} pour s'adresser à un public large).

\begin{textebox}[Modèle de script : « Barrio limpio », épisode 1]
Hola, amigos. Bienvenidos a « Barrio limpio \}, vuestro podcast ecológico semanal. Soy vuestro presentador, Carlos, y hoy les traigo tres consejos esenciales para un barrio más sano.

Primero, \textbf{reciclen} el plástico y el papel, \textbf{no los tiren} a la calle; \textbf{sepárenlos} bien en los contenedores amarillos y azules. Segundo, \textbf{hiervan} el agua de grifo antes de \textbf{beberla} y \textbf{no la desperdicien}; el agua es vida. Tercero, \textbf{planten} un árbol en su patio o en la escuela y \textbf{cuidenlo} con cariño; \textbf{no lo abandonen} tras la primera semana.

\emph{¡Protéjanse} ustedes mismos y \textbf{protejan} su entorno! Hasta la próxima, amigos. ¡Cuídense mucho!
\end{textebox}

\noindent\textbf{Traduction et glossaire :} « Bonjour les amis. Bienvenue sur "Quartier propre", votre podcast écologique hebdomadaire. Je suis votre animateur, Carlos, et aujourd'hui je vous apporte trois conseils essentiels pour un quartier plus sain. Premièrement, recyclez le plastique et le papier, ne les jetez pas dans la rue ; séparez-les bien dans les conteneurs jaunes et bleus. Deuxièmement, faites bouillir l'eau du robinet avant de la boire et ne la gaspillez pas ; l'eau c'est la vie. Troisièmement, plantez un arbre dans votre cour ou à l'école et prenez-en soin avec affection ; ne l'abandonnez pas après la première semaine. Protégez-vous vous-mêmes et protégez votre environnement ! À la prochaine, amis. Prenez bien soin de vous ! »

\cle{Vocabulaire clé} : \esp{el contenedor} (conteneur/poubelle de tri), \esp{el grifo} (robinet), \esp{desperdiciar} (gaspiller), \esp{el patio} (cour), \esp{abandonar} (abandonner), \esp{el entorno} (environnement).

\begin{aretenir}[Analyse grammaticale du modèle]
\textbf{Structure type d'un podcast conseil :}
\begin{enumerate}
    \item \textbf{Introducción} : Salutation + nom de l'émission + identité du locuteur + annonce du sujet (\esp{Hoy les traigo tres consejos...}).
    \item \textbf{Cuerpo (Consejos)} : Énumération (Primero, Segundo, Tercero) + \cle{impératifs} + \cle{justification courte} (souvent introduite par \esp{;} ou \esp{porque}).
    \item \textbf{Conclusión} : Appel à l'action global + rappel identitaire.
    \item \textbf{Despedida} : Formule de politesse + signature.
\end{enumerate}

\textbf{Vérification des contraintes (modèle) :}
\begin{itemize}
    \item \textbf{Impératifs affirmatifs (7)} : \esp{reciclen, hiervan, planten, cuidenlo, protéjanse, protejan, cuídense}. (Note : \esp{cuidenlo, protéjanse, cuídense} portent pronoms enclitiques ; \esp{beberla} est un infinitif avec enclise, non un impératif).
    \item \textbf{Impératifs négatifs (3)} : \esp{no los tiren, no la desperdicien, no lo abandonen}. (Pronoms proclitiques : \esp{los, la, lo}).
    \item \textbf{Pronoms enclitiques (4 sur formes impératives)} : \esp{sepárenlos, cuidenlo, protéjanse, cuídense}.
    \item \textbf{Pronoms proclitiques (3)} : \esp{no los tiren, no la desperdicien, no lo abandonen}.
    \item \textbf{Cohérence personne} : Tout est à \cle{ustedes} (public large, formel/standard média).
\end{itemize}
\end{aretenir}

\courssub{Méthode : Construire un conseil efficace (Impératif + Complément + Justification)}

Un bon conseil en podcast ne se limite pas à l'ordre ; il persuade. La structure \cle{Impératif + Complément d'objet (prononimal ou nominal) + Justification (causale ou finale)} est la plus naturelle.

\begin{methode}[Rédiger une séquence « Conseil » pour le podcast]
\textbf{Étape 1 : Choisir la personne et le verbe d'action.}
Décidez si vous tutoyez (\cle{tú} / \cle{vosotros} — Espagne) ou vous vouvoyez (\cle{usted} / \cle{ustedes}). Pour un podcast grand public : \textbf{ustedes}. Conjuguez le verbe à l'impératif affirmatif ou négatif.
\esp{Reciclen} (ustedes, affirmatif) / \esp{No tiren} (ustedes, négatif).

\textbf{Étape 2 : Ajouter le(s) pronom(s) complément(s) selon la polarité.}
\begin{itemize}
    \item \textbf{Affirmatif} $\rightarrow$ \cle{Enclise} (collé après, accentuation si la règle générale l'exige) : \esp{Reciclen + los $\rightarrow$ Reciclénlos}. \esp{Separen + los $\rightarrow$ Sepárenlos}.
    \item \textbf{Négatif} $\rightarrow$ \cle{Proclise} (avant le verbe, après \esp{no}) : \esp{No + los + tiren $\rightarrow$ No los tiren}.
\end{itemize}
\textbf{Rappel ordre pronoms :} \cle{Se} (le/les) $>$ \cle{Me/Te/Nos/Os} $>$ \cle{Le/Les/Se} $>$ \cle{Lo/La/Los/Las}. Ex: \esp{Dáselo} (Da + se + lo), \esp{No se lo den}.

\textbf{Étape 3 : Formuler la justification courte (connecteur + indicatif/ infinitif).}
Utilisez \esp{porque} + indicatif, \esp{para} + infinitif, \esp{ya que} + indicatif, ou simple \esp{;} (point-virgule) pour lier cause/effet.
\esp{Reciclénlos; protegen el río.} / \esp{No los tiren, porque contaminan el suelo.}

\textbf{Étape 4 : Varier les verbes et les structures.}
Alternez verbes réguliers/irréguliers, transitifs directs/indirects, pronominaux.
\esp{¡Cuídense!} (pronominal), \esp{¡Planténlo!} (transitif direct + enclise), \esp{¡No se lo pierdan!} (double pronom proclitique).
\end{methode}

\begin{exemplebox}[Exemples variés de conseils structurés]
\begin{center}
\begin{tabularx}{\linewidth}{|X|X|X|}
\hline
\rowcolor{popblueL} \textbf{Conseil (Espagnol)} & \textbf{Structure grammaticale} & \textbf{Français} \\
\hline
\esp{Lávense las manos con jabón.} & Imp. affirm. pronominal (enclise \esp{se}) + CD nominal & Lavez-vous les mains avec du savon. \\
\hline
\esp{No las toquen; están sucias.} & Imp. négatif + CD proclitique \esp{las} + justification & Ne les touchez pas ; elles sont sales. \\
\hline
\esp{Separen la basura para reciclarla.} & Imp. affirm. + CD nominal + \esp{para} + inf. + enclise & Triez les déchets pour les recycler. \\
\hline
\esp{No la contaminen, por favor.} & Imp. négatif + CD proclitique \esp{la} + adverbe & Ne la polluez pas, s'il vous plaît. \\
\hline
\esp{Denle una oportunidad al cambio.} & Imp. affirm. + CI \esp{le} (enclise) + CD nominal & Donnez une chance au changement. \\
\hline
\end{tabularx}
\end{center}
\end{exemplebox}

\begin{attention}[Pièges majeurs pour le podcast : Cohérence et Accents]
\begin{enumerate}
    \item \textbf{Incohérence de personne (Le piège n°1).} Ne passez pas de \esp{ustedes} à \esp{tú} ou \esp{vosotros} au milieu du script.
    \begin{itemize}
        \item \esp{\textcolor{red}{Reciclen} (ustedes) el plástico y \textcolor{red}{no lo tires} (tú) al suelo.} $\leftarrow$ \textbf{INCORRECT}.
        \item \esp{\textcolor{popgreen}{Reciclen} el plástico y \textcolor{popgreen}{no lo tiren} al suelo.} $\leftarrow$ \textbf{CORRECT} (tout \cle{ustedes}).
    \end{itemize}
    \item \textbf{Oubli de l'accent écrit à l'enclise (Le piège n°2).} L'ajout de pronoms allonge le mot ; l'accent tonique reste sur la même syllabe du verbe $\rightarrow$ accent écrit obligatoire si la règle générale l'exige (mots de plus de 2 syllabes terminés par voyelle, \esp{n}, \esp{s} dont la tonique n'est pas sur la pénultième).
    \begin{itemize}
        \item \esp{Reciclen + los = \textcolor{popgreen}{Reciclénlos}} (verbe \esp{re-ci-clen} 3 syll., tonique sur \emph{clen} ; +\emph{los} = 4 syll., tonique sur \emph{clen} = antépénultième $\rightarrow$ accent sur \emph{é}).
        \item \esp{Digan + me + lo = \textcolor{popgreen}{Díganmelo}} (verbe \esp{di-gan} 2 syll., tonique sur \emph{di} ; +\emph{me-lo} = 4 syll., tonique sur \emph{di} = pré-antépénultième $\rightarrow$ accent sur \emph{í}).
        \item \esp{Cuida + te = \textcolor{popgreen}{Cuídate}} (verbe \esp{cui-da} 2 syll., tonique sur \emph{cui} ; +\emph{te} = 3 syll., tonique sur \emph{cui} = antépénultième $\rightarrow$ accent sur \emph{í}).
        \item \esp{Hiervan + la = \textcolor{popgreen}{hiervanla}} (verbe \esp{hier-van} 2 syll., tonique sur \emph{van} ; +\emph{la} = 3 syll., tonique sur \emph{van} = pénultième, terme en voyelle $\rightarrow$ \textbf{pas d'accent}).
        \item \esp{Cuídense} (verbe \esp{cuí-den} 2 syll., tonique sur \emph{cuí} ; +\emph{se} = 3 syll., tonique sur \emph{den} = pénultième, terme en voyelle $\rightarrow$ \textbf{pas d'accent} sur le verbe, l'accent de \esp{cuíden} tombe).
    \end{itemize}
    \item \textbf{Ordre des pronoms inversé (Le piège n°3).} Jamais \esp{lo se}, toujours \esp{se lo} / \esp{se la} / \esp{se los} / \esp{se las}.
    \item \textbf{Traduction littérale du français « vous » (Le piège n°4).} « Lavez-vous les mains » $\neq$ \esp{Laven ustedes las manos} (lourd). Préférer \textbf{impératif pronominal} : \esp{Lávense las manos} (ustedes) ou \esp{Lavaos las manos} (vosotros) / \esp{Lávate las manos} (tú).
\end{enumerate}
\end{attention}

\courssub{Activité de production : Rédaction en binôme (8–12 lignes)}

Vous allez rédiger le script complet d'un mini-podcast (1 min 30 env.) sur l'un des trois thèmes imposés. Travaillez en binôme : l'un rédige, l'autre relit avec la grille ; puis inversez les rôles pour un second brouillon si le temps le permet.

\begin{experience}[Atelier d'écriture : « Mon podcast éco-santé »]
\textbf{Consignes :}
\begin{enumerate}
    \item \textbf{Choisir un thème} (un seul par binôme) :
    \begin{itemize}
        \item \textbf{A. Hygiène au quartier} (gestion des ordures, eaux stagnantes, moustiques, propreté des caniveaux).
        \item \textbf{B. Lutte contre le paludisme} (moustiquaires, insecticides, herbes, consultation précoce).
        \item \textbf{C. Zéro plastique au lycée} (gourdes, sacs réutilisables, tri, sensibilisation camarades).
    \end{itemize}
    \item \textbf{Rédiger le script} (8 à 12 lignes de dialogue/texte continu) en respectant \textbf{obligatoirement} :
    \begin{itemize}
        \item Structure : \esp{Introducción} $\rightarrow$ \esp{3-5 Consejos} $\rightarrow$ \esp{Conclusión} $\rightarrow$ \esp{Despedida}.
        \item \textbf{Minimum grammatical} : 4 impératifs affirmatifs + 2 impératifs négatifs + 3 pronoms enclitiques (sur impératifs) + 2 pronoms proclitiques.
        \item Cohérence de personne (\cle{ustedes} recommandé).
        \item Lexique thématique spécifique (voir boîte lexique ci-dessous).
    \end{itemize}
    \item \textbf{Relecture croisée} : Échangez votre brouillon avec un autre binôme. Utilisez la \textbf{Grille d'auto-évaluation} (ci-dessous) pour cocher les items. Surlignez en couleur : \textcolor{popgreen}{Vert} = OK, \textcolor{poporange}{Orange} = À vérifier (accord, accent), \textcolor{red}{Rouge} = Erreur (personne, place pronom).
    \item \textbf{Rédaction finale} : Intégrez les corrections. Préparez la lecture à voix haute (chronométrez : 1 min 30 max).
\end{enumerate}
\end{experience}

\begin{definition}[Lexique thématique utile par thème]
\begin{center}
\begin{tabularx}{\linewidth}{|X|X|X|}
\hline
\rowcolor{popblueL} \textbf{Thème A : Hygiène quartier} & \textbf{Thème B : Paludisme} & \textbf{Thème C : Zéro plastique lycée} \\
\hline
\esp{el desagüe / la cuneta} (caniveau) & \esp{la mosquitera} (moustiquaire) & \esp{la botella reutilizable} (gourde) \\
\esp{el agua estancada} (eau stagnante) & \esp{el repelente} (répulsif) & \esp{la bolsa de tela} (sac en tissu) \\
\esp{la basura / los desechos} & \esp{la malaria / el paludismo} & \esp{el táper / la fiambrera} (boîte repas) \\
\esp{quemar / enterrar} (brûler/enterrer) & \esp{la fiebre / los escalofríos} & \esp{el envoltorio} (emballage) \\
\esp{el cubo de basura} (poubelle) & \esp{el centro de salud} (CS) & \esp{la cantina} (cantine) \\
\esp{el dengue / la malaria} & \esp{la consulta temprana} & \esp{concienciar / concientizar} (sensibiliser) \\
\hline
\end{tabularx}
\end{center}
\textit{Note : \esp{el caniveo} est un régionalisme camerounais (de \emph{caniveau}), à éviter dans l'écrit standard.}
\textbf{Verbes d'action fréquents :} \esp{eliminar, tapar, limpiar, drenar, usar, colocar, acudir, tomar, reducir, reutilizar, rechazar, separar, depositar, evitar, cambiar, promover.}
\end{definition}

\begin{savaistu}[L'oral au Bac et le format Podcast]
L'épreuve d'expression orale du Baccalauréat (LV2) valorise la \cle{production authentique} et la \cle{prise de parole en continu}. Les programmes (MINESEC) mentionnent explicitement les « \emph{nouvelles technologies de l'information et de la communication} » et les « \emph{supports variés} ». Le podcast est un format parfait : il permet d'évaluer la \cle{prononciation}, l'\cle{intelligibilité}, la \cle{cohérence}, la \cle{richesse lexicale} et la \cle{correction grammaticale} (ici l'impératif/pronoms) sans le stress de l'interaction immédiate. Un podcast bien structuré, enregistré ou présenté en direct, est une trace évaluable excellente pour le dossier de l'élève ou l'examen pratique.
\end{savaistu}

\courssub{Grille d'auto-évaluation et d'évaluation formative (Binôme / Professeur)}

Cette grille sert à la relecture croisée (étape 3 de l'atelier) et à l'évaluation finale par le professeur lors de l'oral/enregistrement.

\begin{exoresolu}[Application de la grille : Exemple de correction croisée]
\textbf{Extrait élève (Thème C) :} \esp{« Hola estudiantes. Bienvenidos. Primer consejo: usen botellas de plastico, no los compren. Segundo: traigan su tupper, no lo olviden. Tercero: reciclen las botellas, separenlas bien. ¡Cuíden el medio ambiente! Adiós. »}

\textbf{Analyse avec la grille :}
\begin{itemize}
    \item \textbf{Grammaire :} \esp{usen} (OK affirmatif ustedes), \esp{no los compren} (OK négatif + proclise \esp{los}), \esp{traigan} (OK irrégulier \esp{traer}), \esp{no lo olviden} (OK négatif + proclise \esp{lo}), \esp{reciclen} (OK), \esp{separenlas} (OK enclise + accent \esp{se-pa-ren-las} $\rightarrow$ \esp{sepárenlas} \textcolor{red}{ERREUR ACCENT : doit être \esp{sepárenlas}}), \esp{Cuíden} (OK pronominal \esp{cuiden} + \esp{se} $\rightarrow$ \esp{cuídense} \textcolor{red}{ERREUR PRONOM MANQUANT : \esp{Cuíden} n'est pas pronominal, il faut \esp{Cuídense}}).
    \item \textbf{Lexique :} \esp{plastico} $\rightarrow$ \esp{plástico} (accent). \esp{tupper} (anglicisme, accepter \esp{táper} ou \esp{fiambrera}). \esp{medio ambiente} (OK).
    \item \textbf{Contraintes :} Affirmatifs : \esp{usen, traigan, reciclen, separenlas, cuíden} (5 OK). Négatifs : \esp{no los compren, no lo olviden} (2 OK). Enclitiques : \esp{separenlas} (1 seul, \textcolor{red}{manque 2}). Proclitiques : \esp{los, lo} (2 OK).
    \item \textbf{Action corrective :} Ajouter enclitiques : \esp{úsenlas} (conseil 1), \esp{tráiganlo} (conseil 2). Corriger \esp{sepárenlas}, \esp{cuídense}.
\end{itemize}
\textbf{Version corrigée :} \esp{« ... Primer consejo: \textbf{úsenlas} (botellas reutilizables), \textbf{no las compren} (desechables). Segundo: \textbf{tráiganlo} (su táper), \textbf{no lo olviden}. Tercero: \textbf{recíclenlas}, \textbf{sepárenlas} bien. \textbf{¡Cuídense} y \textbf{cuídenlo} (el planeta)! ... »}
\end{exoresolu}

\begin{center}
\begin{tabularx}{\linewidth}{|>{\hsize=1.2\hsize}X|>{\hsize=0.8\hsize}X|>{\hsize=0.8\hsize}X|>{\hsize=1.2\hsize}X|}
\hline
\rowcolor{popblueL} \textbf{Critère} & \textbf{Indicateur observable} & \textbf{Poids} & \textbf{Commentaire / Note (sur 4)} \\
\hline
\textbf{Exactitude grammaticale} & Respect contraintes : 4 Imp. Aff + 2 Imp. Neg + 3 Enclitiques + 2 Proclitiques. Bonne conjugaison (irréguliers : \esp{ir $\rightarrow$ id, vaya ; ser $\rightarrow$ sé, sea ; haber $\rightarrow$ he, haya ; poner $\rightarrow$ pon, ponga ; tener $\rightarrow$ ten, tenga ; venir $\rightarrow$ ven, venga ; decir $\rightarrow$ di, diga ; salir $\rightarrow$ sal, salga}). Accents écrites enclises corrects. & 30\% & \textcolor{popgreen}{4} / \textcolor{poporange}{3} / \textcolor{red}{2} / \textcolor{red}{1} \\
\hline
\textbf{Cohérence personne} & Une seule personne tout du long (ustedes recommandé). Pas de mélange tú/ustedes/vosotros. & 15\% & \textcolor{popgreen}{4} / \textcolor{poporange}{3} / \textcolor{red}{2} / \textcolor{red}{1} \\
\hline
\textbf{Lexique thématique} & Minimum 8 mots/spécifiques au thème choisi (voir boîte lexique). Variété verbes d'action. Pas de franglais (ex: \esp{reciclar} pas \esp{recycler}). & 20\% & \textcolor{popgreen}{4} / \textcolor{poporange}{3} / \textcolor{red}{2} / \textcolor{red}{1} \\
\hline
\textbf{Structure du podcast} & Introduction claire (Nom, Sujet). 3-5 conseils distincts (connecteurs \esp{Primero, Segundo...}). Justifications courtes. Conclusion engageante. Despedida. & 15\% & \textcolor{popgreen}{4} / \textcolor{poporange}{3} / \textcolor{red}{2} / \textcolor{red}{1} \\
\hline
\textbf{Fluidité orale / Prosodie} & Débit naturel, pas de lecture mot-à-mot. Intelligibilité (accents toniques, \esp{ñ}, \esp{ll}, \esp{r/rr}, \esp{j/g}). Intonation impérative (descendante affirmative, montante-descendante négative/insistance). Respect temps (1'30 $\pm$ 15\%). & 20\% & \textcolor{popgreen}{4} / \textcolor{poporange}{3} / \textcolor{red}{2} / \textcolor{red}{1} \\
\hline
\end{tabularx}
\end{center}

\courssub{Enregistrement ou présentation orale : Mise en voix}

Le passage à l'oral est l'aboutissement. Deux modalités selon l'équipement de l'établissement :

\begin{experience}[Modalités de production orale]
\textbf{Option A : Enregistrement audio (Préféré si matériel dispo).}
\begin{itemize}
    \item Outils : Dictaphone, smartphone (enregistreur vocal), logiciel libre \emph{Audacity} (PC salle info), application \emph{OnBuch+} (module oral).
    \item Consigne : Enregistrer en \textbf{une seule prise} (conditions d'examen). Si erreur majeure $\rightarrow$ recommencer.
    \item Post-production (optionnel) : Ajout d'un jingle d'intro/outro (musique libre de droit), bruitage (oiseaux, eau, foule) pour réalisme.
    \item Remise : Fichier \texttt{.mp3} ou \texttt{.wav} nommé \texttt{Nom\_Prenom\_Podcast\_ThemeX.mp3} sur clé USB / plateforme école / OnBuch+.
\end{itemize}

\textbf{Option B : Présentation en direct devant la classe (Classique).}
\begin{itemize}
    \item Disposition : Binôme face à la classe (micro fictif ou vrai micro sur pied).
    \item Public : Les autres élèves notent sur grille simplifiée (3 items : \cle{Compréhension globale}, \cle{Correction impératifs/pronoms}, \cle{Intérêt/Originalité}).
    \item Professeur : Note sur grille complète (ci-dessus) + feedback oral immédiat (1 min) : « Point fort / Point à travailler ».
\end{itemize}

\textbf{Conseils de prosodie pour l'impératif à l'oral :}
\begin{itemize}
    \item \textbf{Affirmatif} : Intonation \textbf{descendante} finale (ordre/consigne ferme). \esp{Reciclénlos $\downarrow$}.
    \item \textbf{Négatif} : Intonation souvent \textbf{montante sur \esp{no}} puis descendante, ou \textbf{insistante} sur le verbe. \esp{No lo tiren $\downarrow$} (interdiction nette).
    \item \textbf{Enclise} : Bien lier le pronom au verbe (pas de pause). \esp{Se-pá-ren-las} (syllabes liées). L'accent tonique porte sur la syllabe accentuée graphiquement.
    \item \textbf{Rythme} : Marquer les virgules / points-virgules par une micro-pause pour aérer les conseils.
\end{itemize}
\end{experience}

\courssub{Schéma récapitulatif : Structure visuelle du podcast}

\begin{popfigure}
\begin{tikzpicture}[
    node distance=1.2cm and 1.5cm,
    box/.style={draw=popblue, thick, rounded corners, fill=popblueL, text width=3.2cm, align=center, minimum height=1.8cm, font=\small},
    arrow/.style={-Stealth, thick, popdark},
    labelbox/.style={font=\tiny\itshape, text=popmuted, anchor=north}
]
% Nœuds principaux
\node[box, align=center] (intro){\textbf{INTRODUCCIÓN}\\\esp{Saludo + Nombre podcast + Locutor + Tema}};
\node[box, right=of intro, align=center] (c1){\textbf{CONSEJO 1}\\\esp{Imperativo + Pronombre + Justificación}};
\node[box, right=of c1, align=center] (c2){\textbf{CONSEJO 2}\\\esp{Imperativo + Pronombre + Justificación}};
\node[box, right=of c2, align=center] (c3){\textbf{CONSEJO 3}\\\esp{Imperativo + Pronombre + Justificación}};
\node[box, below=of c2, align=center] (conc){\textbf{CONCLUSIÓN}\\\esp{Llamada a la acción global + Valores}};
\node[box, right=of conc, align=center] (desp){\textbf{DESPEDIDA}\\\esp{Fórmula cierre + Próximo episodio}};

% Flèches horizontales
\draw[arrow] (intro) -- (c1) node[midway, above, labelbox] {Presenta};
\draw[arrow] (c1) -- (c2) node[midway, above, labelbox] {Enumera};
\draw[arrow] (c2) -- (c3) node[midway, above, labelbox] {Desarrolla};
% Flèches vers conclusion
\draw[arrow] (c3) -- ++(0,-0.6) -| (conc.north) node[pos=0.25, right, labelbox] {Sintetiza};
\draw[arrow] (conc) -- (desp) node[midway, above, labelbox] {Cierra};

% Annotations grammaticales (côté gauche)
\node[left=0.5cm of intro, align=right, font=\scriptsize, text=poppurple] (g1) {\textcolor{poppurple}{Gramática:}\\{Ustedes / Presente indicativo (traigo)\\\textit{No imperativo aquí}};
\node[left=0.5cm of c1, align=right, font=\scriptsize, text=popgreen] (g2) {\textcolor{popgreen}{Gramática:}\\{Imperativo Aff/Neg\\\textbf{Enclise / Proclise}\\\textit{+ Justificación (pq/para)}};
\node[left=0.5cm of conc, align=right, font=\scriptsize, text=poporange] (g3) {\textcolor{poporange}{Gramática:}\\{Imperativo Aff\\\textbf{Reflexivo (cuídense)}\\\textit{+ Objeto directo (cuídenlo)}};

% Décoration
\draw[popline, line width=1.5pt, rounded corners] ($(intro.north west)+(-0.3,0.3)$) rectangle ($(desp.south east)+(0.3,-0.3)$);
\node[above=0.1cm of intro.north, font=\bfseries\large, text=popblue] {ESTRUCTURA PODCAST « SALUD / ECO »};
\end{tikzpicture}
\legende{Schéma de la structure narrative et grammaticale d'un podcast conseil. Les flèches indiquent la progression logique ; les encadrés latéraux rappellent les points grammaticaux clés attendus par section.}
\end{popfigure}

\courssub{Exercices d'application finale (Entraînement individuel)}

\begin{exercice}[Transformation et complétion : Préparez vos « cartes » pour l'oral]
\textbf{A. Complétez les conseils avec l'impératif \cle{ustedes} et les pronoms requis (affirmatif $\rightarrow$ enclise / négatif $\rightarrow$ proclise). Mettez l'accent écrit si nécessaire.}
\begin{enumerate}
    \item (Reciclar + los) \esp{\_\_\_\_\_\_\_\_\_\_\_\_ las botellas de plástico en el contenedor amarillo.}
    \item (No / tirar / las) \esp{\_\_\_\_\_\_\_\_\_\_\_\_ al suelo; \_\_\_\_\_\_\_\_\_\_\_\_ (depositar + las) en la papelera.}
    \item (Hervir + la) \esp{\_\_\_\_\_\_\_\_\_\_\_\_ el agua diez minutos antes de \_\_\_\_\_\_\_\_\_\_\_\_ (beber + la).}
    \item (No / desperdiciar / la) \esp{\_\_\_\_\_\_\_\_\_\_\_\_; \_\_\_\_\_\_\_\_\_\_\_\_ (reutilizar + la) para regar las plantas.}
    \item (Poner / se) \esp{\_\_\_\_\_\_\_\_\_\_\_\_ repelente al atardecer y \_\_\_\_\_\_\_\_\_\_\_\_ (usar + la) mosquitera.}
    \item (No / olvidar / se) \esp{\_\_\_\_\_\_\_\_\_\_\_\_ de tomar la profilaxis si el médico \_\_\_\_\_\_\_\_\_\_\_\_ (recetar + se la).}
\end{enumerate}

\textbf{B. Réécrivez les phrases suivantes en changeant la personne de \cle{tú} à \cle{ustedes} (adaptez verbes et pronoms).}
\begin{enumerate}
    \item \esp{Cuida tu higiene; lávate las manos y no te toques la cara.}
    \item \esp{Recicla tus botellas; sepáralas y no las tires.}
    \item \esp{Planta un árbol; riégalo y no lo abandones.}
\end{enumerate}
\end{exercice}

\begin{corrige}[Exercice A et B]
\textbf{A.}
\begin{enumerate}
    \item \esp{\textbf{Reciclénlas} las botellas...} (Affirmatif + enclise \esp{los} $\rightarrow$ \esp{las}, accent \esp{Re-ci-clén-las} : antépénultième).
    \item \esp{\textbf{No las tiren}} al suelo; \esp{\textbf{deposítenlas}} en la papelera. (Négatif proclise \esp{las} ; Affirmatif enclise \esp{depositen + las} $\rightarrow$ \esp{deposítenlas}, accent sur \emph{é} : antépénultième).
    \item \esp{\textbf{Hiervanla}} el agua diez minutos antes de \esp{\textbf{beberla}}. (Affirmatif \esp{hiervan + la} $\rightarrow$ \esp{hiervanla} \textbf{sans accent} : pénultième \emph{van}, terme en voyelle ; Infinitif + enclise \esp{beber + la} $\rightarrow$ \esp{beberla} \textbf{sans accent} : pénultième \emph{ber}, terme en voyelle).
    \item \esp{\textbf{No la desperdicien}}; \esp{\textbf{reutilicénla}} para regar las plantas. (Négatif proclise \esp{la} ; Affirmatif \esp{reutilicen + la} $\rightarrow$ \esp{reutilicénla}, accent sur \emph{é} : antépénultième).
    \item \esp{\textbf{Pónganse}} repelente... y \esp{\textbf{usenla}} mosquitera. (Irrégulier \esp{poner} $\rightarrow$ \esp{pongan} + \esp{se} $\rightarrow$ \esp{pónganse}, accent sur \emph{ó} : antépénultième ; \esp{usar} $\rightarrow$ \esp{usen + la} $\rightarrow$ \esp{usenla} \textbf{sans accent} : pénultième \emph{sen}, terme en voyelle).
    \item \esp{\textbf{No se olviden}} de tomar la profilaxis si el médico \esp{\textbf{se la ha recetado}} (passé composé indicatif) / \esp{\textbf{se la recete}} (subjonctif présent si hypothétique futur). \textit{Note : l'impératif ne s'utilise pas dans une subordonnée \esp{si}.}
\end{enumerate}

\textbf{B. (Passage \cle{tú} $\rightarrow$ \cle{ustedes})}
\begin{enumerate}
    \item \esp{\textbf{Cuiden} su higiene; \textbf{lávense} las manos y \textbf{no se toquen} la cara.}
    \item \esp{\textbf{Reciclen} sus botellas; \textbf{sepárenlas} y \textbf{no las tiren}.}
    \item \esp{\textbf{Planten} un árbol; \textbf{rieguenlo} y \textbf{no lo abandonen}.} (\esp{regar} $\rightarrow$ \esp{rieguen} + \esp{lo} $\rightarrow$ \esp{rieguenlo} \textbf{sans accent} : pénultième \emph{guen}, terme en voyelle).
\end{enumerate}
\end{corrige}

\begin{aretenir}[Check-list finale avant l'oral / l'enregistrement]
\begin{itemize}
    \item [\checkmark] Mon script fait 8-12 lignes, structure complète (Intro, 3-5 conseils, Conclusion, Despedida).
    \item [\checkmark] J'ai \textbf{4 impératifs affirmatifs} minimum (surlignés en vert).
    \item [\checkmark] J'ai \textbf{2 impératifs négatifs} minimum (surlignés en rouge).
    \item [\checkmark] J'ai \textbf{3 pronoms enclitiques} (affirmatifs : \esp{...melo, ...selo, ...noslas}...) avec \textbf{accents écrits vérifiés}.
    \item [\checkmark] J'ai \textbf{2 pronoms proclitiques} (négatifs : \esp{no me lo, no se la, no nos los}...).
    \item [\checkmark] \textbf{Une seule personne} (\esp{ustedes}) du début à la fin.
    \item [\checkmark] Lexique thème (8 mots) intégré, verbes variés (irréguliers inclus).
    \item [\checkmark] Relu par un pair (grille remplie) + corrigé.
    \item [\checkmark] Lu à voix haute 3 fois : chrono OK (1'30), fluidité OK, intonation impérative travaillée.
\end{itemize}
\textbf{Bravo ! Vous êtes prêts pour l'émission. ¡En vivo, muchachos!}
\end{aretenir}

\newpage

\coursec{Activité d'intégration}

\courssub{Objectifs de l'activité}

À l'issue de cette activité d'intégration, l'élève sera capable de :
\begin{itemize}
\item mobiliser dans une production orale authentique l'\cle{impératif} affirmatif et négatif (toutes les personnes : \esp{tú}, \esp{vosotros}, \esp{usted}, \esp{ustedes}) ;
\item placer correctement les \cle{pronoms compléments} : enclise à l'impératif affirmatif (\esp{dámelo}, \esp{cuídenlo}) et proclise à l'impératif négatif (\esp{no me lo des}, \esp{no lo tiren}), en respectant l'accentuation écrite ;
\item utiliser le lexique de la \cle{santé} et de l'\cle{environnement} étudié dans la leçon ;
\item collaborer en groupe pour rédiger, répéter et présenter un \cle{podcast} de deux minutes destiné à la radio du lycée ;
\item développer une écoute active et critique grâce à une fiche d'écoute.
\end{itemize}

\courssub{Situation}

Le club environnement de votre lycée de Yaoundé participe à la campagne trimestrielle « \esp{Un barrio sano y limpio} », lancée en partenariat avec une association de Guinée équatoriale. À l'occasion de la semaine culturelle, la radio du lycée souhaite diffuser chaque midi un court podcast de sensibilisation, réalisé par les élèves de Terminale A. Votre groupe a été sélectionné pour enregistrer le premier épisode, qui s'adressera à tous les élèves et au personnel du lycée : il faudra donner des conseils de santé et des consignes écologiques, en espagnol, de manière claire, dynamique et convaincante.

Pour vous inspirer, voici le message officiel de la campagne, affiché au tableau d'information du lycée :

\begin{textebox}[Mensaje oficial de la campaña]
¡Atención, jóvenes! Nuestro barrio necesita tu ayuda. Cada día tiramos plástico al suelo, quemamos basura cerca de las casas y olvidamos lavarnos las manos antes de comer. Pero cambiar es fácil: recoge los residuos, sepáralos, no los quemes. Bebe agua limpia, hiérvela si es necesario. Cuida tu cuerpo: haz deporte, duerme bien, no fumes. Y sobre todo, no te olvides: un barrio limpio es un barrio sano. ¡Participa! Graba tu mensaje y envíalo a la radio del liceo. ¡Te escuchamos!
\end{textebox}

\textbf{Glossaire :} \esp{tirar} = jeter ; \esp{quemar basura} = brûler des ordures ; \esp{recoger los residuos} = ramasser les déchets ; \esp{sepáralos} = sépare-les (impératif + pronom enclitique) ; \esp{hervir el agua} = faire bouillir l'eau ; \esp{hiérvela} = fais-la bouillir ; \esp{fumar} = fumer ; \esp{grabar} = enregistrer ; \esp{envíalo} = envoie-le.

\courssub{Consignes}

\begin{enumerate}
\item \textbf{Compréhension.} Lisez le message officiel de la campagne. Soulignez tous les impératifs et classez-les en deux colonnes : affirmatifs et négatifs. Entourez les pronoms compléments et dites s'ils sont enclitiques ou proclitiques.
\item \textbf{Lexique.} En groupe, constituez deux listes de huit mots chacune : (a) santé (\esp{la salud, el médico, lavarse las manos, la enfermedad...}) ; (b) environnement (\esp{el medio ambiente, la basura, reciclar, la contaminación...}). Vous les réutiliserez dans votre script.
\item \textbf{Grammaire.} Transformez les consignes suivantes à l'impératif avec pronom, selon le modèle : \esp{Debes beber el agua} $\rightarrow$ \esp{Bébela.} / \esp{No debes tirar la basura} $\rightarrow$ \esp{No la tires.} Consignes à transformer : \esp{hay que lavarse las manos ; hay que proteger los árboles ; no hay que quemar los plásticos ; hay que cuidar el jardín del liceo}.
\item \textbf{Rédaction du script.} Par groupes de 3 ou 4, rédigez le script d'un podcast de deux minutes (environ 180 à 220 mots) intitulé \esp{Un barrio sano y limpio}. Il doit obligatoirement contenir : au moins \textbf{4 impératifs affirmatifs}, \textbf{2 impératifs négatifs}, \textbf{3 pronoms enclitiques} (ex. \esp{cuídenlo}, \esp{bébanla}) et \textbf{2 pronoms proclitiques} (ex. \esp{no lo tiren}, \esp{no se olviden}). Prévoyez une introduction avec le nom de l'émission, deux ou trois voix différentes, et une chute percutante.
\item \textbf{Répétition et présentation.} Répétez en veillant à la prononciation (accent tonique respecté malgré l'enclise : \esp{CUI-den-lo}). Enregistrez le podcast sur un téléphone ou présentez-le en direct devant la classe.
\item \textbf{Écoute active.} Pendant la présentation des autres groupes, remplissez la fiche d'écoute : conseils entendus, impératifs relevés (affirmatifs / négatifs), pronoms repérés, erreurs éventuelles.
\item \textbf{Vote.} La classe élit le meilleur podcast selon trois critères : correction grammaticale, richesse du lexique, qualité de l'expression orale. Le podcast gagnant sera « diffusé » lors de la semaine culturelle.
\end{enumerate}

\courssub{Ressources à mobiliser}

\begin{aretenir}[Impératif et pronoms : l'essentiel]
\begin{itemize}
\item Impératif affirmatif : \esp{tú} $\rightarrow$ 3\up{e} personne du singulier du présent (\esp{bebe}, \esp{recoge}) ; irréguliers : \esp{di, haz, ve, pon, sal, sé, ten, ven}. \esp{usted/ustedes} $\rightarrow$ subjonctif présent (\esp{beba}, \esp{beban}) ; \esp{vosotros} $\rightarrow$ infinitif sans -r + -d (\esp{bebed}).
\item Impératif négatif : \esp{no} + subjonctif présent (\esp{no bebas}, \esp{no beba}, \esp{no bebáis}, \esp{no beban}).
\item Affirmatif : le pronom se \textbf{colle} après le verbe (enclise) : \esp{bebe + la} $\rightarrow$ \esp{bébela} ; \esp{cuiden + lo} $\rightarrow$ \esp{cuídenlo}. On ajoute un accent écrit pour conserver l'accent tonique.
\item Négatif : le pronom se \textbf{place devant} (proclise) : \esp{no la bebas}, \esp{no lo tiren}, \esp{no se olviden}.
\item Deux pronoms : l'indirect avant le direct : \esp{dámelo} ; \esp{se} remplace \esp{le/les} devant \esp{lo/la/los/las} : \esp{dígaselo}.
\end{itemize}
\end{aretenir}

\begin{definition}[Lexique utile]
Santé : \esp{la salud, enfermarse, lavarse las manos, el agua potable, hacer deporte, dormir bien, el mosquito, el paludismo}. Environnement : \esp{el medio ambiente, la basura, los residuos, reciclar, reutilizar, plantar árboles, la contaminación, el plástico, proteger, limpiar}. Présentation radio : \esp{Bienvenidos a...}, \esp{Hoy hablamos de...}, \esp{Escuchen nuestros consejos}, \esp{Hasta la próxima}.
\end{definition}

\courssub{Proposition de production et corrigé commenté}

\begin{exoresolu}[Podcast modèle : « Un barrio sano y limpio »]
Rédigez un script de podcast respectant les contraintes de la consigne 4 (4 impératifs affirmatifs, 2 négatifs, 3 pronoms enclitiques, 2 proclitiques).
\tcblower
\textbf{Production modèle (script à trois voix) :}

\esp{LOCUTOR 1 : ¡Hola, queridos oyentes! Bienvenidos a « Salud y medio ambiente », el podcast del liceo. Hoy hablamos de nuestro barrio: ¿cómo vivir mejor? Escuchen nuestros consejos.}

\esp{LOCUTORA 2 : Primero, la salud. Lávense las manos antes de comer, es fundamental. El agua del pozo puede ser peligrosa: hiérbanla antes de beberla. Hagan deporte cada día y duerman al menos ocho horas. Y por favor, no fumen: el tabaco destruye los pulmones.}

\esp{LOCUTOR 3 : Ahora, el medio ambiente. La basura no pertenece a la calle: recójanla y sepárenla. Los plásticos contaminan la tierra y los ríos: no los quemen, recíclenlos. Planten un árbol cerca de su casa y cuídenlo con amor.}

\esp{LOCUTOR 1 : Recuerden: un barrio limpio es un barrio sano. No se olviden de nuestros consejos y pónganlos en práctica hoy mismo. ¡Hasta la próxima semana!}

\textbf{Commentaire des choix de langue :}
\begin{itemize}
\item \textbf{Impératifs affirmatifs (bien plus de 4) :} \esp{escuchen}, \esp{hagan}, \esp{duerman}, \esp{planten}, \esp{recuerden} — tous à la personne \esp{ustedes}, adaptée à un public collectif (les élèves du lycée) ; c'est le subjonctif présent sans \esp{no}.
\item \textbf{Impératifs négatifs (2) :} \esp{no fumen} et \esp{no los quemen} — subjonctif présent précédé de \esp{no}.
\item \textbf{Pronoms enclitiques (4, avec accent écrit) :} \esp{hiérbanla} (l'eau, féminin singulier : \esp{la}), \esp{recójanla} et \esp{sepárenla} (la basura), \esp{cuídenlo} (el árbol, masculin : \esp{lo}), \esp{pónganlos} (los consejos). Remarquez l'accent ajouté : \esp{hiervan} $\rightarrow$ \esp{hiérbanla}, pour conserver l'accent tonique d'origine.
\item \textbf{Pronoms proclitiques (3) :} \esp{no los quemen} (les plásticos), \esp{no se olviden} (verbe pronominal : le pronom réfléchi \esp{se} précède le verbe négatif), \esp{antes de beberla} montre aussi l'enclise sur l'infinitif.
\item \textbf{Lexique thématique :} \esp{el pozo, el tabaco, los pulmones, reciclar, plantar un árbol, contaminar} — champ lexical santé/environnement du module.
\item \textbf{Format radio :} salutation (\esp{queridos oyentes}), annonce du thème, transitions (\esp{primero}, \esp{ahora}), chute mémorable (\esp{un barrio limpio es un barrio sano}, reprise du slogan de la campagne).
\end{itemize}
\end{exoresolu}

\begin{attention}[Erreurs à éviter dans vos scripts]
\begin{itemize}
\item Ne dites pas \esp{no tireslo} : au négatif, le pronom passe \textbf{devant} : \esp{no lo tires}.
\item N'oubliez pas l'accent écrit à l'enclise : \esp{cuidenlo} est fautif, écrivez \esp{cuídenlo}.
\item Attention à l'impératif \esp{tú} irrégulier : \esp{haz} (pas \esp{hace}), \esp{pon} (pas \esp{pone}), \esp{sal} (pas \esp{sale}).
\item \esp{No te olvides} garde le \esp{se/te} devant : jamais \esp{no olvídese} à la forme négative.
\end{itemize}
\end{attention}

\courssub{Bilan de l'activité}

Cette activité d'intégration a permis de réinvestir l'ensemble des savoir-faire de la leçon dans une situation de communication authentique et motivante. Sur le plan grammatical, les élèves ont manipulé l'impératif affirmatif et négatif à plusieurs personnes et vérifié concrètement la règle d'or : \textbf{enclise à l'affirmatif, proclise au négatif}, avec l'accentuation écrite qui en découle. Sur le plan lexical, le vocabulaire de la santé et de l'environnement a été employé dans un contexte réel et proche du vécu des élèves (eau du puits, gestion des ordures, propreté du quartier). Sur le plan méthodologique, la fiche d'écoute a développé l'écoute active et l'esprit critique, tandis que le vote final a valorisé la qualité de la langue autant que la créativité. Enfin, la dimension collaborative et la « diffusion » du podcast gagnant lors de la semaine culturelle donnent un sens concret à l'apprentissage de l'espagnol : parler pour informer, conseiller et convaincre, comme le font chaque jour les médias du monde hispanophone, de l'Espagne à la Guinée équatoriale.

\newpage

\coursec{Méthodes et exercices résolus}

\coursec{Leçon E08 : Traduction : la phrase impérative avec ou sans pronoms compléments ; emisión de salud / podcast ecológico}

\courssub{Mise en route : une émission de santé au lycée}

\begin{textebox}[Script d'intro — Podcast « Salud y planeta », Club d'espagnol, Lycée de Yaoundé]
\esp{¡Buenos días, queridos oyentes! Bienvenidos a una nueva edición de « Salud y planeta », su programa de consejos desde el club de español del liceo. Soy su presentador, Carlos, y hoy tenemos un tema vital: ¿cómo cuidar nuestra salud y nuestro entorno sin complicarnos la vida? Escuchen bien: \textbf{beban} agua nada más levantarse, \textbf{coman} fruta de temporada, \textbf{muévanse} cada día. Y por favor, \textbf{no tiren} plásticos al suelo: \textbf{recíclenlos}. La Tierra y su cuerpo se lo agradecerán. ¡Empezamos!}
\tcblower
\textbf{Traduction :} « Bonjour, chers auditeurs ! Bienvenue dans une nouvelle édition de "Santé et planète", votre programme de conseils depuis le club d'espagnol du lycée. Je suis votre animateur, Carlos, et aujourd'hui nous avons un sujet vital : comment prendre soin de notre santé et de notre environnement sans se compliquer la vie ? Écoutez bien : buvez de l'eau dès le lever, mangez des fruits de saison, bougez chaque jour. Et s'il vous plaît, ne jetez pas de plastiques par terre : recyclez-les. La Terre et votre corps vous en remercieront. On commence !
\textbf{Glossaire :} \esp{oyentes} (auditeurs) ; \esp{entorno} (entourage, environnement) ; \esp{complicarnos la vida} (nous compliquer la vie) ; \esp{nada más levantarse} (dès le lever / juste après s'être levé) ; \esp{fruta de temporada} (fruits de saison) ; \esp{muévanse} (bougez-vous, impératif \esp{ustedes} de \esp{moverse}) ; \esp{recíclenlos} (recyclez-les, enclise COD \esp{los}) ; \esp{se lo agradecerán} (vous en remercieront, futur simple).
\end{textebox}

\courssub{Savoir-faire 1 : Former l'impératif affirmatif et négatif}

\begin{methode}[Conjuguer et traduire l'impératif : la démarche en 5 étapes]
Pour traduire ou produire une phrase impérative en espagnol, procède toujours ainsi :
\begin{enumerate}
\item \cle{Identifier le destinataire} : à qui s'adresse l'ordre ou le conseil ? À un ami (\esp{tú}), à un groupe d'amis (\esp{vosotros}), à une personne de politesse (\esp{usted}), à un public ou une audience (\esp{ustedes}) ? Dans un podcast, on utilise presque toujours \esp{ustedes} ou \esp{tú}.
\item \cle{Repérer le signe} : la consigne est-elle affirmative (« fais ») ou négative (« ne fais pas ») ? C'est le signe qui décide de la forme du verbe.
\item \cle{Choisir la forme verbale} :
\begin{itemize}
\item Affirmatif : \esp{tú} = 3\textsuperscript{e} personne du singulier du présent (\esp{habla, come, escribe}) ; \esp{vosotros} = infinitif dont le \esp{-r} final devient \esp{-d} (\esp{hablad, comed, escribid}) ; \esp{usted/ustedes} = subjonctif présent (\esp{hable, hablen ; coma, coman ; escriba, escriban}).
\item Négatif : \textbf{toujours} le subjonctif présent précédé de \esp{no} : \esp{no hables, no habléis, no hable, no hablen}.
\end{itemize}
\item \cle{Vérifier les irréguliers} de l'impératif affirmatif à \esp{tú} (à mémoriser par cœur) :
\begin{center}
\begin{tabular}{|l|l|l|}\hline
\rowcolor{popblueL} Verbe & Impératif \esp{tú} & Traduction \\ \hline
\esp{decir} & \esp{di} & Dis \\ \hline
\esp{hacer} & \esp{haz} & Fais \\ \hline
\esp{ir} & \esp{ve} & Va \\ \hline
\esp{poner} & \esp{pon} & Mets / Pose \\ \hline
\esp{salir} & \esp{sal} & Sors \\ \hline
\esp{ser} & \esp{sé} & Sois \\ \hline
\esp{tener} & \esp{ten} & Aie \\ \hline
\esp{venir} & \esp{ven} & Viens \\ \hline
\end{tabular}
\end{center}
\item \cle{Traduire en français} avec l'impératif français correspondant, sans sujet exprimé : \esp{Bebe agua} = « Bois de l'eau » ; \esp{No fuméis} = « Ne fumez pas ».
\end{enumerate}
\end{methode}

\begin{propriete}[Impératif affirmatif vs négatif — Synthèse visuelle]
\begin{center}
\begin{tikzpicture}[node distance=1.2cm and 1.5cm, every node/.style={font=\small, align=center}]
\node[draw, rounded corners, fill=popgreenL, thick, align=center] (aff){\textbf{AFFIRMATIF}\\\esp{tú: habla / di}\\\esp{vosotros: hablad}\\\esp{usted: hable}\\\esp{ustedes: hablen}};
\node[draw, rounded corners, fill=poporangeL, thick, right=of aff, align=center] (neg){\textbf{NÉGATIF}\\\esp{no + subjonctif}\\\esp{tú: no hables}\\\esp{vosotros: no habléis}\\\esp{usted: no hable}\\\esp{ustedes: no hablen}};
\draw[->, thick, popdark] (aff) -- node[above] {Signe : \textbf{+}} (neg);
\draw[<-, thick, popdark] (aff) -- node[below] {Signe : \textbf{-}} (neg);
\end{tikzpicture}
\end{center}
\emph{Règle d'or :} Le négatif \textbf{impose} le subjonctif pour toutes les personnes. Pas d'exception.
\end{propriete}

\begin{exoresolu}[Former les impératifs — Type Bac]
Énoncé. Tu animes une émission de santé au lycée de Yaoundé. Mets les verbes entre parenthèses à l'impératif qui convient.
\begin{enumerate}
\item (\esp{beber}, \esp{tú}) al menos dos litros de agua al día.
\item No (\esp{fumar}, \esp{vosotros}) dentro del liceo.
\item (\esp{hacer}, \esp{usted}) ejercicio treinta minutos cada día.
\item No (\esp{tirar}, \esp{ustedes}) basura en la calle.
\item (\esp{venir}, \esp{tú}) a la consulta del médico mañana.
\end{enumerate}
\tcblower
Solution détaillée.
\begin{enumerate}
\item \esp{Bebe al menos dos litros de agua al día.} « Bois au moins deux litres d'eau par jour. » Destinataire \esp{tú}, phrase affirmative : on prend la 3\textsuperscript{e} personne du singulier du présent de \esp{beber} : \esp{bebe}. Verbe régulier.
\item \esp{No fuméis dentro del liceo.} « Ne fumez pas dans le lycée. » Destinataire \esp{vosotros}, phrase \textbf{négative} : subjonctif présent de \esp{fumar} à la 2\textsuperscript{e} personne du pluriel : \esp{fuméis} (accent sur le \esp{é}, radical \esp{fum-} + désinence \esp{-éis}). Erreur à éviter : *\esp{no fumad}.
\item \esp{Haga ejercicio treinta minutos cada día.} « Faites de l'exercice trente minutes par jour. » Destinataire \esp{usted} : subjonctif présent. \esp{Hacer} a un radical irrégulier au subjonctif : \esp{hag-} (comme \esp{hago}), d'où \esp{haga}. Attention : l'impératif \esp{tú} serait \esp{haz}.
\item \esp{No tiren basura en la calle.} « Ne jetez pas d'ordures dans la rue. » Destinataire \esp{ustedes}, négatif : subjonctif présent de \esp{tirar} à la 3\textsuperscript{e} personne du pluriel : \esp{tiren}. Verbe régulier en \esp{-ar} : voyelle thématique \esp{-e-} au subjonctif.
\item \esp{Ven a la consulta del médico mañana.} « Viens à la consultation du médecin demain. » \esp{Venir} fait partie des huit impératifs irréguliers à \esp{tú} : \esp{ven} (et non *\esp{vene}).
\end{enumerate}
Conclusion : la clé est toujours le couple destinataire + signe ; le négatif impose le subjonctif, sans exception.
\end{exoresolu}

\courssub{Savoir-faire 2 : Placer les pronoms compléments avec l'impératif}

\begin{methode}[Enclise et proclise : la règle d'or]
\begin{enumerate}
\item \cle{Observer le signe de la phrase} : c'est lui, et lui seul, qui décide de la place du pronom.
\item \cle{Impératif affirmatif} : le pronom se place \textbf{après} le verbe et s'y \textbf{soude} (enclise) : \esp{levántate, dámelo, cómpralo}. En français, le pronom reste avant : « lève-toi, donne-le-moi ».
\item \cle{Impératif négatif} : le pronom se place \textbf{avant} le verbe, \textbf{détaché} (proclise) : \esp{no te levantes, no me lo des, no lo compres}. Le français fait de même : « ne te lève pas ».
\item \cle{Ordonner les pronoms doubles} : toujours COI avant COD : \esp{me lo, te lo, se lo, nos lo}. À la 3\textsuperscript{e} personne, \esp{le/les} deviennent \esp{se} devant \esp{lo, la, los, las} : \esp{dáselo} (donne-le-lui), jamais *\esp{dálelo}.
\item \cle{Cas particulier \esp{vosotros}} : le \esp{-d} final de l'impératif affirmatif tombe devant \esp{os} : \esp{levantad} + \esp{os} = \esp{levantaos} ; exception : \esp{id} + \esp{os} = \esp{idos}.
\end{enumerate}
Réflexe de vérification : relire la phrase à voix haute ; si l'accent oral tombe au mauvais endroit, l'accent écrit manque.
\end{methode}

\begin{exoresolu}[Thème : impératif et pronoms — Type Bac]
Énoncé. Traduis en espagnol.
\begin{enumerate}
\item Donne-moi la bouteille d'eau.
\item Ne la jette pas par terre !
\item Lavez-vous les mains avant de manger. (à un groupe d'amis)
\item Apportez-les-moi immédiatement. (politesse, pluriel)
\item Ne te lève pas tard le dimanche.
\end{enumerate}
\tcblower
Solution détaillée.
\begin{enumerate}
\item \esp{Dame la botella de agua.} « Donne » = impératif affirmatif de \esp{dar} à \esp{tú} : \esp{da}. Pronom COI \esp{me} soudé après : \esp{dame}. Pas d'accent écrit : \esp{dame} reste paroxytonique terminée par voyelle, accentuée naturellement sur \esp{da}.
\item \esp{¡No la tires por el suelo!} Phrase négative : proclise obligatoire. Le pronom COD \esp{la} (la bouteille, féminin) se détache et précède le subjonctif : \esp{no la tires}. Le calque *\esp{no tírala} est une faute grave : à l'impératif négatif, l'enclise est impossible.
\item \esp{Lavaos las manos antes de comer.} Groupe d'amis = \esp{vosotros}. Verbe pronominal \esp{lavarse} : impératif \esp{lavad} + \esp{os} ; le \esp{-d} tombe : \esp{lavaos}. En espagnol, on utilise l'article défini avec les parties du corps : \esp{las manos}, jamais le possessif (*\esp{vuestras manos}).
\item \esp{Tráiganmelos inmediatamente.} Politesse pluriel = \esp{ustedes} : subjonctif de \esp{traer} : \esp{traigan}. Deux pronoms : COI \esp{me} puis COD \esp{los}, soudés dans cet ordre : \esp{traiganmelos}. Accent écrit obligatoire sur \esp{á} pour conserver l'accent tonique de \esp{traigan} : \esp{tráiganmelos}.
\item \esp{No te levantes tarde el domingo.} Négatif : proclise du pronom réfléchi \esp{te} devant le subjonctif \esp{levantes}. « Le dimanche » se traduit par l'article défini : \esp{el domingo}.
\end{enumerate}
Conclusion : affirmatif = soudure après le verbe avec accent écrit si besoin ; négatif = pronom détaché devant le subjonctif.
\end{exoresolu}

\courssub{Savoir-faire 3 : Maîtriser l'accentuation liée à l'enclise}

\begin{methode}[Placer l'accent écrit quand on soude les pronoms]
\begin{enumerate}
\item \cle{Identifier la syllabe tonique} de l'impératif seul : \esp{DA, com-PRA, es-CRI-be, le-VAN-ta}.
\item \cle{Ajouter le ou les pronoms} : la syllabe tonique \textbf{ne change jamais} de place.
\item \cle{Compter les syllabes} après la soudure : si la syllabe tonique devient l'antépénultième (ou plus), l'accent écrit est \textbf{obligatoire} : \esp{escribe} + \esp{me} = \esp{escríbeme} (tonique sur \esp{es}, 3\textsuperscript{e} syllabe à partir de la fin).
\item \cle{Cas fréquents à mémoriser} :
\begin{center}
\begin{tabularx}{\linewidth}{|X|X|X|}\hline
\rowcolor{popblueL} Impératif seul & Avec pronom(s) & Traduction \\ \hline
\esp{da} & \esp{dame / dámelo} & donne-moi / donne-le-moi \\ \hline
\esp{compra} & \esp{cómpralo} & achète-le \\ \hline
\esp{dile} & \esp{díselo} & dis-le-lui \\ \hline
\esp{haz} & \esp{hazlo} & fais-le \\ \hline
\esp{levanta} & \esp{levántate} & lève-toi \\ \hline
\end{tabularx}
\end{center}
\item \cle{Ne jamais accentuer} quand la soudure ne déplace pas la tonique au-delà de l'avant-dernière syllabe : \esp{dame, hazlo, ponla} s'écrivent sans accent.
\end{enumerate}
\end{methode}

\begin{exoresolu}[Accentuation : à vous de corriger — Type Bac]
Énoncé. Ces phrases extraites de copies d'élèves contiennent chacune une erreur d'accentuation ou de soudure. Corrige-les et justifie.
\begin{enumerate}
\item *\esp{Damelo ahora mismo.}
\item *\esp{Escribeme una carta.}
\item *\esp{Levantaté temprano.}
\item *\esp{Compraló en el mercado de Mvog-Mbi.}
\item *\esp{No te lo dés a tu hermano.}
\end{enumerate}
\tcblower
Solution détaillée.
\begin{enumerate}
\item \esp{Dámelo ahora mismo.} « Donne-le-moi tout de suite. » \esp{da} + \esp{me} + \esp{lo} : la tonique \esp{da} devient antépénultième (mot proparoxytonique) : accent écrit obligatoire sur le \esp{á}.
\item \esp{Escríbeme una carta.} « Écris-moi une lettre. » \esp{escribe} + \esp{me} : tonique sur \esp{es}, antépénultième après soudure : \esp{escríbeme}.
\item \esp{Levántate temprano.} « Lève-toi tôt. » L'accent porte sur \esp{van} (\esp{le-VAN-ta-te}), pas sur \esp{te} : \esp{levántate}. Le pronom réfléchi ne reçoit jamais d'accent.
\item \esp{Cómpralo en el mercado de Mvog-Mbi.} « Achète-le au marché de Mvog-Mbi. » Tonique sur \esp{com}, devenue antépénultième : l'accent va sur le \esp{ó} du radical, jamais sur le pronom \esp{lo}.
\item \esp{No te lo des a tu hermano.} « Ne le lui donne pas. » Phrase négative : subjonctif \esp{des}, monosyllabique, sans accent écrit (contrairement à l'impératif affirmatif \esp{dé}, 1\textsuperscript{re}/3\textsuperscript{e} personne du subjonctif de \esp{dar}, qui en porte un pour se distinguer de la préposition \esp{de}). Ici \esp{des} est la forme de \esp{tú} : pas d'accent.
\end{enumerate}
Conclusion : l'accent écrit sert uniquement à garder la syllabe tonique d'origine ; il ne se met jamais sur le pronom.
\end{exoresolu}

\courssub{Savoir-faire 4 : Traduire les consignes de santé et d'écologie (thème et version)}

\begin{methode}[Traduire une consigne : du français vers l'espagnol et inversement]
\begin{enumerate}
\item \cle{Repérer la structure française} : « il faut + infinitif », « il faut que + subjonctif », impératif direct. Chacune a son équivalent : \esp{hay que + infinitivo} (impersonnel), \esp{es necesario que + subjuntivo}, impératif.
\item \cle{Choisir le registre} : dans une émission de santé, on alterne \esp{hay que} (général) et l'impératif à \esp{ustedes} (direct).
\item \cle{Surveiller les faux amis} du champ lexical :
\begin{center}
\begin{tabularx}{\linewidth}{|X|X|X|}\hline
\rowcolor{popblueL} Français & Español correct & Faux ami / Erreur fréquente \\ \hline
une ordonnance & \esp{una receta} & *\esp{una ordenanza} (ordonnance municipale) \\ \hline
un médicament & \esp{un medicamento} & *\esp{una medicina} (la médecine/science) \\ \hline
la santé & \esp{la salud} & --- \\ \hline
gaspiller & \esp{desperdiciar / malgastar} & *\esp{gastar} (dépenser) \\ \hline
les ordures & \esp{la basura} & *\esp{los desperdicios} (les déchets/résidus) \\ \hline
protéger & \esp{proteger} & *\esp{protejar} (n'existe pas) \\ \hline
\end{tabularx}
\end{center}
\item \cle{Éviter les calques} :
\begin{itemize}
\item « prendre un médicament » = \esp{tomar un medicamento} (pas *\esp{prender}).
\item « faire du sport » = \esp{hacer deporte} (pas *\esp{faire}).
\item « être en bonne santé » = \esp{estar sano / estar en buena salud} (verbe \esp{estar} pour l'état, pas \esp{ser}).
\item « se laver les mains » = \esp{lavarse las manos} (article défini, pas de possessif).
\end{itemize}
\item \cle{Version} : pour traduire vers le français, restituer l'impératif français et les tournures impersonnelles : \esp{hay que reciclar} = « il faut recycler ».
\end{enumerate}
\end{methode}

\begin{exoresolu}[Version et thème — Type Bac]
Énoncé A (version). Traduis en français : \esp{Para cuidar el medio ambiente, hay que separar la basura. No tiren las botellas de plástico al río: recíclenlas. Y para su salud, coman frutas y verduras todos los días y no tomen bebidas azucaradas.}

Énoncé B (thème). Traduis en espagnol : « Pour rester en bonne santé, il faut dormir huit heures. Faites du sport régulièrement et ne gaspillez pas l'eau : fermez le robinet quand vous vous brossez les dents. »
\tcblower
Solution détaillée.

A. Version : « Pour préserver l'environnement, il faut trier les ordures. Ne jetez pas les bouteilles en plastique dans la rivière : recyclez-les. Et pour votre santé, mangez des fruits et des légumes tous les jours et ne buvez pas de boissons sucrées. »
Justifications : \esp{hay que separar} = tournure impersonnelle « il faut trier » ; \esp{no tiren} = impératif négatif à \esp{ustedes} (« ne jetez pas ») ; \esp{recíclenlas} = impératif affirmatif avec enclise du COD \esp{las} (les bouteilles) : « recyclez-les » ; \esp{bebidas azucaradas} = « boissons sucrées ».

B. Thème : \esp{Para estar sano, hay que dormir ocho horas. Hagan deporte regularmente y no desperdicien el agua: cierren el grifo cuando se cepillen los dientes.}
Justifications :
\begin{itemize}
\item « rester en bonne santé » : \esp{estar sano} (\esp{estar} pour l'état, pas \esp{ser}).
\item « il faut dormir » : \esp{hay que dormir}, impersonnel, sans sujet.
\item « faites du sport » : destinataire collectif, registre de politesse du podcast : \esp{hagan deporte} (subjonctif de \esp{hacer}, radical \esp{hag-}).
\item « ne gaspillez pas » : négatif, subjonctif : \esp{no desperdicien}.
\item « fermez le robinet » : \esp{cierren el grifo} ; \esp{cerrar} est à diphtongue (\esp{cierr-}).
\item « quand vous vous brossez les dents » : après \esp{cuando} pour une action future/habituelle dans une consigne générale : \esp{cuando se cepillen los dientes} (subjonctif, valeur générale) ; parties du corps avec article défini : \esp{los dientes}.
\end{itemize}
Conclusion : une consigne bien traduite respecte trois choix : la tournure (impérative ou impersonnelle), la personne (\esp{ustedes} pour un public), et le lexique exact sans faux amis.
\end{exoresolu}

\courssub{Savoir-faire 5 : Réaliser un court podcast santé ou écologique}

\begin{methode}[Structurer une émission orale de conseils en 5 étapes]
\begin{enumerate}
\item \cle{Saluer et présenter} l'émission : \esp{Buenos días, queridos oyentes. Bienvenidos a «Salud y planeta», su programa semanal desde Yaundé.}
\item \cle{Annoncer le thème} : \esp{Hoy hablamos de cómo proteger nuestra salud y nuestro barrio.}
\item \cle{Enchaîner les conseils} avec des impératifs et des tournures impersonnelles variées, reliés par des connecteurs : \esp{primero, además, también, por último}. Alterner \esp{hay que...}, \esp{no olviden...}, \esp{beban..., reciclen...}.
\item \cle{Soigner la langue} : impératifs corrects, pronoms bien placés (\esp{consúmelas, no las tires}), lexique précis (\esp{el médico, la receta, la basura, el medio ambiente, reciclar, ahorrar agua}).
\item \cle{Conclure et remercier} : \esp{Esto es todo por hoy. Cuídense mucho y hasta la próxima semana.}
\end{enumerate}
Réflexe : préparer une fiche de 6 à 8 consignes maximum, répéter à voix haute en surveillant l'accentuation des formes avec pronoms.
\end{methode}

\begin{exoresolu}[Production : rédige le script d'un podcast — Type Bac]
Énoncé. Tu es l'animateur de l'émission « Salud y planeta » diffusée au club d'espagnol de ton lycée à Douala. Rédige le script complet (10 à 12 lignes) d'un podcast donnant des conseils de santé et d'écologie aux élèves. Utilise au moins quatre impératifs différents, deux impératifs négatifs et deux formes avec pronoms enclitiques.
\tcblower
Solution détaillée — script modèle :

\esp{Buenos días, queridos oyentes. Bienvenidos a «Salud y planeta», el programa del club de español del liceo de Duala. Hoy les damos consejos para vivir mejor. Primero, cuiden su cuerpo: beban mucha agua, coman frutas frescas del mercado y hagan deporte cada día. No pasen toda la tarde delante del teléfono: levántense y caminen un poco. Segundo, protejan nuestro barrio: no tiren basura en la calle; recojan las botellas de plástico y recíclenlas. Si ven un grifo abierto, ciérrenlo: el agua es un tesoro. Por último, si están enfermos, vayan al médico y sigan su receta. Esto es todo por hoy. Cuídense mucho y hasta la próxima semana.}

Traduction : « Bonjour, chers auditeurs. Bienvenue à "Santé et planète", l'émission du club d'espagnol du lycée de Douala. Aujourd'hui, nous vous donnons des conseils pour mieux vivre. D'abord, prenez soin de votre corps : buvez beaucoup d'eau, mangez des fruits frais du marché et faites du sport chaque jour. Ne passez pas tout l'après-midi devant le téléphone : levez-vous et marchez un peu. Ensuite, protégez notre quartier : ne jetez pas d'ordures dans la rue ; ramassez les bouteilles en plastique et recyclez-les. Si vous voyez un robinet ouvert, fermez-le : l'eau est un trésor. Enfin, si vous êtes malades, allez chez le médecin et suivez son ordonnance. C'est tout pour aujourd'hui. Prenez bien soin de vous et à la semaine prochaine. »

Justifications grammaticales :
\begin{itemize}
\item Impératifs affirmatifs à \esp{ustedes} (public) : \esp{beban, coman, hagan, caminen, protejan, recojan, vayan, sigan} — tous au subjonctif présent, y compris les irréguliers \esp{hagan} (hacer), \esp{vayan} (ir), \esp{sigan} (seguir, radical \esp{sig-}).
\item Impératifs négatifs : \esp{no pasen, no tiren} — subjonctif précédé de \esp{no}.
\item Pronoms enclitiques avec accent écrit : \esp{levántense} (tonique \esp{van} conservée), \esp{recíclenlas} (tonique \esp{ci} + COD \esp{las}), \esp{ciérrenlo} (tonique \esp{cierr} + COD \esp{lo}), \esp{cuídense} (tonique \esp{cui}).
\item Connecteurs : \esp{primero, segundo, por último} ; conclusion rituelle \esp{cuídense}.
\end{itemize}
Conclusion : le script respecte le cahier des charges (4 impératifs, 2 négatifs, 2 enclises) tout en restant naturel et adapté à l'oral.
\end{exoresolu}

\courssub{Entraînement : Thème et Version guidés}

\begin{exercice}[Entraînement Bac — Thème / Version]
\textbf{A. Version.} Traduisez en français :
\esp{Estimados estudiantes: para evitar el dengue, no dejen agua estancada en los neumáticos viejos. Tapen los tanques de agua y usen mosquiteros en las ventanas. Si tienen fiebre, no se automediquen: vayan al centro de salud. Cuídense.}

\textbf{B. Thème.} Traduisez en espagnol :
« Chers amis, pour protéger la planète, il faut réduire les déchets. Ne jetez pas les piles à la poubelle normale : déposez-les dans les conteneurs spéciaux. Économisez l'eau et l'électricité. Éteignez les lumières en sortant. Faites-le pour votre avenir ! »
\end{exercice}

\begin{corrige}[Exercice Entraînement Bac]
\textbf{A. Version :} « Chers étudiants : pour éviter la dengue, ne laissez pas d'eau stagnante dans les vieux pneus. Couvrez les réservoirs d'eau et utilisez des moustiquaires aux fenêtres. Si vous avez de la fièvre, ne vous soignez pas vous-mêmes : allez au centre de santé. Prenez soin de vous. »
\begin{itemize}
\item \esp{no dejen} (impératif négatif \esp{ustedes} de \esp{dejar}) ; \esp{tapen} (impératif affirmatif \esp{ustedes} de \esp{tapar}) ; \esp{usen} (impératif \esp{ustedes} de \esp{usar}) ; \esp{no se automediquen} (négatif pronominal \esp{ustedes}) ; \esp{vayan} (impératif \esp{ustedes} irrégulier de \esp{ir}) ; \esp{cuídense} (impératif affirmatif pronominal \esp{ustedes} de \esp{cuidarse}, accent sur \esp{uí}).
\end{itemize}

\textbf{B. Thème :} \esp{Queridos amigos: para proteger el planeta, hay que reducir los residuos. No tiren las pilas a la basura normal: deposítenlas en los contenedores especiales. Ahorren agua y electricidad. Apaguen las luces al salir. ¡Háganlo por su futuro!}
\begin{itemize}
\item \esp{hay que reducir} (impersonnel) ; \esp{no tiren} (négatif \esp{ustedes}) ; \esp{deposítenlas} (affirmatif \esp{ustedes} \esp{depositen} + \esp{las} $\rightarrow$ accent sur \esp{é}) ; \esp{ahorren, apaguen} (affirmatif \esp{ustedes}) ; \esp{háganlo} (affirmatif \esp{ustedes} irrégulier \esp{hagan} + \esp{lo} $\rightarrow$ accent sur \esp{á}).
\end{itemize}
\end{corrige}

\courssub{Activité orale : Enregistrement du podcast}

\begin{experience}[Enregistrement « Salud y planeta » — Travail en binôme / groupe]
\textbf{Consigne :} En binôme, préparez et enregistrez (sur téléphone ou ordinateur) un épisode de 1 minute 30 de « Salud y planeta ».
\begin{enumerate}
\item \textbf{Rôles :} Un animateur (présente, conclut), un expert invité (donne 3 conseils précis).
\item \textbf{Contraintes linguistiques obligatoires :}
\begin{itemize}
\item Minimum 5 impératifs à \esp{ustedes} (affirmatifs).
\item Minimum 2 impératifs négatifs.
\item Minimum 2 formes avec pronoms enclitiques (dont 1 double pronom COI+COD type \esp{dáselo, recíclenlas}).
\item Utiliser \esp{hay que} au moins une fois.
\item Lexique : \esp{reciclar, ahorrar, basura, grifo, receta, medicamento, ejercicio, frutas/verduras}.
\end{itemize}
\item \textbf{Scénario suggéré :} Spécial « Rentrée scolaire au Cameroun » : chaleur, moustiques, déchets plastiques, alimentation saine (mangues, papayes, eau).
\item \textbf{Évaluation par les pairs (grille) :} Prononciation / Accentuation des enclises / Correction grammaticale / Fluidité / Respect du cahier des charges.
\end{enumerate}
\end{experience}

\courssub{Culture et ouverture : Le Cameroun et la Guinée équatoriale}

\begin{savaistu}[Santé, environnement et francophonie/hispanophonie]
La Guinée équatoriale (Malabo, Bata) est le seul pays africain où l'espagnol est langue officielle. Elle partage avec le Cameroun des défis sanitaires et écologiques communs : gestion des déchets plastiques à Douala et Bata, lutte contre le paludisme/dengue, accès à l'eau potable.
\begin{itemize}
\item \textbf{Vocabulaire partagé :} \esp{el paludismo} (paludisme), \esp{el dengue}, \esp{la red mosquitera} (moustiquaire), \esp{el agua potable}, \esp{la basura plástica}.
\item \textbf{Initiatives réelles :} Campagnes « \esp{Limpia tu barrio} » (Nettoie ton quartier) dans les lycées de Malabo et Yaoundé ; programmes \esp{Ecoescuelas} (Écoles écologiques) soutenus par l'UNESCO.
\item \textbf{Expression utile :} \esp{« La salud no tiene fronteras »} (La santé n'a pas de frontières) — devise souvent utilisée dans les campagnes transfrontalières CEMAC.
\end{itemize}
\end{savaistu}

\begin{attention}[Les 5 erreurs qui coûtent des points au Bac]
\begin{enumerate}
\item \cle{Enclise au négatif} : écrire *\esp{no tíralo} au lieu de \esp{no lo tires}. À l'impératif négatif, le pronom est toujours \textbf{avant} le verbe et détaché. C'est la faute la plus sanctionnée dans cette leçon.
\item \cle{Accent écrit oublié ou mal placé} : *\esp{escribeme}, *\esp{levantaté}, *\esp{compraló}. L'accent se met sur la syllabe tonique d'origine du verbe (\esp{escríbeme, levántate, cómpralo}), jamais sur le pronom.
\item \cle{Impératif négatif sans subjonctif} : *\esp{no fumas}, *\esp{no fumad}. Le négatif exige le subjonctif : \esp{no fumes, no fuméis}.
\item \cle{Faux amis et calques du français} : « ordonnance » = \esp{receta} (et non *\esp{ordenanza}) ; « se laver les mains » = \esp{lavarse las manos} (article défini, pas de possessif) ; « être en bonne santé » = \esp{estar sano} (verbe \esp{estar}, pas \esp{ser}).
\item \cle{Confusion des personnes} : mélanger \esp{tú} et \esp{ustedes} dans le même texte, ou utiliser *\esp{hace} comme impératif de politesse au lieu de \esp{haga}. Choisir un destinataire au début du podcast et s'y tenir jusqu'au bout.
\end{enumerate}
\end{attention}

\newpage

\coursec{Exercices}

\courssub{Je vérifie mes connaissances}

\begin{exercice}[QCM : formation et place des pronoms à l'impératif]
Pour chaque question, une seule réponse est correcte. Cochez la bonne case.
\begin{enumerate}
\item Quelle est la forme correcte de l'impératif affirmatif pour \emph{usted} du verbe \esp{lavar} avec les pronoms \emph{me} et \emph{lo} (me laver les mains) ?
\begin{itemize}
\item[] A. Lávamelo
\item[] B. Lávemelo
\item[] C. \textbf{Lávemelo} (avec accent sur la voyelle tonique du verbe)
\item[] D. Lávamelo
\end{itemize}
\item À l'impératif négatif \emph{tú}, où se placent les pronoms compléments ?
\begin{itemize}
\item[] A. Après le verbe (enclise)
\item[] B. Avant le verbe (proclise)
\item[] C. Entre \emph{no} et le verbe
\item[] D. À la fin de la phrase
\end{itemize}
\item La forme \esp{no se lo digas} correspond à :
\begin{itemize}
\item[] A. Impératif affirmatif \emph{tú}
\item[] B. Impératif négatif \emph{tú}
\item[] C. Impératif affirmatif \emph{usted}
\item[] D. Impératif négatif \emph{usted}
\end{itemize}
\item Quel verbe suivant garde l'accent écrit à l'impératif affirmatif \emph{tú} avec un pronom enclitique ?
\begin{itemize}
\item[] A. \esp{hablar} → \esp{háblame}
\item[] B. \esp{comer} → \esp{cómele}
\item[] C. \esp{subir} → \esp{súbelo}
\item[] D. \esp{lavar} → \esp{lávalo}
\end{itemize}
\end{enumerate}
\end{exercice}

\begin{exercice}[Vrai / Faux justifié]
Indiquez si l'affirmation est vraie (V) ou fausse (F) et justifiez brièvement en français.
\begin{enumerate}
\item L'impératif affirmatif \emph{vosotros} de \esp{dar} + \emph{me} + \emph{lo} s'écrit \esp{dádmelo}.
\item À l'impératif négatif, l'ordre des pronoms est : \emph{no} + pronom(s) + verbe.
\item L'accentuation de l'impératif affirmatif avec pronoms suit toujours la règle générale des mots \emph{llanos} terminés par voyelle, \emph{n} ou \emph{s}.
\item En français, les pronoms compléments se placent après le verbe à l'impératif affirmatif (ex. \emph{Donne-le-moi}).
\item La forme \esp{levántense} est l'impératif affirmatif \emph{ustedes} du verbe pronominal \esp{levantarse}.
\end{enumerate}
\end{exercice}

\begin{exercice}[Vocabulaire : santé et environnement]
Associez chaque terme espagnol à sa définition française. Écrivez la lettre correspondante.
\begin{center}
\begin{tabularx}{\linewidth}{|X|X|X|}
\hline
\rowcolor{popblueL} \textbf{Espagnol} & \textbf{Français} & \textbf{Lettre} \\
\hline
1. \esp{hervir el agua} & A. Se protéger des moustiques & \\
2. \esp{tirar la basura} & B. Faire bouillir l'eau & \\
3. \esp{protegerse de los mosquitos} & C. Jeter les ordures & \\
4. \esp{usar mosquiteras} & D. Utiliser des moustiquaires & \\
5. \esp{cuidar el barrio} & E. Prendre soin du quartier & \\
6. \esp{lavarse las manos} & F. Se laver les mains & \\
\hline
\end{tabularx}
\end{center}
\end{exercice}

\begin{exercice}[Conjugaison à compléter : impératif avec pronoms]
Conjuguez les verbes entre parenthèses à l'impératif affirmatif ou négatif selon le contexte, en plaçant correctement les pronoms indiqués. Écrivez la forme complète.
\begin{enumerate}
\item (Tú, afirmativo) \esp{\_\_\_\_\_\_\_\_\_\_} (dar + me + lo) el libro, por favor.
\item (Usted, negativo) \esp{\_\_\_\_\_\_\_\_\_\_} (no + tirar + se + la) la basura al río.
\item (Vosotros, afirmativo) \esp{\_\_\_\_\_\_\_\_\_\_} (lavar + os + las) las manos antes de comer.
\item (Ustedes, afirmativo) \esp{\_\_\_\_\_\_\_\_\_\_} (proteger + se + los) de los mosquitos.
\item (Tú, negativo) \esp{\_\_\_\_\_\_\_\_\_\_} (no + dar + se + la) a ellos, es peligrosa.
\end{enumerate}
\end{exercice}

\courssub{Je m'entraîne}

\begin{exercice}[Phrases à traduire : thème (français → espagnol)]
Traduisez en espagnol en utilisant l'impératif approprié (personne indiquée) et en plaçant correctement les pronoms.
\begin{enumerate}
\item (Tú) Lave-toi les mains avant de manger.
\item (Usted) Ne jetez pas les ordures dans la rivière.
\item (Tú) Donne-la-moi, s'il te plaît (la bouteille).
\item (Ustedes) Ne la leur donnez pas.
\item (Vosotros) Levez-vous et prenez soin de votre quartier.
\end{enumerate}
\end{exercice}

\begin{exercice}[Transformation : indicatif → impératif affirmatif et négatif]
Transformez chaque phrase de l'indicatif à l'impératif affirmatif \emph{tu} puis à l'impératif négatif \emph{tu}, en conservant les pronoms compléments.
\begin{enumerate}
\item \esp{Tú me das el libro.}
\item \esp{Tú te lavas las manos.}
\item \esp{Tú nos tiras la pelota.}
\item \esp{Tú se lo dices a ella.}
\item \esp{Tú os levantáis temprano.}
\end{enumerate}
\end{exercice}

\begin{exercice}[Texte à trous : script de podcast écologique]
Complétez le script ci-dessous avec la forme correcte de l'impératif (affirmatif ou négatif) du verbe indiqué, y compris les pronoms nécessaires.
\begin{textebox}[Script podcast « Cuidemos nuestro barrio »]
¡Hola, vecinos! Bienvenidos a nuestro podcast ecológico. Hoy les traigo cinco consejos para un barrio más limpio y sano.
1. \_\_\_\_\_\_\_\_\_\_ (hervir + lo) el agua antes de beberla.
2. \_\_\_\_\_\_\_\_\_\_ (no + tirar + la) la basura en la calle.
3. \_\_\_\_\_\_\_\_\_\_ (usar + las) bolsas reutilizables para la compra.
4. \_\_\_\_\_\_\_\_\_\_ (proteger + se) de los mosquitos: \_\_\_\_\_\_\_\_\_\_ (poner + las) mosquiteras en las ventanas.
5. \_\_\_\_\_\_\_\_\_\_ (cuidar + lo) el parque: \_\_\_\_\_\_\_\_\_\_ (no + ensuciar + lo).
¡Gracias por escuchar! \_\_\_\_\_\_\_\_\_\_ (compartir + lo) con tus amigos.
\end{textebox}
\emph{Glossaire :} \esp{vecinos} = voisins ; \esp{bolsas reutilizables} = sacs réutilisables ; \esp{ensuciar} = salir.
\end{exercice}

\begin{exercice}[Appariement : formes affirmatives et négatives]
Reliez chaque impératif affirmatif (colonne A) à sa forme négative correspondante (colonne B) pour la même personne et les mêmes pronoms.
\begin{center}
\begin{tabularx}{\linewidth}{|X|X|}
\hline
\rowcolor{popgreenL} \textbf{A. Affirmatif} & \textbf{B. Négatif} \\
\hline
1. \esp{Dámelo} & a. \esp{No me lo des} \\
2. \esp{Levántense} & b. \esp{No se lo digas} \\
3. \esp{Dígaselo} & c. \esp{No se levanten} \\
4. \esp{Cuidemosla} & d. \esp{No la cuidemos} \\
5. \esp{Póngansela} & e. \esp{No se la pongan} \\
\hline
\end{tabularx}
\end{center}
\end{exercice}

\courssub{Je me prépare au Bac}

\begin{exercice}[Compréhension et langue : texte de campagne sanitaire]
Lisez le texte suivant et répondez aux questions en français.
\begin{textebox}[Campaña « Agua segura, vida sana » — Ministerio de Salud]
¡Atención, comunidad! Para prevenir enfermedades, sigan estas normas:
1. \esp{Hiervan el agua y bébanla.}
2. \esp{No tiren la basura al río.}
3. \esp{Protéjanse de los mosquitos: usen mosquiteras y cuídense.}
4. \esp{Lávense las manos antes de cocinar y después de usar el baño.}
5. \esp{No compartan utensilios personales.}
6. \esp{Acudan al centro de salud si tienen fiebre o diarrea.}
¡Su salud está en sus manos!
\end{textebox}
\emph{Glossaire :} \esp{utensilios personales} = ustensiles personnels ; \esp{acudan} = rendez-vous (impératif \emph{ustedes} de \esp{acudir}).

\textbf{Questions :}
\begin{enumerate}
\item Relevez les six verbes à l'impératif du texte et indiquez pour chacun : la personne (\emph{tú, usted, ustedes, vosotros, nosotros}) et la polarité (affirmatif / négatif).
\item Expliquez, exemples du texte à l'appui, la place des pronoms compléments (réfléchis, COD, COI) à l'impératif affirmatif et à l'impératif négatif.
\item Pourquoi les formes \esp{bébanla}, \esp{cuídense}, \esp{lávense} portent-elles un accent écrit ? Donnez la règle générale.
\item Traduisez en français les consignes 1, 3 et 4.
\end{enumerate}
\end{exercice}

\begin{exercice}[Thème et version : impératif avec pronoms]
\textbf{A. Thème (français → espagnol)} — Traduisez en espagnol les phrases suivantes. Utilisez l'impératif à la personne indiquée entre parenthèses.
\begin{enumerate}
\item (Tú) Lave-toi les mains avant de manger.
\item (Usted) Ne jetez pas les ordures dans la rivière.
\item (Tú) Donne-la-moi, s'il te plaît (la bouteille).
\item (Ustedes) Ne la leur donnez pas.
\item (Vosotros) Levez-vous et prenez soin de votre quartier.
\end{enumerate}

\textbf{B. Version (espagnol → français)} — Traduisez en français le paragraphe suivant extrait d'une campagne de santé publique.
\begin{textebox}[Extrait de campagne]
\esp{Hiervan el agua y bébanla. No tiren la basura al río. Protéjanse de los mosquitos: usen mosquiteras y cuídense. Lávense las manos frecuentemente. No compartan objetos personales. Acudan al médico ante cualquier síntoma.}
\end{textebox}
\end{exercice}

\begin{exercice}[Rédaction guidée : consignes écologiques pour le lycée]
Rédigez en espagnol \textbf{5 consignes écologiques} destinées à être affichées sur le tableau d'affichage du lycée. 
\begin{itemize}
\item Chaque consigne doit être à l'impératif (choisissez la personne appropriée : \emph{tú, usted, ustedes, nosotros}).
\item Au moins \textbf{3 consignes} doivent contenir un ou plusieurs pronoms compléments (COD, COI, réfléchi) correctement placés (enclise à l'affirmatif, proclise au négatif).
\item Respectez l'accentuation écrite liée à l'enclise.
\item Nombre de mots total : \textbf{entre 60 et 80 mots}.
\item Utilisez du vocabulaire du thème : \esp{reciclar, ahorrar agua, apagar luces, usar transporte público, plantar árboles, cuidar el medio ambiente, reducir residuos, reutilizar, separar la basura, compostar}.
\end{itemize}
\textit{Conseil :} Commencez par lister vos 5 idées en français, puis traduisez-les en vérifiant la place des pronoms et les accents.
\end{exercice}

\newpage

\coursec{Corrigés des exercices}

\begin{corrige}[Exercice 1 — QCM : formation et place des pronoms à l'impératif]
\textbf{Réponses :}
\begin{enumerate}
\item \textbf{Bonne réponse : B/C — \esp{Lávemelo}.} L'impératif \emph{usted} de \esp{lavar} est \esp{lave} (forme de subjonctif présent). Avec les pronoms enclitiques \emph{me} + \emph{lo} : \esp{lave} + \esp{me} + \esp{lo} = \esp{lávemelo}. L'accent écrit est obligatoire : la syllabe tonique reste \emph{la}- (première syllabe), et le mot devient \emph{esdrújulo} (accentué sur l'antépénultième), toujours accentué. \esp{Lávamelo} est doublement faux : mauvaise base verbale (\esp{lava} = \emph{tú}) et accentuation incorrecte.
\item \textbf{Bonne réponse : B.} À l'impératif négatif, les pronoms se placent \emph{avant} le verbe (proclise), entre \esp{no} et le verbe : \esp{no me lo des}. La réponse C décrit d'ailleurs la même réalité : les pronoms se placent bien entre \esp{no} et le verbe ; \textbf{B et C sont donc toutes deux acceptables}, la formulation la plus complète étant C.
\item \textbf{Bonne réponse : B.} \esp{No se lo digas} : présence de \esp{no} + proclise + forme \esp{digas} (subjonctif présent, 2\textsuperscript{e} personne du singulier) = impératif négatif \emph{tú} de \esp{decir}.
\item \textbf{Bonne réponse : A — \esp{háblame}.} Règle : à l'impératif affirmatif avec enclise, la syllabe tonique du verbe ne bouge pas ; si le mot résultant devient \emph{esdrújulo} ou contredit la règle des \emph{llanas}, on écrit un accent. \textbf{Attention : les quatre formes proposées ne portent pas toutes un accent.} \esp{háblame}, \esp{súbelo}, \esp{lávalo} sont esdrújules (accent obligatoires) ; \esp{comele} (\esp{comer} + \emph{le}) est une \emph{llana} terminée par voyelle (tonique sur \emph{me}), elle \textbf{ne s'accentue pas}. La forme \esp{cómele} est donc erronée. La réponse attendue est A, mais retenez la règle générale plutôt qu'un cas isolé.
\end{enumerate}
\textbf{Points de vigilance :} ne jamais écrire \esp{lavemelo} sans accent ; ne jamais dire \esp{no dámelo} (mélange interdit de négation et d'enclise).
\end{corrige}

\begin{corrige}[Exercice 2 — Vrai / Faux justifié]
\begin{enumerate}
\item \textbf{FAUX.} L'impératif \emph{vosotros} de \esp{dar} est \esp{dad} ; avec \emph{me} + \emph{lo} : \esp{dad} + \esp{me} + \esp{lo} = \esp{dadmelo}, \textbf{sans accent écrit} : la syllabe tonique (\emph{me}) est l'avant-dernière et le mot se termine par une voyelle (mot \emph{llano}). La forme correcte est donc \esp{dadmelo}.
\item \textbf{VRAI.} À l'impératif négatif, la structure est : \esp{no} + pronom(s) + verbe au subjonctif présent : \esp{No me lo digas} (« Ne me le dis pas »).
\item \textbf{FAUX.} C'est l'inverse : l'ajout de pronoms enclitiques déplace la position relative de la syllabe tonique ; le mot devient souvent \emph{esdrújulo} ou \emph{sobreesdrújulo}, et l'accent écrit devient obligatoire : \esp{dámelo}, \esp{dígaselo}, \esp{levántense}.
\item \textbf{VRAI.} En français aussi, les pronoms suivent le verbe à l'impératif affirmatif (\emph{Donne-le-moi}), mais l'ordre interne diffère : en espagnol COI avant COD (\esp{dámelo} = me + lo), en français « le » avant « moi » dans \emph{donne-le-moi}. Attention au calque.
\item \textbf{VRAI.} \esp{levantarse} : impératif \emph{ustedes} = \esp{levanten} ; avec le pronom réfléchi enclitique \esp{se} : \esp{levanten} + \esp{se} = \esp{levántense}, avec accent car le mot devient esdrújulo (\emph{le}-\textbf{\emph{van}}-\emph{ten}-\emph{se}).
\end{enumerate}
\end{corrige}

\begin{corrige}[Exercice 3 — Vocabulaire : santé et environnement]
\textbf{Correspondances :}
\begin{center}
\begin{tabularx}{\linewidth}{|X|X|X|}
\hline
\rowcolor{popblueL} \textbf{Espagnol} & \textbf{Français} & \textbf{Lettre} \\
\hline
1. \esp{hervir el agua} & Faire bouillir l'eau & \textbf{B} \\
2. \esp{tirar la basura} & Jeter les ordures & \textbf{C} \\
3. \esp{protegerse de los mosquitos} & Se protéger des moustiques & \textbf{A} \\
4. \esp{usar mosquiteras} & Utiliser des moustiquaires & \textbf{D} \\
5. \esp{cuidar el barrio} & Prendre soin du quartier & \textbf{E} \\
6. \esp{lavarse las manos} & Se laver les mains & \textbf{F} \\
\hline
\end{tabularx}
\end{center}
\textbf{Points de vigilance :} \esp{hervir} est un verbe à changement vocalique (e → ie) : \esp{hierve}, impératif \esp{hierve} / \esp{hierva} / \esp{hiervan}. \esp{mosquitera} = moustiquaire (faux ami possible avec « moustique » = \esp{mosquito}).
\end{corrige}

\begin{corrige}[Exercice 4 — Conjugaison à compléter : impératif avec pronoms]
\begin{enumerate}
\item \esp{\textbf{Dámelo}, por favor.} — Impératif \emph{tú} de \esp{dar} : \esp{da} ; enclise de \emph{me} + \emph{lo} ; accent obligatoire (mot esdrújulo). « Donne-le-moi, s'il te plaît. »
\item \esp{\textbf{No se la tire} al río.} — Impératif négatif \emph{usted} : \esp{no} + proclise (\esp{se la}) + \esp{tire} (subjonctif présent de \esp{tirar}). « Ne la jetez pas dans la rivière. »
\item \esp{\textbf{Lavaóslas} antes de comer.} — Impératif \emph{vosotros} : \esp{lavad} → \esp{lava} + \esp{os} = \esp{lavaós} (accent sur le \emph{o} : hiatus \emph{ao}, mot llano en \emph{s}) ; + \esp{las} = \esp{lavaóslas}, avec accent (mot esdrújulo : \emph{la}-\emph{va}-\textbf{\emph{os}}-\emph{las}). « Lavez-vous-les avant de manger. »
\item \esp{\textbf{Protéjanse} de los mosquitos.} — Impératif \emph{ustedes} de \esp{protegerse} : \esp{protejan} + \esp{se} = \esp{protéjanse} (accent, esdrújulo). « Protégez-vous des moustiques. »
\item \esp{\textbf{No se la des} a ellos, es peligrosa.} — Impératif négatif \emph{tú} : \esp{no} + \esp{se la} + \esp{des} (subjonctif présent de \esp{dar}). « Ne la leur donne pas, elle est dangereuse. »
\end{enumerate}
\textbf{Points de vigilance :} l'ordre des pronoms est toujours COI (\esp{se/me/os}) avant COD (\esp{lo/la/las}) ; \esp{le/les} devient \esp{se} devant \esp{lo/la/los/las}.
\end{corrige}

\begin{corrige}[Exercice 5 — Phrases à traduire : thème]
\begin{enumerate}
\item « Lave-toi les mains avant de manger. » → \esp{\textbf{Lávate las manos antes de comer.}} — Impératif \emph{tú} de \esp{lavarse} : \esp{lava} + \esp{te} = \esp{lávate} (accent, esdrújulo).
\item « Ne jetez pas les ordures dans la rivière. » → \esp{\textbf{No tire la basura al río.}} — Négatif \emph{usted} : \esp{no} + \esp{tire} (subjonctif). On peut aussi écrire \esp{No la tire al río} si « les ordures » est repris par pronom.
\item « Donne-la-moi, s'il te plaît. » → \esp{\textbf{Dámela, por favor.}} — \esp{da} + \esp{me} + \esp{la} (la bouteille = \esp{la botella}, féminin) = \esp{dámela}. Attention à l'ordre : en espagnol COI (\esp{me}) avant COD (\esp{la}), contrairement au français « donne-la-moi ».
\item « Ne la leur donnez pas. » → \esp{\textbf{No se la den.}} — Négatif \emph{ustedes} : \esp{no} + \esp{se la} (\esp{les} → \esp{se} devant \esp{la}) + \esp{den}.
\item « Levez-vous et prenez soin de votre quartier. » → \esp{\textbf{Levántaos y cuidad vuestro barrio.}} — \esp{levantad} → \esp{levanta} + \esp{os} = \esp{levántaos} (accent obligatoire : mot esdrújulo \emph{le}-\textbf{\emph{van}}-\emph{ta}-\emph{os}) ; \esp{cuidad} = impératif \emph{vosotros} de \esp{cuidar} (sans accent, llana terminée en \emph{d}).
\end{enumerate}
\textbf{Points de vigilance :} vérifier le genre du COD (\esp{la botella} → \esp{la}) ; ne pas calquer l'ordre français des pronoms ; accentuer \esp{lávate}, \esp{dámela}, \esp{levántaos}.
\end{corrige}

\begin{corrige}[Exercice 6 — Transformation : indicatif → impératif]
\begin{enumerate}
\item \esp{Tú me das el libro.} → Affirmatif : \esp{\textbf{Dame el libro.}} (ou \esp{Dámelo.}) — Négatif : \esp{\textbf{No me des el libro.}} (ou \esp{No me lo des.})
\item \esp{Tú te lavas las manos.} → Affirmatif : \esp{\textbf{Lávate las manos.}} — Négatif : \esp{\textbf{No te laves las manos.}}
\item \esp{Tú nos tiras la pelota.} → Affirmatif : \esp{\textbf{Tíranos la pelota.}} — Négatif : \esp{\textbf{No nos tires la pelota.}} (\esp{tira} + \esp{nos} = \esp{tíranos}, accent obligatoire.)
\item \esp{Tú se lo dices a ella.} → Affirmatif : \esp{\textbf{Díselo a ella.}} — Négatif : \esp{\textbf{No se lo digas a ella.}} (\esp{di} + \esp{se} + \esp{lo} = \esp{díselo}, accent ; forme irrégulière \esp{di} de \esp{decir}.)
\item \esp{Os levantáis temprano.} → Affirmatif : \esp{\textbf{Levántaos temprano.}} — Négatif : \esp{\textbf{No os levantéis temprano.}} (La phrase de départ mélangeait \esp{tú} et \esp{vosotros} ; on corrige en traitant le verbe à la personne \emph{vosotros}. \esp{Levántaos} porte l'accent car esdrújulo.)
\end{enumerate}
\textbf{Règle récapitulative :} affirmatif = impératif + pronoms soudés après (enclise) + accent si nécessaire ; négatif = \esp{no} + pronoms détachés avant + subjonctif présent.
\end{corrige}

\begin{corrige}[Exercice 7 — Texte à trous : script de podcast écologique]
\textbf{Script complété :}
\begin{enumerate}
\item \esp{\textbf{Hiérvelo} el agua antes de beberla.} — Impératif \emph{tú} de \esp{hervir} (e → ie) : \esp{hierve} + \esp{lo} = \esp{hiérvelo} (accent, esdrújulo). Remarque : on dira plus naturellement \esp{Hierve el agua}, le pronom faisant double emploi avec le nom ; la forme avec enclise seule serait \esp{Hiérvela} (l'eau, féminin). Acceptez \esp{Hiérvelo} comme forme mécanique attendue.
\item \esp{\textbf{No la tires} en la calle.} — Négatif \emph{tú} : \esp{no} + \esp{la} + \esp{tires}.
\item \esp{\textbf{Úsalas} para la compra.} — \esp{usa} + \esp{las} = \esp{úsalas} (accent, esdrújulo).
\item \esp{\textbf{Protégete} de los mosquitos: \textbf{ponlas} en las ventanas.} — \esp{protege} + \esp{te} = \esp{protégete} (accent, esdrújulo) ; \esp{pon} + \esp{las} = \esp{ponlas} (sans accent : mot \textbf{aguda} \emph{pon}-\textbf{\emph{las}}, tonique sur la dernière syllabe, terminée en \emph{s} → pas d'accent).
\item \esp{\textbf{Cuídalo}: \textbf{no lo ensucies}.} — \esp{cuida} + \esp{lo} = \esp{cuídalo} (accent, esdrújulo) ; négatif : \esp{no lo ensucies}.
\end{enumerate}
Conclusion : \esp{\textbf{Compártelo} con tus amigos.} — \esp{comparte} + \esp{lo} = \esp{compártelo} (accent, esdrújulo).
\textbf{Points de vigilance :} \esp{ponlas} (aguda en \emph{-s}) et \esp{dadmelo} (llana en voyelle) ne prennent pas d'accent ; \esp{hiérvelo}, \esp{úsalas}, \esp{cuídalo}, \esp{compártelo}, \esp{protégete} en prennent un (esdrújulos).
\end{corrige}

\begin{corrige}[Exercice 8 — Appariement : formes affirmatives et négatives]
\textbf{Correspondances : 1 — a ; 2 — c ; 3 — b ; 4 — d ; 5 — e.}
\begin{center}
\begin{tabularx}{\linewidth}{|X|X|X|}
\hline
\rowcolor{popgreenL} \textbf{Affirmatif} & \textbf{Négatif} & \textbf{Explication} \\
\hline
1. \esp{Dámelo} & a. \esp{No me lo des} & \emph{tú} : \esp{da} + enclise / \esp{no} + \esp{des} (subjonctif) \\
2. \esp{Levántense} & c. \esp{No se levanten} & \emph{ustedes}, verbe pronominal \\
3. \esp{Díselo} & b. \esp{No se lo digas} & \emph{tú} : \esp{di} + \esp{se} + \esp{lo} = \esp{díselo} / \esp{no se lo digas} (cohérence de personne) \\
4. \esp{Cuidemosla} & d. \esp{No la cuidemos} & \emph{nosotros} : \esp{cuidemos} + \esp{la} / \esp{no la cuidemos} \\
5. \esp{Póngansela} & e. \esp{No se la pongan} & \emph{ustedes} : \esp{pongan} + \esp{se} + \esp{la} = \esp{póngansela} (accent) / \esp{no se la pongan} \\
\hline
\end{tabularx}
\end{center}
\textbf{Méthode :} repérez d'abord la personne (désinence verbale), puis vérifiez que les pronoms sont identiques des deux côtés ; la transformation mécanique consiste à détacher les pronoms, à les placer après \esp{no}, et à passer le verbe au subjonctif présent.
\end{corrige}

\begin{corrige}[Exercice 9 — Compréhension et langue : texte de campagne sanitaire]
\textbf{Question 1 — Les formes impératives du texte :}
\begin{center}
\begin{tabularx}{\linewidth}{|X|X|X|}
\hline
\rowcolor{popblueL} \textbf{Forme} & \textbf{Personne} & \textbf{Polarité} \\
\hline
\esp{hiervan} & \emph{ustedes} & affirmatif \\
\esp{bébanla} & \emph{ustedes} & affirmatif (avec COD enclitique \esp{la}) \\
\esp{no tiren} & \emph{ustedes} & négatif \\
\esp{protéjanse} & \emph{ustedes} & affirmatif (réfléchi enclitique \esp{se}) \\
\esp{usen} & \emph{ustedes} & affirmatif \\
\esp{cuídense} & \emph{ustedes} & affirmatif (réfléchi enclitique) \\
\esp{lávense} & \emph{ustedes} & affirmatif (réfléchi enclitique) \\
\esp{no compartan} & \emph{ustedes} & négatif \\
\esp{acudan} & \emph{ustedes} & affirmatif \\
\hline
\end{tabularx}
\end{center}
(Les six consignes numérotées contiennent neuf formes verbales à l'impératif.)

\textbf{Question 2 — Place des pronoms :} à l'\textbf{impératif affirmatif}, les pronoms compléments (réfléchis, COD, COI) se placent \emph{après} le verbe et s'y soudent (enclise) : \esp{bébanla} (= \esp{beban} + \esp{la}), \esp{protéjanse}, \esp{cuídense}, \esp{lávense}. À l'\textbf{impératif négatif}, ils se placent \emph{avant} le verbe, entre \esp{no} et celui-ci (proclise) : le texte ne contient pas de négatif avec pronom, mais sur le modèle de \esp{no tiren la basura}, on aurait \esp{no la tiren} (« ne la jetez pas »).

\textbf{Question 3 — Accentuation :} \esp{bébanla}, \esp{cuídense}, \esp{lávense} portent un accent écrit parce que l'ajout du ou des pronoms enclitiques allonge le mot sans déplacer sa syllabe tonique : celle-ci se retrouve en troisième position à partir de la fin (mot \emph{esdrújulo}), et tout mot esdrújulo prend obligatoirement un accent écrit. Règle générale : \textbf{à l'impératif affirmatif avec enclise, on conserve la tonique d'origine et l'on ajoute un accent écrit chaque fois que la règle naturelle d'accentuation est enfreinte}.

\textbf{Question 4 — Traduction :}
\begin{enumerate}
\item \esp{Hiervan el agua y bébanla.} → « Faites bouillir l'eau et buvez-la. »
\item \esp{Protéjanse de los mosquitos: usen mosquiteras y cuídense.} → « Protégez-vous des moustiques : utilisez des moustiquaires et prenez soin de vous. »
\item \esp{Lávense las manos antes de cocinar y después de usar el baño.} → « Lavez-vous les mains avant de cuisiner et après être allé aux toilettes. »
\end{enumerate}
\textbf{Barème indicatif (sur 10) :} Q1 : 3 pts (relevé complet + personne + polarité) ; Q2 : 2,5 pts ; Q3 : 2 pts ; Q4 : 2,5 pts.
\end{corrige}

\begin{corrige}[Exercice 10 — Thème et version]
\textbf{A. Thème (français → espagnol) :}
\begin{enumerate}
\item \esp{\textbf{Lávate las manos antes de comer.}}
\item \esp{\textbf{No tire la basura al río.}}
\item \esp{\textbf{Dámela, por favor.}}
\item \esp{\textbf{No se la den.}}
\item \esp{\textbf{Levántaos y cuidad vuestro barrio.}}
\end{enumerate}
(Justifications détaillées : voir le corrigé de l'exercice 5, phrases identiques ; attention à \esp{Levántaos} avec accent.)

\textbf{B. Version (espagnol → français) :}
\begin{quote}
« Faites bouillir l'eau et buvez-la. Ne jetez pas les ordures dans la rivière. Protégez-vous des moustiques : utilisez des moustiquaires et prenez soin de vous. Lavez-vous les mains fréquemment. Ne partagez pas les objets personnels. Consultez le médecin au moindre symptôme. »
\end{quote}
\textbf{Barème indicatif (sur 10) :} thème : 5 pts (1 pt par phrase : forme verbale 0,5 + place des pronoms et accent 0,5) ; version : 5 pts (fidélité 3 pts, qualité du français 2 pts).

\textbf{Points de vigilance :}
\begin{itemize}
\item \esp{bébanla} : ne pas traduire « buvez-lui » ; \esp{la} reprend \esp{el agua} (féminin en espagnol !) → « buvez-la ». Faux ami de genre : \emph{l'eau} est masculin en espagnol mais se comporte en féminin pour l'accord du pronom (\esp{el agua} → \esp{la}).
\item \esp{ante cualquier síntoma} : « au moindre symptôme / dès le premier symptôme », non « avant n'importe quel symptôme ».
\item \esp{cuídense} : « prenez soin de vous », non « cuidez-vous ».
\end{itemize}
\end{corrige}

\begin{corrige}[Exercice 11 — Rédaction guidée : consignes écologiques pour le lycée]
\textbf{Proposition de rédaction modèle (environ 70 mots) :}
\begin{quote}
\esp{¡Atención, alumnos! Para cuidar nuestro liceo y el medio ambiente, sigan estas consignas:}
\esp{1. Apaguen las luces al salir del aula: apáguenlas siempre.}
\esp{2. No tiren la basura al suelo; sepárenla y recíclenla.}
\esp{3. Ahorren agua: cierren bien los grifos, no la desperdicien.}
\esp{4. Plantemos árboles y cuidémoslos juntos.}
\esp{5. Usen el transporte público o vengan a pie: protéjanse y protejan el planeta.}
\esp{¡Su barrio y su liceo están en sus manos!}
\end{quote}
\textbf{Traduction :} « Attention, élèves ! Pour prendre soin de notre lycée et de l'environnement, suivez ces consignes : 1. Éteignez les lumières en sortant de la salle : éteignez-les toujours. 2. Ne jetez pas les ordures par terre ; triez-les et recyclez-les. 3. Économisez l'eau : fermez bien les robinets, ne la gaspillez pas. 4. Plantons des arbres et prenons-en soin ensemble. 5. Utilisez les transports en commun ou venez à pied : protégez-vous et protégez la planète. Votre quartier et votre lycée sont entre vos mains ! »

\textbf{Vérification des contraintes :} 5 consignes à l'impératif (\emph{ustedes} et une à \emph{nosotros}) ; 6 formes avec pronoms : \esp{apáguenlas}, \esp{sepárenla}, \esp{recíclenla}, \esp{no la desperdicien} (proclise au négatif), \esp{cuidémoslos}, \esp{protéjanse} ; accents corrects sur toutes les formes esdrújules.

\textbf{Barème indicatif (sur 10) :} respect des contraintes (5 consignes, 60-80 mots, 3 pronoms minimum) : 3 pts ; correction morphologique de l'impératif : 3 pts ; place des pronoms et accentuation : 2 pts ; richesse du lexique thématique : 2 pts.

\textbf{Points de vigilance :} \esp{recíclenla}, \esp{apáguenlas}, \esp{sepárenla}, \esp{cuidémoslos}, \esp{protéjanse} prennent l'accent (esdrújulos) ; au négatif, jamais d'enclise (\esp{no la desperdicien}, non pas \esp{no desperdícenla}) ; \esp{el aula} est féminin (\esp{del aula}).
\end{corrige}

\newpage

\coursec{Fiche bilan}

\begin{popfigure}
\begin{tikzpicture}[mindmap, grow cyclic, every node/.style={concept, font=\small},
concept color=popblue, text=white,
level 1/.append style={level distance=3.4cm, sibling angle=60},
level 2/.append style={level distance=2.2cm, sibling angle=45, font=\footnotesize}]
\node[concept, font=\small\bfseries, align=center]{E08\\Impératif\\+ pronoms\\Podcast santé}
  child[concept color=poporange] { node[align=center]{Formes de\\l'impératif}
    child { node[align=center]{tú : 3e pers.\\sing. présent} }
    child { node[align=center]{négatif =\\subjonctif} }
  }
  child[concept color=popgreen] { node[align=center]{Pronoms\\compléments}
    child { node[align=center]{affirmatif :\\enclise (dámelo)} }
    child { node[align=center]{négatif :\\proclise (no me lo des)} }
  }
  child[concept color=popgold] { node[align=center]{Accentuation\\écrite}
    child { node[align=center]{dímelo,\\hazlo, dámelo} }
  }
  child[concept color=poppurple] { node[align=center]{Lexique\\santé}
    child { node[align=center]{la salud,\\el médico} }
  }
  child[concept color=poppink] { node[align=center]{Lexique\\environnement}
    child { node[align=center]{reciclar,\\el medio ambiente} }
  }
  child[concept color=popdark] { node[align=center]{Production\\podcast}
    child { node[align=center]{consignes,\\consejos} }
  };
\end{tikzpicture}
\legende{Carte mentale de la leçon E08 : impératif, pronoms et émission de santé.}
\end{popfigure}

\begin{bilan}
\textbf{1. Lexique clé (santé et environnement)}

\begin{center}
\begin{tabularx}{\linewidth}{|X|X|}\hline
\rowcolor{popblueL} \textbf{Español} & \textbf{Français} \\ \hline
la salud / estar sano & la santé / être en bonne santé \\ \hline
cuidarse, hacer ejercicio & se soigner, faire du sport \\ \hline
llevar una dieta equilibrada & suivre un régime équilibré \\ \hline
beber agua, dormir bien & boire de l'eau, bien dormir \\ \hline
el medio ambiente & l'environnement \\ \hline
reciclar, reutilizar, reducir & recycler, réutiliser, réduire \\ \hline
no tirar basura & ne pas jeter de déchets \\ \hline
ahorrar agua y energía & économiser l'eau et l'énergie \\ \hline
plantar árboles & planter des arbres \\ \hline
proteger la naturaleza & protéger la nature \\ \hline
\end{tabularx}
\end{center}

\textbf{2. Grammaire à connaître par cœur}

\begin{center}
\begin{tabularx}{\linewidth}{|X|X|}\hline
\rowcolor{popgreenL} \textbf{Règle} & \textbf{Exemple} \\ \hline
Impératif affirmatif \esp{tú} = 3\textsuperscript{e} pers. sing. du présent & \esp{habla, come, escribe} \\ \hline
8 irréguliers à \esp{tú} & \esp{di, haz, ve, pon, sal, sé, ten, ven} \\ \hline
\esp{vosotros} : infinitif, \textbf{-r} devient \textbf{-d} & \esp{hablad, comed, id} \\ \hline
\esp{usted / ustedes} = subjonctif présent & \esp{hable, hablen; vaya, vayan} \\ \hline
Impératif négatif = \esp{no} + subjonctif & \esp{no hables, no vengas, no fumen} \\ \hline
Affirmatif : pronoms \textbf{après}, soudés (enclise) & \esp{dámelo, hazlo, levántate} \\ \hline
Négatif : pronoms \textbf{avant}, détachés (proclise) & \esp{no me lo des, no lo hagas} \\ \hline
Ordre des pronoms : COI avant COD & \esp{dímelo} (me + lo) \\ \hline
Accent écrit pour garder la tonique & \esp{dame} mais \esp{dámelo} ; \esp{haz} mais \esp{házmelo} \\ \hline
\esp{vosotros} + \esp{os} : chute du \textbf{-d} & \esp{levantaos, callaos} (sauf \esp{idos}) \\ \hline
\end{tabularx}
\end{center}

\textbf{3. Méthodes express}
\begin{itemize}
  \item \textbf{Traduire un impératif français :} identifier la personne visée (\esp{tú / vosotros / usted / ustedes}) $\rightarrow$ choisir la forme $\rightarrow$ placer les pronoms (après-soudés si affirmatif, avant-détachés si négatif) $\rightarrow$ vérifier l'accent.
  \item \textbf{Version :} repérer \esp{no} + verbe = impératif négatif (subjonctif) ; un pronom soudé = impératif affirmatif.
  \item \textbf{Podcast :} saluer (\esp{Hola, oyentes}), annoncer le thème, enchaîner 4 à 6 conseils à l'impératif (\esp{Bebe agua, recicla, no tires basura}), conclure (\esp{¡Hasta pronto!}).
\end{itemize}
\end{bilan}

\begin{aretenir}[Checklist avant le Bac]
\begin{itemize}
  \item[$\square$] Je connais les 8 impératifs irréguliers de \esp{tú} : \esp{di, haz, ve, pon, sal, sé, ten, ven}.
  \item[$\square$] Je sais que l'impératif négatif utilise toujours le subjonctif présent.
  \item[$\square$] Je forme \esp{vosotros} en changeant \textbf{-r} en \textbf{-d} : \esp{hablad, comed, escribid}.
  \item[$\square$] À l'affirmatif, je soude les pronoms après le verbe : \esp{dámelo}.
  \item[$\square$] Au négatif, je place les pronoms avant le verbe : \esp{no me lo des}.
  \item[$\square$] Je mets le COI avant le COD : \esp{dímelo}, jamais \esp{dilome}.
  \item[$\square$] Je n'oublie pas l'accent écrit avec l'enclise : \esp{escríbeme, hazlo, dámelo}.
  \item[$\square$] Je sais que \esp{vosotros} + \esp{os} perd le \textbf{-d} : \esp{levantaos}.
  \item[$\square$] Je distingue \esp{sé} (impératif de \esp{ser}) et \esp{se} (pronom).
  \item[$\square$] Je connais le lexique de la santé et de l'environnement pour le thème.
  \item[$\square$] Je sais structurer un podcast : salutation, conseils à l'impératif, conclusion.
  \item[$\square$] Je relis ma traduction : personne, place des pronoms, accents.
\end{itemize}
\end{aretenir}

\findecours

\end{document}
